\chapter{Ab initio methods}\label{app:E}

A dark-matter rate in silicon needs a few dozen numbers about the crystal: a gap good to tens of meV, a loss function out to \SI{100}{eV}, density vertices at momentum transfers of tens of keV, electron--phonon couplings on grids of $10^9$ points, and the energy a struck atom leaves behind. Almost none of them is measured everywhere it is needed. Each is computed by a method with its own approximation, numerical parameters and typical failure. \Crefrange{app:E.1}{app:E.13} follow each method through one calculation, almost always of silicon, so that a number in the text can be traced to what it neglects and to the check that would expose it. \Crefrange{app:E.14}{app:E.27} turn this around and specify, quantity by quantity, the calculation behind each number the chapters use, with its inputs, its method and its convergence targets. Every method is placed on the truncation label $\mathfrak{T}$ of \cref{sec:8.5}, whose components include the Born--Oppenheimer order $n_B$, the electron--phonon order $n_g$, the anharmonic order $n_A$, the thermal order $n_T$ and the many-body approximation $\mathcal{X}$; a calculation at clamped nuclei, with independent particles and no phonons, sits at its base point $\mathfrak{T}_0$.

\section{Density-functional theory in practice}\label{app:E.1}

A density-functional calculation of silicon at its measured lattice constant returns, in seconds on a few processor cores, a total energy, an electron density and a set of Kohn--Sham bands with an indirect gap near \SI{0.6}{eV}. The measured gap at low temperature is \SI{1.17}{eV}~\cite[Eq.~(11)]{ramanathan2020ionization}. Both numbers matter for a dark-matter rate, and the work of a calculation lies in knowing which outputs to trust, how far, and what each numerical parameter does to them. \Cref{sec:3.4} stated the Kohn--Sham equations, derived in \cref{app:D.5}, and explained why their eigenvalues are not quasiparticle energies. We follow one calculation of silicon from its input to its bands and masses, and test which of its numbers can be used as they stand: the PBE functional, the norm-conserving silicon pseudopotential of the SG15 library~\cite{schlipf2015optimization} (four valence electrons per atom, \cref{app:E.2}), and the diamond structure with two atoms per cell.

\subsection{Computed quantities}

The Kohn--Sham equations of \cref{sec:3.4} are solved for the periodic parts $u_{n\vect{k}}$ of the Bloch states on a grid of crystal momenta. The states define the density
\begin{equation}
  \nd(\vect{r}) = \frac{2}{n_k}\sum_{\vect{k}}\sum_{n}f_{n\vect{k}}\,|u_{n\vect{k}}(\vect{r})|^2 ,
  \label{eq:E.1-density}
\end{equation}
in which the factor 2 counts spin, $n_k$ is the number of grid points and $f_{n\vect{k}}$ are the occupations. Since $u_{n\vect{k}}$ is normalized over the cell (\cref{app:A.5}), $\nd$ integrates over $\Omega$ to the number of valence electrons per cell. The density fixes the Kohn--Sham potential $V_{\mathrm{KS}}$, the potential fixes the states, and the loop is iterated until input and output agree. The outputs deserve different levels of trust: ground-state quantities, eigenvalues and states, and the responses built from both. The total energy, the density, and the forces and stresses derived from the energy are ground-state quantities, exact in principle for the exact functional (\cref{sec:3.4}); their error is that of the approximate functional and of the numerics. The eigenvalues $\varepsilon_{n\vect{k}}$ and states $\psi_{n\vect{k}}$ are auxiliary. Apart from the highest occupied level, they have no exact physical meaning, and the band gap they give is systematically too small. Part of the difference would remain even with the exact functional: the exact exchange--correlation potential jumps by a constant when an electron is added across the gap, and the Kohn--Sham gap leaves that derivative discontinuity out. The response quantities built from the states, such as $\varepsilon_\infty$, the phonons of \cref{app:E.4} or the density vertex $\rho_{mn}(\vect{k}, \vect{q})$ of \cref{app:E.7}, inherit the errors of both.

\subsection{Place on the hierarchy}

On the expansion hierarchy of \cref{sec:8.5}, a ground-state calculation sits at clamped nuclei, $n_B = 0$, and at zeroth thermal order, $n_T = 0$, with no thermal occupation of phonons: the nuclei are fixed at their equilibrium positions, and neither phonons nor zero-point motion enter. The Kohn--Sham states are the independent particles, $\mathcal{X} = \mathrm{IP}$, of the base point $\mathfrak{T}_0$, and every further many-body approximation (RPA, local fields, TDDFT, BSE; \cref{app:E.6}) is built on them. The exchange--correlation functional, however, is not a component of $\mathfrak{T}$. It is not an expansion in a small number, and no next order exists that would estimate its error. As \cref{sec:8.5} explains for the many-body approximation $\mathcal{X}$, the error is estimated instead by comparing functionals, by comparing with a quasiparticle calculation (\cref{app:E.3}), and by checking measured quantities such as the lattice constant, the phonon frequencies and the gap. \Cref{sec:8.8} turns that comparison into a nuisance parameter.

\subsection{Functionals}

A calculation chooses one approximation to the exchange--correlation energy $E_{\mathrm{xc}}$, and the choice sets the accuracy of everything downstream. The local-density approximation (LDA) takes the exchange--correlation energy per electron of the homogeneous electron gas at the local density, from the quantum Monte Carlo energies of Ceperley and Alder~\cite{ceperley1980ground} in the parametrization of Perdew and Zunger~\cite{perdew1981self}. The generalized-gradient approximation of Perdew, Burke and Ernzerhof (PBE)~\cite{perdew1996generalized} adds a dependence on $|\nabla\nd|/\nd^{4/3}$, built to satisfy exact constraints of the electron gas. Hybrid functionals mix in a fraction of Hartree--Fock exchange computed with the Kohn--Sham orbitals. The screened hybrid of Heyd, Scuseria and Ernzerhof (HSE)~\cite{heyd2003hybrid} takes one quarter of the exchange from Hartree--Fock, but only at short range. The Coulomb kernel $e^2/4\pi r$ is split, with an error function, into a part that decays on a length of a few \AA{} and a long-range remainder, and the exact exchange acts with the short-range part only. The screening reflects what happens in a semiconductor, where the exchange interaction between distant electrons is screened by the others (\cref{sec:4.8}).

The trends are systematic. LDA binds too strongly and gives lattice constants that are slightly too small, by \SIrange{0.2}{0.6}{\percent} for the materials of \cref{tab:E-functionals}; PBE binds too weakly and gives them too large, by \SIrange{0.9}{2.1}{\percent}. Over the elemental crystals, PBE equilibrium volumes are too large by \SI{3.6}{\percent} on average~\cite[Table~IV]{lejaeghere2014error}, about \SI{1.2}{\percent} in the lattice constant. The semilocal functionals underestimate gaps, LDA and PBE alike, and close the small gap of germanium altogether; HSE recovers most of the gap. Against the measured gaps of 41 semiconductors and insulators, HSE has a mean absolute error of \SI{0.38}{eV}, falling to \SI{0.18}{eV} for the typical semiconductors of the set, whose gaps are measured at low temperature~\cite[Table~I and text]{garza2016predicting}. \Cref{tab:E-functionals} collects the numbers for the book's running materials. In silicon the hybrid lattice constant is \SI{0.6}{\percent} too large and its gap \SI{0.05}{eV} too large. In germanium, where LDA gives a metal, the hybrid reopens the gap but leaves it \SI{0.2}{eV} too small.

\begin{table}[t]
  \centering
  \caption{Lattice constants and band gaps of silicon, germanium and gallium arsenide for three functionals, without spin--orbit coupling. Lattice constants from LDA and PBE are from Ref.~\cite[Table~II]{csonka2009assessing}, those from HSE from Ref.~\cite[Table~I]{lucero2012improved}. Gaps are evaluated at each functional's own relaxed lattice constant, from Ref.~\cite[Table~I]{lucero2012improved}; a zero gap means that the functional makes the crystal a metal. The static-lattice constants are the measured ones extrapolated to zero temperature with the zero-point anharmonic expansion removed~\cite[Table~II]{csonka2009assessing}; they are not measured at any temperature, and they lie below the lattice constants measured at low temperature. They are used only inside calculations, as the cell of a static lattice. The room-temperature lattice constants are measured: silicon~\cite[lattice parameter of silicon]{mohr2025codata}, germanium~\cite[Sect.~II\,A]{rideau2006strained} and GaAs~\cite{vurgaftman2001band}; they are the values for observables and the cell of the shared lattice setting of \cref{app:E.17}. The measured gaps are low-temperature values~\cite[Table~I]{lucero2012improved}.}
  \label{tab:E-functionals}
  \begin{tabular}{l ccccc ccc}
    \toprule
    & \multicolumn{5}{c}{lattice constant (\AA)} & \multicolumn{3}{c}{gap (eV)} \\
    \cmidrule(lr){2-6}\cmidrule(lr){7-9}
    & LDA & PBE & HSE & static lattice & room $T$ (expt.) & LDA & HSE & measured \\
    \midrule
    Si   & 5.403 & 5.466 & 5.451 & 5.416 & 5.431 & 0.51 & 1.22 & 1.17 \\
    Ge   & 5.623 & 5.759 & 5.706 & 5.640 & 5.658 & 0.00 & 0.54 & 0.74 \\
    GaAs & 5.605 & 5.746 & 5.711 & 5.638 & 5.653 & 0.43 & 1.18 & 1.52 \\
    \bottomrule
  \end{tabular}
\end{table}

The price of a hybrid is the exact exchange. A semilocal potential acts on a state $\psi_{n\vect{k}}$ by a multiplication in real space. The Fock operator instead needs, for every pair of the state with an occupied state $\psi_{m\vect{k}'}$, the potential of the pair density $\psi^*_{m\vect{k}'}\psi_{n\vect{k}}$. Each such potential is one Poisson solve, a fast Fourier transform (FFT) of the pair density, a division by $|\vect{k} - \vect{k}' + \vect{G}|^2$ and an FFT back. With $N_{\mathrm{occ}}$ occupied bands, applying the Fock operator to all $N_{\mathrm{b}}$ bands at one $\vect{k}$ therefore costs about $n_kN_{\mathrm{occ}}N_{\mathrm{b}}$ FFTs, against $N_{\mathrm{b}}$ for a semilocal potential, a factor $n_kN_{\mathrm{occ}}$ more. For silicon on an $8^3$ grid this factor is $512\times4 \approx 2\times10^3$. In practice the inner sum over $\vect{k}'$ runs over a coarser grid than the outer one, which reduces the factor by one or two orders of magnitude. A hybrid calculation of a two-atom cell therefore costs $10$ to $10^3$ times a semilocal one, minutes to hours where PBE takes seconds.

In this book the choice is usually PBE (or LDA) for structures, phonons and electron--phonon couplings, where semilocal functionals are accurate to a few percent (\cref{app:E.4}). Wherever an excitation energy sets a threshold, a gap correction, a hybrid functional or a quasiparticle calculation is used instead (\cref{app:E.3}). Note that the two choices do not mix freely: a phonon computed with PBE and a threshold computed with HSE describe two slightly different crystals, and the lattice constant at which both are evaluated must be stated.

\subsection{Plane waves and the cutoff}

The periodic part of each Bloch state is expanded in plane waves,
\begin{equation*}
  u_{n\vect{k}}(\vect{r}) = \Omega^{-1/2}\sum_{\vect{G}}c_{n\vect{k}}(\vect{G})\,\ee^{\ii\vect{G}\cdot\vect{r}}, \qquad \frac{|\vect{k} + \vect{G}|^2}{2m_e} \le E_{\mathrm{cut}} ,
\end{equation*}
truncated at the kinetic-energy cutoff $E_{\mathrm{cut}}$ (\cref{sec:3.5}). Codes quote $E_{\mathrm{cut}}$ in rydberg ($\SI{1}{Ry} = \SI{13.606}{eV}$) or hartree (two rydberg). The number of plane waves follows from counting reciprocal-lattice points in a sphere of radius $k_{\mathrm{cut}} = (2m_eE_{\mathrm{cut}})^{1/2}$. Each point occupies a volume $(2\pi)^3/\Omega$ of reciprocal space, so
\begin{equation}
  N_{\mathrm{PW}} \simeq \frac{4\pi k_{\mathrm{cut}}^3/3}{(2\pi)^3/\Omega} = \frac{\Omega\,k_{\mathrm{cut}}^3}{6\pi^2} .
  \label{eq:E.1-npw}
\end{equation}
For silicon, $\Omega = \SI{40.05}{\angstrom^3}$ (\cref{sec:3.1}) is 270.3 cubic bohr. At $E_{\mathrm{cut}} = \SI{60}{Ry}$, with $\hbar^2/2m_e = \SI{1}{Ry.bohr^2}$, $k_{\mathrm{cut}} = \sqrt{60}$ inverse bohr and \cref{eq:E.1-npw} gives 2121 plane waves per state. The density is a product of two states and has Fourier components up to $2k_{\mathrm{cut}}$, eight times as many, about $1.7\times10^4$. These set the size of the FFT grid on which the potential is applied. A component $\vect{G}$ has integer coordinates $m_i = \vect{G}\cdot\vect{a}_i/2\pi$ with $|m_i| \le 2k_{\mathrm{cut}}|\vect{a}_i|/2\pi$, which is 17.9 for $|\vect{a}_i| = a/\sqrt{2}$. The indices then run from $-17$ to 17, so the grid needs at least 35 points along each primitive vector, in practice 36, the next size that the FFT handles efficiently, and $36^3 \approx 4.7\times10^4$ points in all.

The cutoff is a variational parameter. For a fixed pseudopotential, enlarging the basis can only lower the total energy, so the energy converges from above and monotonically. The change between successive cutoffs estimates the error only once the scan has reached its asymptotic regime: in germanium the change from 50 to \SI{60}{Ry}, \SI{1.3}{meV}, is less than half the error that remains at \SI{60}{Ry}, \SI{3.1}{meV}. \Cref{tab:E-cutoff} shows the scan for silicon and germanium. Silicon converges to \SI{1}{meV} per atom at \SI{50}{Ry}, and to \SI{0.1}{meV} at about \SI{60}{Ry}. Germanium, whose pseudopotential carries the ten 3d electrons in the valence (\cref{app:E.2}), needs \SI{80}{Ry} for \SI{0.2}{meV}, and at \SI{30}{Ry} it is wrong by almost \SI{3}{eV} per atom. The pressure behaves differently in the two cases: in silicon it is stable to \SI{0.3}{kbar} from \SI{30}{Ry}, while in germanium it is not even of the right sign at \SI{30}{Ry}. The positive pressures at the measured lattice constant, \SI{24}{kbar} and \SI{38}{kbar}, are the PBE tendency to expand the lattice seen from the other side.

\begin{table}[t]
  \centering
  \caption{Convergence of the total energy and pressure with the plane-wave cutoff in the representative PBE calculations of silicon and germanium at their measured room-temperature lattice constants, \SI{5.431}{\angstrom} and \SI{5.658}{\angstrom} (\cref{tab:E-functionals}). The norm-conserving pseudopotentials carry four (Si) and fourteen (Ge, with 3d) valence electrons, and the grid has $8^3$ points $\vect{k}$. $\Delta E$ is the total energy per atom minus its value at \SI{120}{Ry}.}
  \label{tab:E-cutoff}
  \begin{tabular}{r rr rr}
    \toprule
    & \multicolumn{2}{c}{Si} & \multicolumn{2}{c}{Ge} \\
    \cmidrule(lr){2-3}\cmidrule(lr){4-5}
    $E_{\mathrm{cut}}$ (Ry) & $\Delta E$ (meV) & $P$ (kbar) & $\Delta E$ (meV) & $P$ (kbar) \\
    \midrule
    30 & 8.36 & 23.6 & 2880 & $-1615$ \\
    40 & 2.52 & 23.4 & 11.7 & 27.5 \\
    50 & 0.71 & 23.6 & 4.42 & 37.6 \\
    60 & 0.14 & 23.7 & 3.10 & 37.2 \\
    70 & 0.02 & 23.7 & 0.49 & 37.7 \\
    80 & 0.01 & 23.7 & 0.22 & 38.0 \\
    100 & 0.00 & 23.7 & 0.04 & 38.0 \\
    \bottomrule
  \end{tabular}
\end{table}

Two things go wrong with the cutoff: different quantities converge at different rates, and the stress picks up a Pulay term. Energy differences converge faster than energies, since most of the cutoff error sits near the nuclei and is the same in both structures, whereas absolute quantities, and quantities that probe the region near the nuclei, converge more slowly. The density vertex at large momentum transfer is the extreme case, treated in \cref{app:E.2}. At a fixed cutoff the number of plane waves changes with the cell volume, so the basis is not the same at two lattice constants, and the computed stress carries a spurious negative contribution that vanishes only at convergence. A structural relaxation must therefore be done at a cutoff where the stress, and not only the energy, is converged.

\subsection{Brillouin-zone sampling}

Sums over the zone become averages over a uniform grid of $n_k = n_1n_2n_3$ points, possibly shifted from the origin as proposed by Monkhorst and Pack~\cite{monkhorst1976special}. \Cref{app:C.2} derived the error of such a grid, \cref{eq:C.2-grid}: it is the sum of the real-space Fourier components of the integrand at the superlattice vectors $\vect{R}$ of the $n_1\times n_2\times n_3$ block of cells that the grid implies. For an insulator these decay exponentially, so a modest grid suffices. In the representative silicon calculation at \SI{60}{Ry}, the total energy changes by \SI{1.6}{meV} per atom from $8^3$ to $10^3$, by \SI{0.3}{meV} from $10^3$ to $12^3$ and by \SI{0.06}{meV} from $12^3$ to $14^3$. Symmetry reduces the work: on a $\Gamma$-centered $12^3$ grid, the 48 operations of the diamond structure leave 72 inequivalent points.

The grid needed depends on the quantity, however, and a converged total energy says little about a converged response. The electronic dielectric constant $\varepsilon_\infty$ (\cref{sec:4.4}), computed with density-functional perturbation theory (\cref{app:E.4}) at the PBE lattice constant, shows this. It is 14.36 on an $8^3$ grid, 13.57 on $10^3$, 13.29 on $12^3$ and 13.20 on $14^3$; at the measured lattice constant it reaches 13.04 on $18^3$. The total energy has converged to \SI{1}{meV} per atom long before $\varepsilon_\infty$ converges to one percent. The reason is that a response involves transitions between occupied and empty states, energy denominators $1/(\varepsilon_{c\vect{k}} - \varepsilon_{v\vect{k}})$ that are largest near the minimum direct gap, and a grid must resolve the region of the zone where they are largest. The same holds, more severely, for any quantity that probes a small region of the zone, such as the conduction-band minimum of silicon at $0.84$ of the way to X or the valley edges of a phonon-assisted window (\cref{app:E.5}).

\subsection{Metals and smearing}

In a metal the occupations jump from one to zero at the Fermi surface. The integrand of \cref{eq:C.2-grid} then has a discontinuity, its Fourier components decay only as a power of $|\vect{R}|$, and a grid converges slowly and erratically. The remedy is to replace the step by a smooth function of width $\sigma_{\mathrm{sm}}$, $f_{n\vect{k}} = \tilde f\bigl((\mu_{\mathrm{ch}} - \varepsilon_{n\vect{k}})/\sigma_{\mathrm{sm}}\bigr)$, which makes the integrand smooth on the scale of the grid. The smearing changes the answer, and the change must be controlled. The cleanest case is Fermi--Dirac smearing, $\tilde f(x) = 1/(\ee^{-x} + 1)$, where $\sigma_{\mathrm{sm}}$ is a fictitious temperature. The quantity that a self-consistent calculation minimizes is then the free energy $F_{\mathrm{sm}} = E - \sigma_{\mathrm{sm}}S_{\mathrm{sm}}$, with the entropy $S_{\mathrm{sm}}$ of the occupations. Let the density of states be smooth on the scale of $\sigma_{\mathrm{sm}}$, with $\mathcal{N}_{\mathrm{F}}$ per spin and volume at the Fermi level (\cref{sec:2.5}), so $2\Omega\mathcal{N}_{\mathrm{F}}$ per cell. The Sommerfeld expansion then gives
\begin{equation}
  E(\sigma_{\mathrm{sm}}) = E(\sigma_{\mathrm{sm}} = 0) + \frac{\pi^2}{3}\,\Omega\mathcal{N}_{\mathrm{F}}\,\sigma_{\mathrm{sm}}^2, \qquad
  F_{\mathrm{sm}}(\sigma_{\mathrm{sm}}) = E(\sigma_{\mathrm{sm}} = 0) - \frac{\pi^2}{3}\,\Omega\mathcal{N}_{\mathrm{F}}\,\sigma_{\mathrm{sm}}^2 ,
  \label{eq:E.1-sommerfeld}
\end{equation}
up to terms of order $\sigma_{\mathrm{sm}}^4$, with $E(\sigma_{\mathrm{sm}} = 0)$ the energy per cell of the unsmeared metal. The first relation is the electronic heat of a metal, with $\sigma_{\mathrm{sm}}$ in the role of the temperature; the second follows from the electronic entropy at the same order, $S_{\mathrm{sm}} = (2\pi^2/3)\Omega\mathcal{N}_{\mathrm{F}}\sigma_{\mathrm{sm}}$. The errors of $E$ and $F_{\mathrm{sm}}$ are equal and opposite, and the average $\tfrac12(E + F_{\mathrm{sm}})$ is correct to order $\sigma_{\mathrm{sm}}^4$. The forces, however, are derivatives of $F_{\mathrm{sm}}$, not of $E$, so they carry the full $\sigma_{\mathrm{sm}}^2$ error.

Other smearing functions remove the $\sigma_{\mathrm{sm}}^2$ term from the start. The functions of Methfessel and Paxton~\cite{methfessel1989high} are Gaussians multiplied by Hermite polynomials, whose low moments vanish, at the price of occupations that are not confined to $[0, 1]$. The ``cold smearing'' of Marzari and Vanderbilt~\cite[Eq.~(1)]{marzari1999thermal},
\begin{equation}
  \frac{\dd\tilde f}{\dd x} = \frac{1}{\sqrt{\pi}}\,\exp\!\Bigl[-\Bigl(x - \frac{1}{\sqrt{2}}\Bigr)^2\Bigr]\bigl(2 - \sqrt{2}\,x\bigr),
  \label{eq:E.1-cold}
\end{equation}
is the broadening that replaces the delta function. The occupations it integrates to are never negative, although they can exceed one, and it is chosen so that the free energy has no quadratic dependence on $\sigma_{\mathrm{sm}}$, and neither do its derivatives, forces and stresses included~\cite{marzari1999thermal}. A width of a few tenths of an eV is common for a structural calculation of a simple metal. Such a smearing must not be confused with a physical temperature. For a quantity that depends on the occupations near the Fermi surface, such as the density of states at $\mu_{\mathrm{ch}}$ of aluminum (\cref{sec:7.2}) or a Fermi-surface average of the electron--phonon coupling (\cref{app:E.5}), the smearing and the grid must be converged together. The tetrahedron method of \cref{app:C.2}~\cite{blochl1994improved} avoids smearing altogether. For an insulator such as silicon, all of this is unnecessary: occupations are fixed at one and zero, and smearing an insulator is an error, discussed at the end of this section.

\subsection{Self-consistency}

The density that goes into $V_{\mathrm{KS}}$ and the density that comes out of the states must agree. A calculation starts from a guess, usually a superposition of atomic densities, and iterates. Feeding the output straight back in diverges in all but the smallest systems, and the reason is screening. Let the input density differ from the self-consistent one by $\delta\nd_{\mathrm{in}}$. Its Hartree potential $v_{\mathrm{C}}\,\delta\nd_{\mathrm{in}}$ changes the states, and the output density differs by $\delta\nd_{\mathrm{out}} = R^0v_{\mathrm{C}}\,\delta\nd_{\mathrm{in}}$, with $R^0$ the static independent-particle response of \cref{sec:4.6} (we drop the exchange--correlation kernel, which does not change the argument). The dielectric matrix of the random-phase approximation is $\varepsilon = 1 - v_{\mathrm{C}}R^0$ (\cref{sec:4.4}); the same product in the other order, $\varepsilon^{\mathrm{T}} = 1 - R^0v_{\mathrm{C}}$, acts on densities and has the same eigenvalues. The output error is $(1 - \varepsilon^{\mathrm{T}})\,\delta\nd_{\mathrm{in}}$. Linear mixing takes the new input as $\nd_{\mathrm{in}} + \beta_{\mathrm{mix}}(\nd_{\mathrm{out}} - \nd_{\mathrm{in}})$, and the error then evolves as
\begin{equation}
  \delta\nd_{\mathrm{in}}^{(i+1)} = (1 - \beta_{\mathrm{mix}}\,\varepsilon^{\mathrm{T}})\,\delta\nd_{\mathrm{in}}^{(i)} .
  \label{eq:E.1-mixing}
\end{equation}
The iteration converges if $|1 - \beta_{\mathrm{mix}}\varepsilon_\lambda| < 1$ for every eigenvalue $\varepsilon_\lambda$ of the static dielectric matrix, that is for $\beta_{\mathrm{mix}} < 2/\max_\lambda\varepsilon_\lambda$. In an insulator the largest eigenvalue stays finite as the cell grows, of the order of $\varepsilon_\infty$, and a fixed mixing parameter works. In a metal, the Thomas--Fermi form $\varepsilon(\vect{G}) = 1 + k_{\mathrm{TF}}^2/G^2$ (\cref{sec:4.4}) grows without limit at small $G$, and the smallest $G$ of a cell of length $L$ is $2\pi/L$. A large metallic cell therefore needs $\beta_{\mathrm{mix}} \lesssim 2(2\pi/k_{\mathrm{TF}}L)^2$, and long-wavelength charge oscillates back and forth between iterations, a failure known as charge sloshing. The cure of Kerker~\cite{kerker1981efficient} is to mix each Fourier component with $\beta_{\mathrm{mix}}G^2/(G^2 + k_{\mathrm{TF}}^2)$, with a wavevector of the order of $k_{\mathrm{TF}}$ in the denominator, which cancels the $1/G^2$ of the Thomas--Fermi dielectric function, so that $\beta_{\mathrm{mix}}\varepsilon(\vect{G}) \approx \beta_{\mathrm{mix}}$ for all $\vect{G}$. Production codes combine such a preconditioner with a quasi-Newton update that uses the history of previous iterations to build an approximate inverse of $\varepsilon$, the modified Broyden method of Johnson~\cite{johnson1988modified} or the closely related Pulay mixing.

The convergence threshold of the loop must match the use of the result. A total energy to \SI{0.1}{meV} needs a loose threshold, since the energy is variational and its error is quadratic in the density error. A response calculation started from the converged ground state needs a tight one, since the first-order change of the states is computed against the ground-state Hamiltonian, and a residual error in the density appears linearly in the response (\cref{app:E.4}).

\subsection{Cost}

The work of an iteration is dominated by applying the Hamiltonian, local and nonlocal parts, to the $N_{\mathrm{b}}$ bands at each of the $n_k^{\mathrm{irr}}$ inequivalent $\vect{k}$ points, and by keeping those bands orthonormal. Applying the local part of the Hamiltonian to one band takes two FFTs, of order $N_{\mathrm{PW}}\log N_{\mathrm{PW}}$ each. The nonlocal projectors of the pseudopotential (\cref{app:E.2}) take of order $N_{\mathrm{PW}}N_{\mathrm{proj}}$, with $N_{\mathrm{proj}}$ the number of projectors in the cell. Keeping the bands orthonormal, with the diagonalization in the subspace they span, takes of order $N_{\mathrm{b}}^2N_{\mathrm{PW}} + N_{\mathrm{b}}^3$~\cite{payne1992iterative}. The cost per iteration is therefore
\begin{equation}
  t \propto n_k^{\mathrm{irr}}\bigl(N_{\mathrm{b}}N_{\mathrm{PW}}\log N_{\mathrm{PW}} + N_{\mathrm{b}}N_{\mathrm{PW}}N_{\mathrm{proj}} + N_{\mathrm{b}}^2N_{\mathrm{PW}}\bigr) .
  \label{eq:E.1-cost}
\end{equation}
Each term grows with the size of the cell. The number of bands and of projectors is proportional to the number of atoms $N_{\mathrm{at}}$, and so is $N_{\mathrm{PW}}$ at fixed cutoff, by \cref{eq:E.1-npw}. The last two terms therefore grow as $N_{\mathrm{at}}^3$ and dominate in large cells, while $n_k^{\mathrm{irr}}$ falls as $1/N_{\mathrm{at}}$ for a fixed sampling density. Raising the cutoff increases $N_{\mathrm{PW}}$ as $E_{\mathrm{cut}}^{3/2}$, which is why the 3d electrons of germanium, at \SI{80}{Ry} against \SI{60}{Ry}, cost a factor $(4/3)^{3/2} \approx 1.5$ in plane waves on top of the larger number of bands. Spin--orbit coupling doubles both the length of each state, now a two-component spinor, and the number of bands, and multiplies the $N_{\mathrm{b}}^2N_{\mathrm{PW}}$ term by eight. For a two-atom cell all this is small: the self-consistent calculation of silicon at \SI{60}{Ry} on a $12^3$ grid converges in eight iterations and about four seconds on four cores. Cost becomes a constraint only for supercells (\cref{app:E.10}), hybrids, the many empty states of a quasiparticle calculation (\cref{app:E.3}), and the dense grids of response and electron--phonon calculations (\cref{app:E.4,app:E.5}).

\subsection{Bands and effective masses}

Once the density is converged, it is held fixed and the Kohn--Sham Hamiltonian is diagonalized at any set of $\vect{k}$ points, on lines through the zone for a band structure, or on a fine grid around an extremum. This non-self-consistent step is cheap, because no iteration is needed, and it adds no approximation: the Kohn--Sham Hamiltonian is fixed by the converged density, and once the density is known it can be evaluated at any $\vect{k}$.

An effective mass is the inverse curvature of a band at an extremum, $(m^{*-1})_{\alpha\beta} = \partial^2\varepsilon_{n\vect{k}}/\partial k_\alpha\partial k_\beta$ (\cref{sec:3.3}). It is computed by finite differences, by $\vect{k}\cdot\vect{p}$ theory or by Wannier interpolation. The three-point finite-difference formula along a direction,
\begin{equation}
  \frac{1}{m^*} \simeq \frac{\varepsilon(k_0 + h) - 2\varepsilon(k_0) + \varepsilon(k_0 - h)}{h^2} ,
  \label{eq:E.1-fd}
\end{equation}
has a truncation error $h^2\varepsilon''''/12$ from the next term of the Taylor series, and an error of about $4\delta\varepsilon/h^2$ if the eigenvalues carry an error $\delta\varepsilon$. A small step amplifies the noise of the eigenvalues, a large step samples the departure from a parabola, and the step must be chosen between the two. A least-squares parabola through several points around the extremum is more robust than \cref{eq:E.1-fd}, and varying the fitting range tests both errors at once. The $\vect{k}\cdot\vect{p}$ route (\cref{sec:3.3}) uses perturbation theory with the computed velocity matrix elements, which gives the full tensor in one calculation at the extremum, provided enough bands are included and the velocities are corrected for the nonlocal pseudopotential (\cref{app:A.5,app:E.2}). Wannier interpolation (\cref{sec:3.6,app:E.5}) yields the band at any $\vect{k}$ and its derivatives analytically.

The masses inherit the errors of the functional. In $\vect{k}\cdot\vect{p}$ theory the curvature at an extremum comes from the coupling to the other bands at the same $\vect{k}$, each term weighted by the inverse energy separation. A mass set by a small direct gap, such as the conduction-band mass of GaAs at $\Gamma$, therefore follows the gap. A semilocal functional that puts the gap too low gives a mass that is too light, and a method that overestimates the gap gives one that is too heavy (\cref{app:E.3}). Masses are checked against cyclotron resonance, and gaps against optical absorption or photoemission, before they are used in a rate.

\begin{workedexample}{Silicon from input to bands}\label{we:E-si-bands}
The representative calculation uses the PBE functional, the SG15 norm-conserving pseudopotential with four valence electrons, a cutoff of \SI{60}{Ry} (\SI{816}{eV}) for the states and \SI{240}{Ry} for the density, fixed occupations of the four valence bands, a $\Gamma$-centered $12^3$ grid (72 points by symmetry), scalar-relativistic potentials without spin--orbit coupling, and a stopping criterion of $10^{-12}$\,Ry on the change of the total energy. A relaxation of the lattice at this cutoff gives $a = \SI{5.477}{\angstrom}$, \SI{0.85}{\percent} above the measured room-temperature \SI{5.431}{\angstrom} (\cref{sec:3.1}; the zero-temperature value of \cref{tab:E-functionals}, with the zero-point expansion removed, is \SI{5.416}{\angstrom}), in line with the PBE row of that table. At the relaxed lattice constant the valence-band maximum is the threefold $\Gamma$ state, and the conduction-band minimum lies on the line from $\Gamma$ to X. A scan of that line in steps of $0.005\times2\pi/a$ puts the minimum at $k_c = (0.840 \pm 0.003)\times2\pi/a$, against the measured position near $0.85$ (\cref{sec:3.3}), and the indirect gap at $\SI{0.602}{eV}$, about half of the measured \SI{1.17}{eV}. The direct gap at $\Gamma$ is \SI{2.55}{eV}. At the measured lattice constant, with the same pseudopotential at \SI{40}{Ry}, where band energies are converged to about \SI{2}{meV}, the indirect gap is \SI{0.557}{eV}: the \SI{0.85}{\percent} expansion moves it by about \SI{45}{meV}.

The longitudinal mass of the conduction valley follows from the same scan. A parabola fitted to the nine eigenvalues within $\pm0.02\times2\pi/a$ of the minimum gives $m^*_\parallel = 0.955\,m_e$; widening the window to 13 and 17 points gives $0.958$ and $0.951$. The spread, below one percent, is the combined truncation and rounding error of the fit, and $m^*_\parallel = 0.95\,m_e$ is the PBE value at this lattice constant. The measured value is $0.92\,m_e$~\cite{hensel1965cyclotron}. The PBE mass is \SI{4}{\percent} too heavy, while the gap is \SI{49}{\percent} too small: the longitudinal mass of the silicon valleys is insensitive to the error of the fundamental gap, and the functional error that matters for this material is the gap itself.
\end{workedexample}

\subsection{Pitfalls}

A semilocal functional can get the order of the bands wrong, and no later correction of the gap alone repairs that. In germanium, PBE without spin--orbit coupling at the measured lattice constant places the $\Gamma$ conduction state \SI{0.047}{eV} above the valence-band maximum and the L state at \SI{0.191}{eV}, while in experiment the L valley is lowest, at \SI{0.74}{eV}~\cite[Table~I]{lucero2012improved}, and $\Gamma$ lies above it. At the larger PBE lattice constant the gap closes and the crystal becomes a metal. A rigid upward shift of the conduction bands (\cref{app:E.3}) keeps the wrong order, so any germanium band input needs bands that pass the check of \cref{app:E.14} (valley order, both gaps, masses at L and $\Gamma$), with spin--orbit coupling. A hybrid at the standard exchange fraction inverts $\Gamma$ and L even with spin--orbit coupling, and a shift of the 4s-like orbitals fitted to both gaps gets the masses wrong (\cref{app:E.14}); a quasiparticle calculation is the next route under that check. The spin--orbit splitting of the valence-band maximum is \SI{0.296}{eV} in germanium against \SI{0.044}{eV} in silicon (\cref{sec:3.3}). The scalar-relativistic potentials of the worked example omit the splitting altogether; a calculation that includes it needs fully relativistic pseudopotentials and two-component spinors.

The gap error does not stay in the band structure: the dielectric constant of the representative calculation, 13.04, is \SI{11}{\percent} above the measured 11.7, the overscreening of \cref{app:E.4}. In germanium the bias becomes a failure: the near-metallic PBE state gives a dielectric constant of order $10^4$ that does not converge with the grid. Either a corrected gap or the measured screening is used in such cases.

Smearing hides a metal. A semiconductor computed with smeared occupations, as some codes do by default, converges smoothly even when the functional has closed its gap, and the result looks normal unless the band structure is inspected. With fixed occupations, the same calculation struggles to converge, which is the useful symptom.

The pseudopotential behind every number here is tested first, in \cref{app:E.2}; the first correction the bands need, for most dark-matter thresholds, is to the gap, which \cref{app:E.3} supplies.

\section{Pseudopotentials and all-electron reconstruction}\label{app:E.2}

The 1s orbital of silicon has a radius of about $a_0/Z = \SI{0.038}{\angstrom}$. To resolve it, a plane-wave basis needs wavevectors up to a few times $Z/a_0$. With $3Z/a_0 = 42$ inverse bohr, \cref{eq:E.1-npw} asks for a cutoff near \SI{1800}{Ry} and $3\times10^5$ plane waves per state, 160 times the 2121 of the converged calculation of \cref{app:E.1}, and the cost of each band grows at least in proportion. The ten core electrons do none of the bonding, and the shallowest of them, the 2p shell, is bound by about \SI{100}{eV} (\cref{app:E.7}). A pseudopotential removes them, together with the rapid oscillations that the valence states make in the core region to stay orthogonal to them, and replaces the nucleus and core by a soft effective potential that acts on smooth valence states. \Cref{sec:3.5} introduced the idea, the norm-conservation conditions and the projector-augmented-wave (PAW) transformation. What the pseudopotential gives up is the true valence wavefunction near the nucleus, exactly where the density vertex $\rho_{mn}(\vect{k}, \vect{q})$ of a heavy-mediator rate is evaluated, at momentum transfers of tens of keV. There a pseudopotential that is excellent for total energies, bands and phonons can fail without any sign in those quantities.

\subsection{The frozen core}

A pseudopotential acts on the valence states through a local radial potential $V_{\mathrm{loc}}(r)$ and a nonlocal part in the separable form of Kleinman and Bylander~\cite{kleinman1982efficacious},
\begin{equation}
  V_{\mathrm{NL}} = \sum_{\kappa}\sum_{ij}|\tilde p_{\kappa i}\rangle D_{\kappa,ij}\langle\tilde p_{\kappa j}| .
  \label{eq:E.2-KB}
\end{equation}
The projector functions $\tilde p_{\kappa i}(\vect{r})$, the nonlocal counterparts of the PAW projectors of \cref{sec:3.5}, are localized within a core radius $r_c$ of nucleus $\kappa$ and carry an angular momentum $l$ each. The coefficients $D_{\kappa,ij}$ form a small matrix. $V_{\mathrm{NL}}$ is the nonlocal pseudopotential of \cref{app:A.5}. The pseudo-wavefunctions it produces equal the all-electron valence wavefunctions outside the spheres of radius $r_c$ and are smooth and nodeless inside them. For a norm-conserving pseudopotential~\cite{hamann1979norm} the charge of each reference state inside $r_c$ is also the same, which fixes the scattering properties of the potential to first order in energy around the reference energies (\cref{sec:3.5}).

The approximation behind all of this is the frozen core. The core electrons are assumed to keep the orbitals they have in the isolated atom, unaffected by the crystal, the valence electrons, or a displacement of the nucleus. It is not a component of the truncation label $\mathfrak{T}$ of \cref{sec:8.5}, since it is not an expansion in a small parameter, and its errors are of three kinds: transferability, the missing core region of the valence states, and the missing core states. Transferability means that a potential built for the atom must describe the atom in the crystal, at other bond lengths and charge states. Tests control this error. The missing core region of the valence states leaves total energies and bands almost untouched but changes matrix elements of operators that probe short distances, such as $\ee^{\ii\vect{q}\cdot\vect{r}}$ at large $q$. The all-electron reconstruction of this section repairs it. The core states themselves are absent from the spectrum, so a pseudopotential calculation of a response has no core edges. They are added back as atomic orbitals (\cref{app:E.7,sec:10.6}).

\subsection{Construction}

A pseudopotential is built from an all-electron calculation of the isolated atom, followed by the construction of pseudo-wavefunctions, the removal of the valence screening and the separable form. The atomic calculation is scalar-relativistic (that is, with the mass--velocity and Darwin terms of \cref{sec:3.5} but without spin--orbit coupling) and uses a reference configuration such as $[\mathrm{Ne}]\,3s^23p^2$ for silicon. The choice of valence comes with it. Silicon has an obvious split, four valence electrons above a core bound by about \SI{100}{eV} and more. Germanium does not: its 3d shell lies about \SI{29}{eV} below the valence-band maximum in the crystal~\cite[Table~1-1]{williams2009electron}, overlaps the 4s and 4p orbitals in space, and is polarized by the crystal. A germanium potential either freezes the 3d shell, with four valence electrons, or puts it in the valence, with fourteen.

The pseudo-wavefunction is constructed for each angular momentum $l$. It equals the all-electron one beyond $r_c$, joins it smoothly at $r_c$, and is nodeless inside. The radius sets the trade-off. A large $r_c$ gives a smooth function and a low cutoff, but a potential that must describe a large region of the atom with few parameters, and that transfers worse. Within these constraints, the shape inside $r_c$ is chosen to make the plane-wave expansion converge fast. In the optimized norm-conserving construction of Hamann~\cite{hamann2013optimized}, the pseudo-wavefunction is expanded in a few spherical Bessel functions inside $r_c$. The coefficients minimize the residual kinetic energy above a wavevector $q_{\mathrm{opt}}$~\cite[Eq.~(1)]{hamann2013optimized},
\begin{equation}
  E^{\mathrm{res}}(q_{\mathrm{opt}}) = \int_{q_{\mathrm{opt}}}^\infty\dd q\,q^4\,|\phi(q)|^2 , \qquad \phi(q) = 4\pi\int_0^\infty\dd r\,r^2j_l(qr)\,\phi(r) ,
  \label{eq:E.2-residual}
\end{equation}
subject to norm conservation and continuity at $r_c$. Here $\phi(r)$ is the radial pseudo-wavefunction and $j_l$ the spherical Bessel function of \cref{app:C.3}. Up to a constant, $q^4|\phi(q)|^2$ is the kinetic-energy density in momentum space, $q^2/2m_e$ times the radial measure $q^2|\phi(q)|^2$, so $E^{\mathrm{res}}$ is the kinetic energy missed by a basis truncated at $q_{\mathrm{opt}}$. The cutoff of the solid is then set by $q_{\mathrm{opt}}$, and the tail above it decays fast. The same construction uses two projectors per angular momentum, which reproduce the scattering of the all-electron atom at two energies, typically a bound state and a scattering state above it, rather than one.

Next, the valence part of the potential is removed. The potential that reproduces the pseudo-wavefunctions is a screened one: it contains the Hartree and exchange--correlation potentials of the valence density of the reference atom, which must be subtracted, since the crystal supplies its own. The Hartree potential is linear in the density and subtracts cleanly. The exchange--correlation potential is not, and where the core and valence densities overlap, subtracting the valence part alone leaves an error that changes with the valence configuration. The nonlinear core correction of Louie, Froyen and Cohen~\cite{louie1982nonlinear} keeps a smoothed partial core density in the exchange--correlation term for this reason. Finally, the semilocal potential is cast into the separable form \cref{eq:E.2-KB}, which reduces the cost of applying it from a double sum over plane waves to a few projections (\cref{eq:E.1-cost}).

Ultrasoft pseudopotentials and the PAW method relax these conditions to lower the cutoff. Ultrasoft pseudopotentials~\cite{vanderbilt1990soft} drop norm conservation and restore the missing charge through augmentation functions localized in the spheres, at the price of a generalized eigenvalue problem. The PAW method~\cite{blochl1994projector} keeps an explicit linear map from the pseudo-wavefunctions to the all-electron ones (\cref{sec:3.5}), so that all-electron quantities can be evaluated after the calculation; Kresse and Joubert showed how ultrasoft potentials follow from it as an approximation~\cite{kresse1999ultrasoft}. Norm-conserving potentials remain the simplest to use in density-functional perturbation theory, in electron--phonon calculations and in many-body perturbation theory (\cref{app:E.3,app:E.4,app:E.5}), where every extra term of the ultrasoft and PAW formalisms has to be carried through each derivation~\cite{hamann2013optimized}.

The potentials of the representative calculations of \cref{app:E.1} come from the SG15 library of Schlipf and Gygi~\cite{schlipf2015optimization}, generated with Hamann's construction, with parameters optimized automatically against all-electron results. The construction parameters of the two potentials used here, as listed in their files, are inputs of the book's calculation. The silicon potential has four valence electrons, s and p channels with $r_c = 2.40$ and $3.56$ bohr (\SI{1.27}{\angstrom} and \SI{1.88}{\angstrom}), two projectors per channel, $q_{\mathrm{opt}} = 5.11$ and $3.50$ inverse bohr, and no nonlinear core correction. The germanium potential has fourteen valence electrons, s, p and d channels with $r_c$ between 2.54 and 2.67 bohr, two projectors each, and $q_{\mathrm{opt}}$ up to 6.36 inverse bohr for the 3d channel. The squares of these wavevectors, \SI{26}{Ry} and \SI{12}{Ry} for silicon and \SI{40}{Ry} for the germanium 3d, are the scale at which the residual kinetic energy has been made small. They are below the cutoffs at which the total energy converges in \cref{tab:E-cutoff}, \SI{60}{Ry} and \SI{80}{Ry}, because the residual \cref{eq:E.2-residual} decays above $q_{\mathrm{opt}}$ but does not vanish.

\subsection{Testing}

A pseudopotential is tested on the atom, for transferability, and in the solid before it is used. The logarithmic derivative $\dd\ln R_l(r)/\dd r$ of the radial solution at a radius just beyond $r_c$, plotted against energy for the all-electron and the pseudo atom, measures how well the potential scatters at energies other than its reference energies. The two curves must agree over the whole window of energies the calculation will use, which for a response calculation includes empty states tens of eV above the gap. A separable potential can also bind a spurious state, a ghost, below or among the physical ones, because the eigenstates of \cref{eq:E.2-KB} are not ordered in energy by their number of nodes~\cite{hamann2013optimized}. A ghost appears as a spurious step in the logarithmic derivative, or in the crystal as a flat band. The PseudoDojo tables were scanned for ghosts up to about \SI{200}{eV} above the Fermi level, a range chosen for optical and $GW$ calculations, and a few heavy elements still show one in that range, between 20 and \SI{80}{eV}~\cite[Table~8]{vansetten2018pseudodojo}. Ghosts at still higher energies are unavoidable~\cite{vansetten2018pseudodojo}.

For transferability, the potential, built in one configuration, is applied to excited and ionized configurations of the atom, and the excitation energies are compared with all-electron values. In the solid, the most systematic test is the $\Delta$ gauge of Lejaeghere and coworkers. It is the root-mean-square difference between two equations of state $E(\Vol)$ of the same crystal~\cite[Eq.~(15)]{lejaeghere2014error},
\begin{equation}
  \Delta_{\mathrm{gauge}} = \biggl[\frac{1}{\Delta \Vol}\int_{\Vol_0 - \Delta \Vol/2}^{\Vol_0 + \Delta \Vol/2}\dd \Vol\,\bigl(E_1(\Vol) - E_2(\Vol)\bigr)^2\biggr]^{1/2} , \qquad \Delta \Vol = 0.12\,\Vol_0 ,
  \label{eq:E.2-delta}
\end{equation}
with both curves set to zero at their minima. The window spans $\pm\SI{6}{\percent}$ of the reference volume $\Vol_0$. It folds the errors of the volume, the bulk modulus and its pressure derivative into one number in meV per atom. Different all-electron codes agree with each other to 0.3--\SI{0.5}{meV} per atom on average over the periodic table. The norm-conserving PseudoDojo tables reach a mean of about \SI{0.7}{meV} per atom against an all-electron reference (excluding the magnetic Cr, Mn and Fe), against \SI{1.4}{meV} for earlier norm-conserving tables~\cite{vansetten2018pseudodojo}. Direct comparisons with all-electron lattice constants and bulk moduli are the older form of the same test. For silicon in LDA, Hamann's optimized potential reproduces the all-electron lattice constant to the precision quoted, 10.21 bohr, and the bulk modulus to \SI{1.3}{\percent}, \SI{95.8}{GPa} against \SI{97.1}{GPa}~\cite[Table~I]{hamann2013optimized}. A potential must also be used with the functional with which it was generated; Hamann found that using a PBE-generated potential in an LDA calculation increased the differences from the all-electron results~\cite{hamann2013optimized}.

The standard tests are necessary but not sufficient for a dark-matter calculation. They probe the total energy and its derivatives with respect to volume, and the bands near the gap. Such a calculation also needs phonons (\cref{app:E.4}), which probe the potential under small displacements of the nucleus. It needs empty states far above the gap, which probe the logarithmic derivatives at high energy. And it needs the density vertex at momentum transfers up to tens of keV, which probes the core region. The last of these is not part of any standard test.

\subsection{Semicore states}

When a shell sits between the deep core and the valence, its treatment changes both the accuracy and the cost. The germanium 3d shell is the case the book needs. Keeping it in the valence improves transferability, since the shell responds to its environment; it puts the 3d states into the band structure, where they contribute to the response above an energy deposit of about \SI{30}{eV} (\cref{sec:10.6}); and it costs. The fourteen electrons need seven occupied bands per atom instead of two, and the localized d states need a higher cutoff. In \cref{tab:E-cutoff} germanium with the 3d in valence needs \SI{80}{Ry} where silicon needs \SI{60}{Ry}, and a cutoff of \SI{30}{Ry}, enough to give silicon to \SI{8}{meV} per atom, leaves the germanium energy almost \SI{3}{eV} per atom too high. Gallium arsenide, with the Ga 3d in the valence, is harder still: in the same kind of calculation the energy still changes by \SI{25}{meV} per atom between 100 and \SI{120}{Ry}. The alternative, a frozen 3d shell with a nonlinear core correction, gives a cheaper calculation of the valence and conduction bands and a response without the 3d channel, which must then be added as core states.

\subsection{The density vertex at large momentum}

The crystal form factors of \cref{sec:10.1} are built from the density vertex
\begin{equation}
  \rho_{mn}(\vect{k}, \vect{q}) = \langle\psi_{m,\vect{k}+\vect{q}}|\ee^{\ii\vect{q}\cdot\vect{r}}|\psi_{n\vect{k}}\rangle ,
  \label{eq:E.2-rho}
\end{equation}
between all-electron states. A pseudopotential calculation supplies the pseudo-states $\tilde\psi_{n\vect{k}}$ instead. Within PAW, the two are related by the transformation of \cref{sec:3.5}~\cite{blochl1994projector},
\begin{equation}
  |\psi_{n\vect{k}}\rangle = |\tilde\psi_{n\vect{k}}\rangle + \sum_{\kappa j}\bigl(|\phi^{\mathrm{AE}}_{\kappa j}\rangle - |\phi^{\mathrm{PS}}_{\kappa j}\rangle\bigr)\langle\tilde p_{\kappa j}|\tilde\psi_{n\vect{k}}\rangle .
  \label{eq:E.2-PAW}
\end{equation}
Here $\phi^{\mathrm{AE}}_{\kappa j}$ and $\phi^{\mathrm{PS}}_{\kappa j}$ are all-electron and pseudo partial waves that coincide outside the sphere of atom $\kappa$, and the projectors $\tilde p_{\kappa j}$ are dual to the pseudo partial waves inside it, $\langle\tilde p_{\kappa i}|\phi^{\mathrm{PS}}_{\kappa j}\rangle = \delta_{ij}$. The same map, with the partial waves of the atom the potential was built from, serves as an approximate reconstruction for a norm-conserving potential.

Inserting \cref{eq:E.2-PAW} into \cref{eq:E.2-rho} gives the vertex as a pseudo part plus one-center corrections. The result is compact because the partial-wave differences vanish outside the spheres and the pseudo partial waves are complete inside them. The differences vanishing outside means that cross terms between a difference and the smooth part of the other state live inside the spheres only. Completeness inside a sphere means $\sum_j|\phi^{\mathrm{PS}}_{\kappa j}\rangle\langle\tilde p_{\kappa j}| = 1$ when acting on $\tilde\psi$, which lets the smooth part of each state be replaced by its partial-wave expansion there. Collecting terms~\cite{blochl1994projector},
\begin{equation}
  \rho_{mn}(\vect{k}, \vect{q}) = \langle\tilde\psi_{m,\vect{k}+\vect{q}}|\ee^{\ii\vect{q}\cdot\vect{r}}|\tilde\psi_{n\vect{k}}\rangle
  + \sum_{\kappa ij}\langle\tilde\psi_{m,\vect{k}+\vect{q}}|\tilde p_{\kappa i}\rangle\,\ee^{\ii\vect{q}\cdot\vect{\tau}_\kappa}\,\Delta_{\kappa,ij}(\vect{q})\,\langle\tilde p_{\kappa j}|\tilde\psi_{n\vect{k}}\rangle ,
  \label{eq:E.2-rhoPAW}
\end{equation}
with the one-center matrix
\begin{equation}
  \Delta_{\kappa,ij}(\vect{q}) = \langle\phi^{\mathrm{AE}}_{\kappa i}|\ee^{\ii\vect{q}\cdot\vect{x}}|\phi^{\mathrm{AE}}_{\kappa j}\rangle - \langle\phi^{\mathrm{PS}}_{\kappa i}|\ee^{\ii\vect{q}\cdot\vect{x}}|\phi^{\mathrm{PS}}_{\kappa j}\rangle ,
  \label{eq:E.2-onecentre}
\end{equation}
where $\vect{x} = \vect{r} - \vect{\tau}_\kappa$ is measured from the nucleus. The sum over $\kappa$ runs over the atoms of one cell: the states are normalized over the cell (\cref{app:A.5}), by Bloch's theorem every cell contributes the same term, and the phase is written for the atom in the reference cell. The plane wave about the nucleus expands in spherical harmonics and spherical Bessel functions, $\ee^{\ii\vect{q}\cdot\vect{x}} = 4\pi\sum_{lm}\ii^lj_l(qx)Y^*_{lm}(\hat{\vect{q}})Y_{lm}(\hat{\vect{x}})$ (\cref{app:C.3}), so each element of $\Delta_{\kappa,ij}(\vect{q})$ is a sum of radial integrals of $j_l(qx)$ against products of partial waves, times Gaunt coefficients (integrals of three spherical harmonics) from the angular parts. These are computed once per species and per $|q|$.

The size of the correction follows from the small-$q$ expansion of \cref{eq:E.2-onecentre}. Its zeroth-order term is the difference of overlaps inside the sphere, which vanishes for norm-conserving potentials and is the augmentation charge of an ultrasoft or PAW potential. The next terms are differences of dipole and quadrupole moments of the partial-wave products inside the sphere, of relative size $q r_c$ and $(qr_c)^2$, weighted by the fraction of the state that lies inside the sphere. The pseudo vertex is therefore accurate for $qr_c \ll 1$ and must be checked beyond. In the opposite limit, $q \gg q_{\mathrm{opt}}$, the pseudo-wavefunctions have almost no Fourier components by construction, \cref{eq:E.2-residual}. The pseudo vertex then falls off much faster than the true one, whose large-$q$ tail comes from the nodes and the cusp of the valence states near the nucleus. In between lies the range where most electron-recoil rates are decided.

A plane-wave calculation adds one more limit, which is exact. A state at $\vect{k}$ has components only for $|\vect{k} + \vect{G}| \le k_{\mathrm{cut}}$, so the overlap of $\tilde\psi_{m,\vect{k}+\vect{q}}$ with $\ee^{\ii\vect{q}\cdot\vect{r}}\tilde\psi_{n\vect{k}}$ needs a pair of plane waves that both lie inside the cutoff sphere and differ by $\vect{q}$. Two points of a sphere of radius $k_{\mathrm{cut}}$ are at most $2k_{\mathrm{cut}}$ apart, and the pseudo vertex vanishes identically for $|\vect{q}| > 2k_{\mathrm{cut}}$. At \SI{60}{Ry} this limit is \SI{58}{keV}. Well below it, the vertex is already unconverged, and \cref{sec:3.5} gives the cutoff needed for a given $q$. Increasing the cutoff of a pseudopotential calculation makes the pseudo vertex converge to the pseudo answer, not to the all-electron one: beyond $q \sim 1/r_c$ only the reconstruction \cref{eq:E.2-rhoPAW} moves the result toward the true vertex.

Griffin and coworkers carried out the reconstruction for silicon and germanium with PAW potentials~\cite{griffin2021extended}. They computed the pseudo-states with a cutoff of \SI{1}{keV} (\SI{73.5}{Ry}) on a $10^3$ grid and reconstructed the all-electron states up to \SI{2}{keV} (\SI{147}{Ry})~\cite[Sect.~2.1]{griffin2021extended}, which holds momentum components up to $(2m_e\times\SI{2}{keV})^{1/2} \approx \SI{45}{keV}$. The reconstructed states have large-momentum components that the pseudo-states lack. For a heavy mediator, which weights large $q$, the reconstruction makes the differential rate much higher than the pseudo result at energy deposits above about \SI{15}{eV}, so that the rate integrated above a threshold is enhanced even for thresholds as low as about \SI{10}{eV}; in germanium it is what makes the 3d electrons dominate the heavy-mediator rate. For a light mediator, dominated by small $q$, the effect is smaller~\cite[Sects.~1 and~3.1]{griffin2021extended}. The deeper core states, 2p and below in silicon and 3d and below in germanium in their treatment, were modeled as atomic orbitals in Bloch sums, and the atomic description agreed with the reconstructed crystal states up to the reconstruction cutoff~\cite[Sect.~2.2]{griffin2021extended}.

\subsection{Nonlocal sum rules}

A nonlocal potential changes the commutators on which sum rules rest, and a test of a pseudopotential calculation against a sum rule must account for it. For a local potential, the double commutator of the Hamiltonian with $\ee^{\ii\vect{q}\cdot\vect{r}}$ involves only the kinetic energy and gives the $f$-sum rule of \cref{sec:4.3}. The plane wave $\ee^{\ii\vect{q}\cdot\vect{r}}$ is unitary. Inserting a complete set of states $|m\rangle$ of the one-electron Hamiltonian $H$ gives $\sum_m(\varepsilon_m - \varepsilon_n)|\langle m|\ee^{\ii\vect{q}\cdot\vect{r}}|n\rangle|^2 = \langle n|\ee^{-\ii\vect{q}\cdot\vect{r}}[H, \ee^{\ii\vect{q}\cdot\vect{r}}]|n\rangle = \langle n|\ee^{-\ii\vect{q}\cdot\vect{r}}H\ee^{\ii\vect{q}\cdot\vect{r}} - H|n\rangle$, and the same with $-\vect{q}$ in place of $\vect{q}$. Adding the two,
\begin{equation}
  \tfrac12\sum_m(\varepsilon_m - \varepsilon_n)\bigl(|\langle m|\ee^{\ii\vect{q}\cdot\vect{r}}|n\rangle|^2 + |\langle m|\ee^{-\ii\vect{q}\cdot\vect{r}}|n\rangle|^2\bigr) = \tfrac12\langle n|\ee^{-\ii\vect{q}\cdot\vect{r}}H\ee^{\ii\vect{q}\cdot\vect{r}} + \ee^{\ii\vect{q}\cdot\vect{r}}H\ee^{-\ii\vect{q}\cdot\vect{r}} - 2H|n\rangle .
  \label{eq:E.2-dc}
\end{equation}
The kinetic energy is shifted by the plane wave, $\ee^{-\ii\vect{q}\cdot\vect{r}}\vect{p}\,\ee^{\ii\vect{q}\cdot\vect{r}} = \vect{p} + \vect{q}$, so its contribution is $\tfrac12\langle(\vect{p} + \vect{q})^2 + (\vect{p} - \vect{q})^2 - 2\vect{p}^2\rangle/2m_e = q^2/2m_e$. A local potential commutes with the plane wave and drops out. A nonlocal one does not, and leaves
\begin{equation}
  \frac{q^2}{2m_e} + \delta_{\mathrm{NL}}(\vect{q}), \qquad
  \delta_{\mathrm{NL}}(\vect{q}) = \tfrac12\langle n|\ee^{-\ii\vect{q}\cdot\vect{r}}V_{\mathrm{NL}}\ee^{\ii\vect{q}\cdot\vect{r}} + \ee^{\ii\vect{q}\cdot\vect{r}}V_{\mathrm{NL}}\ee^{-\ii\vect{q}\cdot\vect{r}} - 2V_{\mathrm{NL}}|n\rangle ,
  \label{eq:E.2-fsumNL}
\end{equation}
on the right-hand side of \cref{eq:E.2-dc}. In the separable form \cref{eq:E.2-KB}, the conjugated potential $\ee^{-\ii\vect{q}\cdot\vect{r}}V_{\mathrm{NL}}\ee^{\ii\vect{q}\cdot\vect{r}}$ has the projectors $\ee^{-\ii\vect{q}\cdot\vect{r}}\tilde p_{\kappa i}$, so $\delta_{\mathrm{NL}}$ is a sum of projections and costs almost nothing to evaluate. At small $q$ it is of order $q^2$, like the kinetic term, and it does not vanish when all bands are summed: the pseudo-Hamiltonian obeys a modified $f$-sum rule. The all-electron Hamiltonian, whose potential is local, obeys the unmodified one. For a single Bloch state the same sum rule, without the symmetrization in $\pm\vect{q}$, keeps a term linear in $q$ set by the group velocity; \cref{eq:E.7-fsum} gives that form, and summed over a filled band it reduces to \cref{eq:E.2-dc}.

The representative silicon calculation shows the size of the effect. With the SG15 potential and the states at $\vect{k} = 0$, the left-hand side of \cref{eq:E.2-dc}, divided by $q^2/2m_e$, is 0.81 with 16 bands, 0.85 with 30, 0.86 with 50 and 0.88 with 100 at $|\vect{q}| = \SI{0.1}{\angstrom^{-1}}$ along [100], and nearly the same at \SI{0.5}{\angstrom^{-1}}. The sum still rises slowly at 100 bands, so part of the deficit is band truncation, and part is $\delta_{\mathrm{NL}}$. The two must be separated, by computing \cref{eq:E.2-fsumNL} directly, before the sum rule can be used as a convergence test of the number of bands (\cref{app:E.6,app:E.7}).

The limit $q \to 0$ of the same algebra gives the velocity, which carries the nonlocal commutator $\ii[V_{\mathrm{NL}}, \vect{r}]$ of \cref{app:A.5}; \cref{eq:E.7-vNL} gives its plane-wave form. The density vertex, on the other hand, needs no such correction: a scalar potential couples to the density by multiplication in their result, as in a local Hamiltonian~\cite[Eq.~(5)]{ismailbeigi2001coupling}. The pseudo vertex is therefore the correct vertex of the pseudo-Hamiltonian, and its only error is the missing core region, which \cref{eq:E.2-rhoPAW} restores.

\subsection{Silicon and the vertex}

The SG15 silicon potential, with $r_c = 2.40$ and $3.56$ bohr, has $1/r_c = 1.6$ and \SI{1.0}{keV}, and its optimization wavevectors $q_{\mathrm{opt}}$ correspond to 19 and \SI{13}{keV}. A dark-matter particle of \SI{10}{MeV} at the maximum speed in the Earth frame, $v_{\mathrm{max}} = 2.65\times10^{-3}$ (\cref{sec:1.3}), transfers at most $2m_\chi v_{\mathrm{max}} \approx \SI{53}{keV}$, and the typical momentum transfer of electron recoil is of the order of $\alpha m_e = \SI{3.7}{keV}$ (\cref{sec:1.3}). Much of the rate therefore comes from momentum transfers between $1/r_c$ and $q_{\mathrm{opt}}$, where the pseudo vertex of this potential is neither guaranteed accurate nor negligible, and a heavy-mediator rate weights the upper end of that range. A potential with a smaller $r_c$ pushes the problem to larger $q$ at the price of a higher cutoff. Griffin and coworkers handle this range by reconstructing the all-electron states up to a cutoff set by the largest $q$ needed, checking them against atomic orbitals for the shallowest core shell, and treating the deeper shells as atomic orbitals.

\subsection{Pitfalls}

Most failures of a pseudopotential in a dark-matter calculation come from using it outside the conditions it was tested for. One consequence of the frozen core lies outside the electronic response: core polarization by a displaced nucleus, which enters the large-$q$ lattice response, is identically zero in a frozen-core calculation and has to be computed separately (\cref{app:E.4}). With the representation of the states settled, the remaining error of the bands near the gap is that of the functional itself, the subject of \cref{app:E.3}.

\section{Quasiparticle corrections}\label{app:E.3}

A one-shot $GW$ calculation started from LDA puts the silicon gap at 1.24 to \SI{1.29}{eV} in the calculations of the 1980s and 1990s~\cite[Table~I]{aryasetiawan1998gw}, and at \SI{1.10}{eV} once it is converged with current methods~\cite[Sect.~IV\,C]{klimes2014predictive}, where LDA gives about \SI{0.5}{eV} and the measured gap at low temperature is \SI{1.17}{eV}~\cite[Eq.~(11)]{ramanathan2020ionization}. For an electronic threshold this correction of more than half an eV is the largest a calculation makes. Its effect on the minimum mass a target can detect, the reason Kohn--Sham gaps are too small, and what a rigid shift of the conduction bands does to the velocities are in \cref{sec:3.4}; the $GW$ self-energy $\Sigma = \ii\GF W^{\mathrm{scr}}$, with $W^{\mathrm{scr}} = \varepsilon^{-1}v_{\mathrm{C}}$, and the electron--electron lifetime from $\Imag\Sigma$ are defined in \cref{sec:4.8}. This section computes the corrections in practice, from the simplest, a scissor shift, to a full one-shot $G_0W_0$ calculation and its self-consistent successors, and shows what they do to the energies, the lifetimes and the matrix elements of a dark-matter rate and which parameters must be converged, at what cost.

\subsection{Computed quantities}

The quantities are the quasiparticle energies $E_{\mathrm{qp}}(n\vect{k})$, the positions of the poles of the one-particle Green's function $\GF$ (\cref{sec:4.7}), with their residues $z_{\mathrm{qp}}$ and their widths. In practice they are computed as corrections to the Kohn--Sham energies of \cref{app:E.1}, and the Kohn--Sham states are kept. A quasiparticle calculation changes the inputs of every later step (the threshold $E_{\mathrm{min}}(\vect{q})$ of \cref{sec:3.3}, the energy denominators of a response, the screening) but not its approximation level.

\subsection{Place on the hierarchy}

On the hierarchy of \cref{sec:8.5}, a response built from quasiparticle energies and Kohn--Sham states is still at $\mathcal{X} = \mathrm{IP}$ (or RPA, when the Coulomb interaction is resummed): the electron and the hole it leaves behind do not interact. Their attraction, the exciton, enters at $\mathcal{X} = \mathrm{BSE}$ (\cref{sec:4.9,app:E.6}), which uses the quasiparticle energies as input. The $GW$ approximation is itself not an expansion in a small number. It is the first term of Hedin's equations~\cite{hedin1965new}, with the vertex function set to one~\cite[Eqs.~(43)--(46) and (60)]{aryasetiawan1998gw}, and like the exchange--correlation functional its error is estimated by comparison, with measured gaps, with self-consistent variants and with vertex corrections, rather than by a next order.

\subsection{The scissor shift}

The simplest correction shifts every conduction band rigidly upward by $\Delta_{\mathrm{scis}}$, chosen so that the computed gap equals a measured one. For silicon in the PBE calculation of \cref{we:E-si-bands}, $\Delta_{\mathrm{scis}} = 1.170 - 0.602 = \SI{0.568}{eV}$ at the PBE lattice constant, and $\SI{0.613}{eV}$ at the measured one: the shift depends on the calculation it corrects. The approximation is better than it looks in silicon. A rigid shift reproduces the full $GW$ correction of silicon to about \SI{0.1}{eV} across the bands, and those of GaAs and diamond to 0.2 and \SI{0.4}{eV}~\cite{godby1988self}~\cite[Sect.~V\,B]{aryasetiawan1998gw}. The accuracy in silicon is partly fortuitous: the energy dependence of $\Sigma$ and its nonlocality would each make the correction strongly state-dependent, and in silicon the two almost cancel~\cite[Sect.~V\,B]{aryasetiawan1998gw}.

The scissor operator $H + \Delta_{\mathrm{scis}}\sum_{c}|\psi_{c\vect{k}}\rangle\langle\psi_{c\vect{k}}|$, summed over the conduction states, has the eigenstates of $H$ and shifted eigenvalues, and that fixes what it does to each ingredient of a rate. The threshold $E_{\mathrm{min}}(\vect{q})$ moves up by $\Delta_{\mathrm{scis}}$, and so does the energy of every transition. The density vertex $\rho_{cv}(\vect{k}, \vect{q})$ is unchanged, since the states are. The velocity matrix elements are not: the operator is nonlocal, it adds $\ii[\Delta_{\mathrm{scis}}\sum_c|\psi_c\rangle\langle\psi_c|, \vect{r}]$ to the velocity, and \cref{sec:3.4} showed that interband velocities are rescaled by $(\varepsilon_{c\vect{k}} - \varepsilon_{v\vect{k}} + \Delta_{\mathrm{scis}})/(\varepsilon_{c\vect{k}} - \varepsilon_{v\vect{k}})$. The $f$-sum rule is changed in the same way. With $\rho_{cv}$ fixed and every interband energy raised, the sum $\sum_c(\varepsilon_{c,\vect{k}+\vect{q}} - \varepsilon_{v\vect{k}} + \Delta_{\mathrm{scis}})|\rho_{cv}(\vect{k}, \vect{q})|^2$ exceeds its unshifted value by $\Delta_{\mathrm{scis}}\sum_c|\rho_{cv}|^2$, the extra term \cref{eq:E.2-fsumNL} generated by a nonlocal operator. A scissor-corrected response therefore violates the $f$-sum rule unless the transition strengths are rescaled with the energies, and a response used in an $f$-sum check must be built consistently (\cref{app:E.6}).

The shift also changes the screening. In the independent-particle response each transition contributes $|\rho_{cv}|^2$ divided by its energy, so with $\rho_{cv}$ fixed and the energies raised, $\varepsilon_\infty - 1$ decreases. The PBE value of silicon, 13.04 against the measured 11.7 (\cref{app:E.1}), moves in the direction of the measurement. Whether it lands above or below 11.7 depends on the rest of the calculation, local fields included (\cref{app:E.6}), so a scissored screening is not guaranteed to be closer to the measured one.

A scissor correction stands or falls with the target gap, the kind of gap and the band order. A measured gap at temperature $T$ contains the renormalization of the gap by the lattice vibrations, including their zero-point motion. Shifting a static-lattice calculation to the measured gap and then adding the electron--phonon renormalization of \cref{sec:13.7,app:E.8} counts it twice; in silicon the zero-point part is about \SI{60}{meV} (\cref{we:E-zpr}), small next to the gap but not against the window widths of \cref{chap:13}. As for the kind of gap, an optical absorption edge sits below the fundamental gap by the exciton binding energy (\cref{sec:4.9}), and a scissor to the optical edge puts the free-pair continuum too low. Finally, a scissor keeps the order of the Kohn--Sham bands, and where the functional gets it wrong, as PBE does in germanium (\cref{app:E.1}), no shift repairs it. There the quasiparticle calculation itself is needed.

\subsection{One-shot \texorpdfstring{$G_0W_0$}{G0W0}}

The workhorse for semiconductors is the one-shot $G_0W_0$ method of Hybertsen and Louie~\cite{hybertsen1986electron}: $\GF$ and $W^{\mathrm{scr}}$ are built once from the Kohn--Sham states, and the energies are corrected to first order.

The input is a Kohn--Sham calculation as in \cref{app:E.1}, extended to many empty bands, of the order of a hundred per atom (the count is made below). The screening is built from the static independent-particle response $R^0_{\vect{G}\vect{G}'}(\vect{q}, \omega = 0)$ of \cref{sec:4.6}, a sum over pairs of occupied and empty states of products of the plane-wave matrix elements
\begin{equation}
  \rho_{mn}(\vect{k}, \vect{q} + \vect{G}) = \langle\psi_{m,\vect{k}+\vect{q}}|\ee^{\ii(\vect{q}+\vect{G})\cdot\vect{r}}|\psi_{n\vect{k}}\rangle ,
  \label{eq:E.3-M}
\end{equation}
divided by the energy differences~\cite[Eqs.~(7) and (8)]{deslippe2012berkeleygw}. These are the density vertices of \cref{sec:3.2} at the unreduced momentum $\vect{q} + \vect{G}$ (in the periodic gauge, $\psi_{m,\vect{k}+\vect{q}+\vect{G}} = \psi_{m,\vect{k}+\vect{q}}$). The same objects that enter a dark-matter rate enter the screening of the electron's own interaction, and the large-$q$ caveats of \cref{app:E.2} apply to both. From $R^0$ follow the dielectric matrix of the random-phase approximation, $\varepsilon_{\vect{G}\vect{G}'} = \delta_{\vect{G}\vect{G}'} - v_{\mathrm{C}}(\vect{q} + \vect{G})R^0_{\vect{G}\vect{G}'}$ (\cref{sec:4.6}), its inverse, and the screened interaction $W^{\mathrm{scr}}_{\vect{G}\vect{G}'}(\vect{q}) = (\varepsilon^{-1})_{\vect{G}\vect{G}'}(\vect{q})\,v_{\mathrm{C}}(\vect{q} + \vect{G}')$~\cite[Eq.~(12)]{deslippe2012berkeleygw}.

The frequency dependence of $W^{\mathrm{scr}}$ comes from a model or from a frequency grid. In the generalized plasmon-pole model of Hybertsen and Louie, each element of $\Imag(\varepsilon^{-1})_{\vect{G}\vect{G}'}(\vect{q}, \omega)$ at $\omega > 0$ is a single delta function at a frequency $\tilde\omega_{\vect{G}\vect{G}'}(\vect{q})$, and the real part follows by the Kramers--Kronig relation. The two parameters of each element, the weight and the frequency of its pole, are fixed by two exact constraints: the static value $(\varepsilon^{-1})_{\vect{G}\vect{G}'}(\vect{q}, 0)$ just computed, and the $f$-sum rule for $\varepsilon^{-1}$, $\int_0^\infty\dd\omega\,\omega\,\Imag(\varepsilon^{-1})_{\vect{G}\vect{G}'} = -(\pi/2)\,\omega_p^2\,[(\vect{q} + \vect{G})\cdot(\vect{q} + \vect{G}')/|\vect{q} + \vect{G}|^2]\,\nd(\vect{G} - \vect{G}')/\nd(0)$, with $\nd(\vect{G})$ the Fourier components of the density~\cite[Eqs.~(83)--(86)]{aryasetiawan1998gw}. No parameter is adjusted. The model works because in sp semiconductors most of the spectral weight of $\Imag\varepsilon^{-1}$ sits in the plasmon, near the free-electron plasma frequency of the valence electrons, which for the eight valence electrons in the \SI{40.05}{\angstrom^3} cell of silicon is $\omega_p = \SI{16.6}{eV}$ (\cref{sec:4.3}). The alternative is to compute $\varepsilon^{-1}$ at many frequencies, on the real axis or on the imaginary axis followed by an analytic continuation. The continuation is usually a rational (Pad\'e) fit, and \cref{app:C.5} explains why such fits are unstable and must be checked against sum rules.

The self-energy is a sum over bands, $\vect{q}$ and $\vect{G}$ of products of the matrix elements \cref{eq:E.3-M}, weighted by $v_{\mathrm{C}}$ and by the frequency structure of $\varepsilon^{-1}$. It is organized as a screened-exchange part, a sum over occupied states only, and a Coulomb-hole part, a sum over all bands, occupied and empty~\cite[Eqs.~(19)--(22)]{deslippe2012berkeleygw}. Codes evaluate separately the bare exchange $\Sigma_{\mathrm{x}}$, the Fock exchange of \cref{sec:3.4} built from the occupied Kohn--Sham states, and the correlation part $\Sigma_{\mathrm{c}}$ that $W^{\mathrm{scr}} - v_{\mathrm{C}}$ adds to it.

The quasiparticle equation, which corrects the energies, replaces $v_{\mathrm{xc}}$ in the Kohn--Sham equation by $\Sigma(E)$, evaluated at the quasiparticle energy itself. Keeping only the diagonal matrix elements in the Kohn--Sham basis, which assumes that the Kohn--Sham states are close to the quasiparticle states, and expanding $\Sigma(E)$ to first order about $\varepsilon_{n\vect{k}}$ gives~\cite[Eqs.~(57) and (58)]{aryasetiawan1998gw}
\begin{equation}
  E_{\mathrm{qp}}(n\vect{k}) = \varepsilon_{n\vect{k}} + z_{\mathrm{qp}}\,\Real\langle\psi_{n\vect{k}}|\Sigma(\varepsilon_{n\vect{k}}) - v_{\mathrm{xc}}|\psi_{n\vect{k}}\rangle, \qquad
  z_{\mathrm{qp}} = \biggl[1 - \frac{\partial\Real\langle\psi_{n\vect{k}}|\Sigma(\omega)|\psi_{n\vect{k}}\rangle}{\partial\omega}\biggr]^{-1}_{\omega = \varepsilon_{n\vect{k}}} .
  \label{eq:E.3-qp}
\end{equation}
The residue $z_{\mathrm{qp}}$ is 0.7 to 0.8 in semiconductors~\cite[Sect.~V\,B]{aryasetiawan1998gw}. The slope of $\Real\Sigma$ is negative, so $z_{\mathrm{qp}} < 1$ reduces the correction: neglecting the energy dependence of $\Sigma$, as the static Coulomb-hole plus screened-exchange (COHSEX) approximation does, overestimates the gap correction~\cite[Sect.~V\,B]{aryasetiawan1998gw}. The off-diagonal elements, when kept, turn \cref{eq:E.3-qp} into a small matrix diagonalization in the Kohn--Sham basis~\cite[Eqs.~(17) and (18)]{deslippe2012berkeleygw}, which changes the states as well as the energies.

\begin{table}[t]
  \centering
  \caption{Minimum band gaps of silicon, germanium and GaAs in LDA, in one-shot $GW$ calculations started from LDA, and from experiment, in eV, as collected in Ref.~\cite[Table~I]{aryasetiawan1998gw}. The $GW$ column gives the range of the independent calculations of Refs.~\cite{hybertsen1986electron,godby1988self,rohlfing1993quasiparticle,zhu1991quasiparticle}. The LDA gap of germanium is negative (the crystal is a metal); for GaAs the LDA gap is the one of \cref{tab:E-functionals}, \SI{0.43}{eV}. The last two columns are the $GW$ calculations of Shirley, Zhu and Louie~\cite{shirley1992core}, in which the core--valence exchange and correlation are treated with the local-density approximation (``core in LDA'') or with a core-polarization potential, as collected in Ref.~\cite[Table~II]{aryasetiawan1998gw}: the silicon gap from $\Gamma$ to the conduction minimum at $0.85\,$X, and the direct ($\Gamma$) and indirect (L) gaps of germanium, with spin--orbit coupling. The measured direct gap of germanium is \SI{0.89}{eV}~\cite[Table~II]{aryasetiawan1998gw}.}
  \label{tab:E-gw}
  \begin{tabular}{l ccc cc}
    \toprule
    & LDA & $GW$ & measured & $GW$, core in LDA & $GW$, core polarization \\
    \midrule
    Si & 0.52 & 1.24--1.29 & 1.17 & $0.85\,$X 1.29 & $0.85\,$X 1.16 \\
    Ge & $<0$ & 0.65--0.75 & 0.744 & $\Gamma$ 0.53, L 0.58 & $\Gamma$ 0.85, L 0.73 \\
    GaAs & -- & 1.22--1.58 & 1.52 & -- & -- \\
    \bottomrule
  \end{tabular}
\end{table}

\Cref{tab:E-gw} shows what these $G_0W_0$ calculations achieve for the book's running materials. In silicon they correct most of the LDA error and overshoot the measured gap by 0.07 to \SI{0.12}{eV}. In germanium and GaAs the independent calculations scatter by 0.1 and \SI{0.36}{eV}, a measure of the sensitivity of the method to its numerical and physical choices rather than of its intrinsic accuracy. One such choice is the treatment of the core. The calculation of Shirley, Zhu and Louie~\cite{shirley1992core} includes spin--orbit coupling and resolves both gaps. In it, approximating the exchange and correlation of the valence electrons with the 3d and deeper shells of germanium within LDA puts the direct gap below the indirect one, the wrong order. Treating the core polarization explicitly restores the order and brings both gaps within \SI{0.04}{eV} of experiment, and it lowers the silicon gap from 1.29 to \SI{1.16}{eV} (\cref{tab:E-gw}, last two columns). The core electrons must therefore be treated beyond LDA in the self-energy, either in the valence or through a core-polarization model, and not only in the density.

\subsection{Self-consistency and the starting point}

A one-shot calculation inherits its starting point. $G_0W_0$ on LDA states underestimates gaps systematically once core states are treated adequately, and it can fail qualitatively in transition-metal compounds~\cite{vanschilfgaarde2006quasiparticle}. Updating the energies in $\GF$ and $W^{\mathrm{scr}}$ until they are self-consistent (eigenvalue self-consistency) removes part of the dependence. The quasiparticle self-consistent $GW$ method of van Schilfgaarde, Kotani and Faleev goes further and determines the best static, Hermitian one-particle Hamiltonian from the self-energy itself, so that the result no longer depends on the starting functional~\cite{vanschilfgaarde2006quasiparticle}. Its errors are small and systematic. It overestimates the gaps of sp semiconductors by about \SI{0.2}{eV} for transitions between $\Gamma$ points and \SI{0.1}{eV} for the others, and it slightly underestimates $\varepsilon_\infty$. The reason given is that the random-phase approximation of $W^{\mathrm{scr}}$ omits the electron--hole attraction, which makes the screening too weak and the gap too large~\cite{vanschilfgaarde2006quasiparticle}. A starting point that is metallic, LDA or PBE germanium, is the case where the choice matters most: the first $R^0$ then contains intraband transitions that an insulator does not have, and either a hybrid-functional start or a self-consistent scheme is used.

Current practice differs from the calculations of \cref{tab:E-gw} in three ways~\cite[Sects.~IV\,C--F, V and~XI\,D]{golze2019gw}. The quasiparticle energies converge only as the inverse of the number of plane waves in the basis for the empty states and the response, so that the few hundred empty bands per atom of the early calculations can leave errors of tenths of an eV. Modern calculations extrapolate in $1/N_{\mathrm{PW}}$, or avoid the sum over empty states with a Sternheimer equation or a static remainder; converged in this way, $G_0W_0$ started from LDA gives \SI{1.10}{eV} for silicon~\cite[Sect.~IV\,C]{klimes2014predictive}, below the measured gap rather than above it. The frequency integral is done by contour deformation or on the imaginary axis rather than with a plasmon pole. And the missing electron--hole attraction in $W^{\mathrm{scr}}$ is put back as a vertex correction: adding the ladder diagrams of the Bethe--Salpeter equation (\cref{sec:4.9}) to the polarizability of quasiparticle self-consistent $GW$ removes most of its overestimate of the gaps and its underestimate of $\varepsilon_\infty$, and gives \SI{1.08}{eV} for the silicon gap near X against the measured \SI{1.17}{eV}~\cite[Table~1]{cunningham2023qsgw}.

\subsection{Lifetimes}

The imaginary part of the self-energy gives the width of the quasiparticle. Near the pole, the linearized quasiparticle energy of \cref{eq:E.3-qp} acquires an imaginary part $z_{\mathrm{qp}}\Imag\Sigma$, and the probability that the electron is still in the state decays as $\exp(-2z_{\mathrm{qp}}|\Imag\Sigma|\,t)$. The decay rate of the electron, its inverse lifetime, is therefore $2z_{\mathrm{qp}}|\Imag\Sigma(E_{\mathrm{qp}}(n\vect{k}))|$. By the cutting rule of \cref{sec:4.8}, $\Imag\Sigma$ in the $GW$ approximation is the golden-rule rate of the decay of the electron into a lower state with the emission of an excitation of the screened interaction: an electron--hole pair or a plasmon. An electron in the conduction band of an insulator can create an electron--hole pair only if it has the energy to cross the gap. Its energy above the conduction-band minimum must exceed $\Eg$, and the conservation of crystal momentum raises the true threshold above this bound. Below it the electron--electron width vanishes, and only phonon emission (\cref{sec:6.3}) relaxes the electron. Above it, impact ionization opens, and its rate grows with the phase space of the final pairs. These two rates, as functions of energy, are the input of the hot-carrier cascades that set the mean energy per electron--hole pair $\varepsilon_{eh}$ (\cref{app:E.13}).

The plasmon-pole model, with all the weight of $\Imag\varepsilon^{-1}$ in delta functions at the plasmon energies, gives $\Imag\Sigma = 0$ except at those energies, and it cannot produce lifetimes~\cite[Sect.~III\,C]{aryasetiawan1998gw}. A lifetime needs the full frequency dependence of $\varepsilon^{-1}$ near the energies of the pairs that are created, which for silicon starts at the gap. A full-frequency calculation costs more than a plasmon-pole one by about the number of frequencies at which $\varepsilon^{-1}$ is computed.

\subsection{Matrix elements}

The diagonal approximation of \cref{eq:E.3-qp} keeps the Kohn--Sham states and changes the energies only, and its effect on matrix elements follows from that alone. The density vertex $\rho_{mn}(\vect{k}, \vect{q})$ is unchanged. The velocity changes, and its change can be computed exactly within the approximation. For eigenstates of a Hamiltonian $H$, the interband velocity is $\vect{v}_{mn} = \ii\langle m|[H, \vect{r}]|n\rangle = \ii(\varepsilon_m - \varepsilon_n)\langle m|\vect{r}|n\rangle$ for $m \neq n$. Between Bloch states at the same $\vect{k}$ the interband position matrix element is well defined, and this relation, applied to the Kohn--Sham Hamiltonian, gives it as $\langle m|\vect{r}|n\rangle = -\ii\vect{v}^{\mathrm{KS}}_{mn}/(\varepsilon_m - \varepsilon_n)$. A quasiparticle Hamiltonian that is diagonal in the Kohn--Sham basis has the same $\langle m|\vect{r}|n\rangle$ and the energies $E_{\mathrm{qp}}$, so
\begin{equation}
  \vect{v}^{\mathrm{qp}}_{mn}(\vect{k}) = \vect{v}^{\mathrm{KS}}_{mn}(\vect{k})\,\frac{E_{\mathrm{qp}}(m\vect{k}) - E_{\mathrm{qp}}(n\vect{k})}{\varepsilon_{m\vect{k}} - \varepsilon_{n\vect{k}}} , \qquad m \neq n .
  \label{eq:E.3-vrescale}
\end{equation}
For a rigid shift this is the scissor rescaling of \cref{sec:3.4}. For a state-dependent correction it rescales each transition by its own factor. The velocity vertex enters the absorption of dark matter at small momentum transfer and the optical limit of $\rho_{mn}$ (\cref{sec:3.2}). \Cref{eq:E.3-vrescale} keeps the optical limit of the density vertex, $\vect{q}\cdot\vect{v}_{mn}/(E_m - E_n)$, equal to that of the uncorrected calculation, as it must be when the states are unchanged. The intraband velocities, the group velocities, follow from the dispersion of the corrected bands. Since the corrections are computed on a coarse grid, they are interpolated to fine grids with Wannier functions (\cref{app:E.5}), with the same gauge as every other interpolated quantity (\cref{sec:3.6}).

Effective masses change with the dispersion. The $\vect{k}\cdot\vect{p}$ argument of \cref{app:E.1} connects a mass at a direct gap to the gap, and the quasiparticle self-consistent calculation of GaAs shows the connection: it overestimates the gap, \SI{1.77}{eV} against \SI{1.52}{eV}, and the conduction-band mass at $\Gamma$, $0.077\,m_e$~\cite[Fig.~2 and text]{vanschilfgaarde2006quasiparticle} against the measured $0.067\,m_e$ (\cref{sec:3.3}).

\subsection{Convergence and cost}

A $G_0W_0$ calculation is converged in the number of empty bands, the dielectric cutoff, the grid of crystal momenta and the frequency treatment, and these parameters are coupled. The number of empty bands $N_{\mathrm{b}}$ enters the sums for $R^0$ and for $\Sigma_{\mathrm{c}}$. The dielectric cutoff is the kinetic energy $E^{\varepsilon}_{\mathrm{cut}}$ below which the vectors $\vect{q} + \vect{G}$ of the matrix $\varepsilon_{\vect{G}\vect{G}'}$ are kept. The two are not independent. The highest empty band included should have an energy of about the dielectric cutoff, since a plane wave $\vect{q} + \vect{G}$ resolves transitions into states up to that energy~\cite[Sect.~4]{deslippe2012berkeleygw}. A free-electron count makes the link concrete: each plane wave with kinetic energy below $E$ contributes one band below $E$, so the number of bands below $E$ is about $N_{\mathrm{PW}}(E)$ of \cref{eq:E.1-npw}. For silicon, a dielectric cutoff of \SI{10}{Ry} keeps 144 vectors $\vect{G}$ and asks for about 150 bands; \SI{15}{Ry} keeps 265 and asks for about 270. The sum for the Coulomb-hole part of $\Sigma_{\mathrm{c}}$ converges much more slowly with $N_{\mathrm{b}}$ than the sum for $R^0$. A test that converges $R^0$ alone, or one parameter with the other held fixed, can therefore stop at a value far from the converged one~\cite[Sect.~4 and Fig.~8]{deslippe2012berkeleygw}. The grid of $\vect{k}$ and $\vect{q}$ points is the same for both. The divergence of $v_{\mathrm{C}}(\vect{q})$ at $\vect{q} \to 0$ is handled by averaging $v_{\mathrm{C}}$ and $W^{\mathrm{scr}}$ over the small cell around $\vect{q} = 0$ and by computing the head and wings of $\varepsilon_{\vect{G}\vect{G}'}$ (its $\vect{G} = \vect{G}' = 0$ element and the rest of the $\vect{G} = 0$ row and column) at $\vect{q} \to 0$ from the optical limit of the vertex (\cref{app:E.6}). The frequency treatment is either a plasmon pole or a frequency grid. The bare exchange $\Sigma_{\mathrm{x}}$ is a separate case: it involves the unscreened interaction, whose $1/|\vect{q} + \vect{G}|^2$ falls off slowly, and its sum over $\vect{G}$ runs to the cutoff of the pair densities rather than to the dielectric cutoff.

The cost follows from the structure of the sums~\cite[Sect.~4]{deslippe2012berkeleygw}. The number of occupied and of empty bands, and the number of plane waves, each grow in proportion to the number of atoms $N_{\mathrm{at}}$. Computing the matrix elements \cref{eq:E.3-M} for all pairs of bands is $N_{\mathrm{at}}^2$ pairs times an FFT of size $N_{\mathrm{at}}\log N_{\mathrm{at}}$, so $N_{\mathrm{at}}^3\log N_{\mathrm{at}}$. Summing them into $R^0_{\vect{G}\vect{G}'}$ is $N_{\mathrm{at}}^2$ pairs for each of $N_{\mathrm{at}}^2$ pairs of $\vect{G}$ vectors, so $N_{\mathrm{at}}^4$, done as one dense matrix product; the inversion of $\varepsilon$ is $N_{\mathrm{at}}^3$ per $\vect{q}$ and per frequency. The self-energy of one state is $N_{\mathrm{at}}^3$, and all states in an energy window add a factor $N_{\mathrm{at}}$. A $G_0W_0$ calculation thus scales formally as the fourth power of the system size, against the third power of \cref{eq:E.1-cost}, with a prefactor set by the hundreds of empty bands. For the two-atom cell of silicon all of this remains modest. For the supercells of defects (\cref{app:E.10}) and for the dense grids of a finite-momentum exciton calculation (\cref{app:E.6}), it is the step that limits what can be computed.

\begin{workedexample}{The silicon gap from three corrections}\label{we:E-si-gap}
A scissor, a hybrid functional and a $G_0W_0$ calculation each lead from the PBE calculation of \cref{we:E-si-bands} to a gap for a silicon threshold. The scissor to the measured gap at zero temperature uses $\Delta_{\mathrm{scis}} = \SI{0.568}{eV}$ at the PBE lattice constant, or \SI{0.613}{eV} at the measured one. It reproduces the measured gap by construction, includes the zero-point renormalization through the measurement, and is accurate to about \SI{0.1}{eV} for other bands~\cite{godby1988self}. The hybrid functional HSE gives \SI{1.22}{eV} at its own lattice constant (\cref{tab:E-functionals}). A one-shot $G_0W_0$ calculation from LDA gives 1.24 to \SI{1.29}{eV} in the older calculations of \cref{tab:E-gw} and \SI{1.10}{eV} converged, from a static lattice. Two corrections separate it from the measurement: spin--orbit coupling, which these calculations omit, and the zero-point renormalization. Spin--orbit coupling splits the sixfold valence-band maximum (with spin) into a fourfold level at $+\Delta_{\mathrm{SO}}/3$ and a twofold level at $-2\Delta_{\mathrm{SO}}/3$; the weighted mean, $4\times\Delta_{\mathrm{SO}}/3 - 2\times2\Delta_{\mathrm{SO}}/3 = 0$, is the level without spin--orbit coupling. The gap therefore closes by $\Delta_{\mathrm{SO}}/3$, \SI{15}{meV} in silicon and, with $\Delta_{\mathrm{SO}} = \SI{0.296}{eV}$, \SI{99}{meV} in germanium (\cref{sec:3.3}). The zero-point renormalization lowers the gap of a real crystal relative to the static lattice by about \SI{60}{meV} and is computed separately (\cref{app:E.8}). For a dark-matter threshold in silicon, the three agree to about \SI{0.1}{eV}. That spread is the theory uncertainty of the threshold if the scissor is not anchored to a measurement, and \cref{sec:3.4} showed what it does to the minimum mass.
\end{workedexample}

What a quasiparticle calculation does not include, the interaction of the excited electron with its hole, is the subject of the Bethe--Salpeter equation of \cref{sec:4.9} and of \cref{app:E.6}.

\section{Density-functional perturbation theory in practice}\label{app:E.4}

Every number that \cref{sec:5.2,sec:5.3,sec:5.4,sec:5.5,sec:5.6} attach to the lattice of silicon, from the \SI{15.5}{THz} optical phonon at the zone center~\cite{nilsson1972study} to the Debye--Waller exponent at a momentum transfer of tens of keV, comes out of one kind of calculation: the first-order change of the Kohn--Sham potential when the nuclei of one species are displaced in a wave of crystal momentum $\vect{Q}$. \Cref{sec:5.4} derived the equations of density-functional perturbation theory (DFPT). Here we use them: the same linear-response equations that give a phonon also give, with an electric field or a strain in place of the displacement, the Born charges, $\varepsilon_\infty$ and the elastic constants, and, extended slightly, the changes of the empty states that a phonon-assisted amplitude needs (\cref{sec:13.2}).

\subsection{Outputs of a linear-response run}

One perturbation is labeled by a phonon wavevector $\vect{Q}$, a nucleus $\kappa$ and a Cartesian direction $\alpha$. It is the bare electron--nucleus potential change $\partial_{\kappa\alpha,\vect{Q}}V_{\mathrm{ext}}$ of a displacement wave $u^\alpha_{\vect{R}\kappa} \propto \ee^{\ii\vect{Q}\cdot\vect{R}}$. DFPT solves, for every occupied state $\psi_{v\vect{k}}$, the projected Sternheimer equation of \cref{sec:5.4},
\begin{equation}
  \bigl(H_{\mathrm{KS}} + a_{\mathrm{sh}}\op{P}^{\mathrm{occ}}_{\vect{k}+\vect{Q}} - \varepsilon_{v\vect{k}}\bigr)\,|\partial\psi_{v\vect{k}}\rangle
  = -\bigl(1 - \op{P}^{\mathrm{occ}}_{\vect{k}+\vect{Q}}\bigr)\,\partial_{\kappa\alpha,\vect{Q}}V_{\mathrm{KS}}\,|\psi_{v\vect{k}}\rangle ,
  \label{eq:E.4-sth}
\end{equation}
where $\op{P}^{\mathrm{occ}}_{\vect{k}+\vect{Q}} = \sum_{v'}|\psi_{v',\vect{k}+\vect{Q}}\rangle\langle\psi_{v',\vect{k}+\vect{Q}}|$ projects on the occupied states at $\vect{k}+\vect{Q}$, and the shift $a_{\mathrm{sh}}$, any energy larger than the valence bandwidth, makes the operator on the left positive definite without changing the solution~\cite[Eqs.~(30)--(31)]{baroni2001phonons}, \cite[Eq.~(48)]{giustino2017electron}. The solution $|\partial\psi_{v\vect{k}}\rangle$ has crystal momentum $\vect{k}+\vect{Q}$ and lies entirely in the empty manifold. From it the first-order density $\partial n_{\kappa\alpha,\vect{Q}}(\vect{r})$ is a sum of $\psi^*_{v\vect{k}}\,\partial\psi_{v\vect{k}}$ over occupied states and $\vect{k}$~\cite[Eq.~(49)]{giustino2017electron}, and the screened potential is
\begin{equation}
  \partial_{\kappa\alpha,\vect{Q}}V_{\mathrm{KS}} = \partial_{\kappa\alpha,\vect{Q}}V_{\mathrm{ext}} + v_{\mathrm{C}}\star\partial n_{\kappa\alpha,\vect{Q}} + \frac{\delta v_{\mathrm{xc}}}{\delta n}\,\partial n_{\kappa\alpha,\vect{Q}} ,
  \label{eq:E.4-dV}
\end{equation}
with $v_{\mathrm{C}}\star$ the convolution with the Coulomb kernel and $\delta v_{\mathrm{xc}}/\delta n$ the static exchange--correlation kernel of the ground-state functional. \Cref{eq:E.4-sth,eq:E.4-dV} are iterated to self-consistency, exactly as a ground-state calculation is, and the cost of one perturbation is of the same order as that of the ground-state density~\cite[Sect.~II.C]{baroni2001phonons}.

A converged run returns the screened potential, from which the vertex follows; the dynamical matrix; and, at $\vect{Q} = 0$, the response to a uniform field. The screened potential $\partial_{\kappa\alpha,\vect{Q}}V_{\mathrm{KS}}$ comes for all $3N_{\mathrm{at}}$ perturbations at this $\vect{Q}$, with $N_{\mathrm{at}}$ the number of atoms per cell, and its matrix elements between Bloch states, weighted by the polarization vectors, are the electron--phonon vertex $g_{mn\nu}(\vect{k},\vect{Q})$ of \cref{eq:A.6-g}. The dynamical matrix $D^{\mathrm{dyn}}(\vect{Q})$ of \cref{eq:A.6-Ddyn} is a second derivative of the energy that needs only the first-order density: by the $2n+1$ theorem of \cref{sec:5.4}, the first-order wavefunctions fix the energy to third order. Diagonalizing it yields $\omega_{\vect{Q}\nu}$ and $\vect{e}_{\kappa\nu}(\vect{Q})$. The response to a uniform electric field gives the electronic dielectric tensor $\varepsilon_\infty$ and the Born effective charges $Z^*_\kappa$, which we take up next.

\subsection{Place on the hierarchy}

On the truncation label $\mathfrak{T}$ of \cref{sec:8.5}, a DFPT phonon is a point with $n_B = 0$ and $n_A = 0$. The nuclei are clamped while the electrons relax: the Born--Oppenheimer surface is expanded to second order about its minimum, and the electrons respond statically, as if the displacement were frozen. The electronic part of the response is the static limit of the time-dependent density-functional response of \cref{sec:4.6}, with the adiabatic kernel of whatever functional the ground state used. Within density-functional theory this static response is formally exact for the exact functional, so the error of a DFPT phonon is the error of the functional, not of a perturbative truncation. As a result, DFPT frequencies inherit the strengths and the failures of the ground state (\cref{app:E.1}). The calculation is at zeroth thermal order, $n_T = 0$, in the sense that the lattice constant, the force constants and the electronic screening are those of the static lattice; thermal expansion and the temperature dependence of the frequencies need the quasi-harmonic and anharmonic methods of \cref{app:E.8}. The vertex $g$ from \cref{eq:E.4-dV} is the first-order coupling, $n_g = 1$ in the hierarchy; a rate that contains $|g|^2$ sits at $n_g = 2$.

\subsection{Born charges and dielectric tensors}

A uniform electric field $\vect{\mathcal{E}}$ cannot be written as a periodic potential, since $-e\vect{\mathcal{E}}\cdot\vect{r}$ grows without bound. DFPT avoids the problem because \cref{eq:E.4-sth} needs only matrix elements of the perturbation between different eigenstates, and those of $\vect{r}$ are well defined through the commutator,
\begin{equation}
  \langle\psi_{c\vect{k}}|\vect{r}|\psi_{v\vect{k}}\rangle = \frac{\langle\psi_{c\vect{k}}|[H_{\mathrm{KS}},\vect{r}]|\psi_{v\vect{k}}\rangle}{\varepsilon_{c\vect{k}} - \varepsilon_{v\vect{k}}} ,
  \label{eq:E.4-rcomm}
\end{equation}
which is the relation between the position and velocity matrix elements of \cref{sec:3.2}~\cite[Eq.~(37)]{baroni2001phonons}. For a local potential $[H_{\mathrm{KS}},\vect{r}] = -\ii\vect{p}/m_e$; a nonlocal pseudopotential (\cref{sec:3.5,app:E.2}) adds its own commutator, and omitting it biases every quantity built from the optical limit, the dielectric constant first among them. The projected state $\op{P}^{\mathrm{c}}\vect{r}|\psi_{v\vect{k}}\rangle$ is itself the solution of a Sternheimer equation with $[H_{\mathrm{KS}},\vect{r}]|\psi_{v\vect{k}}\rangle$ on the right-hand side~\cite[Eq.~(44)]{baroni2001phonons}. With it, the field enters \cref{eq:E.4-sth} like any other perturbation, and its self-consistent solution is the induced macroscopic polarization $\vect{P}$.

The macroscopic quantities then follow by differentiation. In the Heaviside--Lorentz units of the book (\cref{app:A.2}), with $\vect{D} = \vect{\mathcal{E}} + \vect{P}$,
\begin{equation}
  \varepsilon_\infty^{\alpha\beta} = \delta_{\alpha\beta} + \left.\frac{\partial P^\alpha}{\partial\mathcal{E}^\beta}\right|_{\vect{u} = 0},
  \qquad
  e\,Z^{*,\alpha\beta}_\kappa = \Omega\left.\frac{\partial P^\beta}{\partial u^\alpha_\kappa}\right|_{\vect{\mathcal{E}} = 0}
  = \left.\frac{\partial F^\alpha_\kappa}{\partial\mathcal{E}^\beta}\right|_{\vect{u} = 0} ,
  \label{eq:E.4-epsZ}
\end{equation}
where $\Omega$ is the cell volume and $\vect{u}_\kappa$ a uniform displacement of sublattice $\kappa$; Ref.~\cite[Eqs.~(94)--(96)]{baroni2001phonons} writes the same relations in Gaussian units, with $4\pi$ in front of $\partial P/\partial\mathcal{E}$. The equality of the two forms of $Z^*_\kappa$ in \cref{eq:E.4-epsZ}, as the polarization per displacement and as the force per field, holds because both are the mixed second derivative of one energy, so a code may compute whichever is cheaper. The quantities $\varepsilon_\infty$ and $Z^*_\kappa$ are all that the nonanalytic part of the dynamical matrix needs (\cref{sec:5.5}): in Heaviside--Lorentz units
\begin{equation}
  D^{\mathrm{dyn,na},\alpha\beta}_{\kappa\kappa'}(\vect{Q}\to 0) = \frac{e^2}{\Omega\sqrt{M_\kappa M_{\kappa'}}}\,
  \frac{(Z^*_\kappa\cdot\vect{Q})_\alpha\,(Z^*_{\kappa'}\cdot\vect{Q})_\beta}{\vect{Q}\cdot\varepsilon_\infty\cdot\vect{Q}} ,
  \label{eq:E.4-na}
\end{equation}
which is Ref.~\cite[Eq.~(98)]{baroni2001phonons} with $4\pi e^2_{\mathrm{G}} \to e^2$ and the mass factors of \cref{eq:A.6-Ddyn}. The same two tensors build the long-range vertex $g^{\mathrm{L}}$ of \cref{sec:6.2} and of \cref{app:E.5}.

Measured frequencies give an independent check on $Z^*$. For a cubic crystal with two atoms per cell, $Z^*_1 = -Z^*_2 = Z^*$ times the identity, and the optical mode at $\vect{Q}\to0$ has $\vect{e}_{1} = (M_2/M_{\mathrm{cell}})^{1/2}\hat{\vect{e}}$ and $\vect{e}_{2} = -(M_1/M_{\mathrm{cell}})^{1/2}\hat{\vect{e}}$. Its expectation value of \cref{eq:E.4-na} is $(e^2Z^{*2}/\Omega\varepsilon_\infty)\,[e_1/\sqrt{M_1} - e_2/\sqrt{M_2}]^2 = e^2Z^{*2}/(\Omega\varepsilon_\infty\mu_{\mathrm{red}})$, with $\mu_{\mathrm{red}} = M_1M_2/M_{\mathrm{cell}}$ the reduced mass of the pair, for a longitudinal mode and zero for a transverse one. Hence
\begin{equation}
  \omega_{\mathrm{L}}^2 - \omega_{\mathrm{T}}^2 = \frac{e^2Z^{*2}}{\Omega\,\varepsilon_\infty\,\mu_{\mathrm{red}}} ,
  \label{eq:E.4-LOTO}
\end{equation}
and with the measured $\omega_{\mathrm{L}}$, $\omega_{\mathrm{T}}$ and $\varepsilon_\infty$ (or $\varepsilon_0$, through the Lyddane--Sachs--Teller relation of \cref{sec:5.5}) it fixes $|Z^*|$. Where a rate depends on a near-cancellation, as the electron-coupled scalar in GaAs does through the vertex $Z_\kappa - Z^*_\kappa$ (\cref{sec:11.8}), this check matters more than any convergence parameter. A rate $\propto (Z_\kappa - Z^*_\kappa)^2$ has relative error $2\,\delta Z^*_\kappa/|Z_\kappa - Z^*_\kappa|$, so a one-percent error in $Z^*_\kappa$ becomes a large error in the rate when $|Z_\kappa - Z^*_\kappa|$ is a small fraction of $Z^*_\kappa$. The computed Born charge and $\varepsilon_\infty$ must then be tested against \cref{eq:E.4-LOTO} with measured frequencies before they are used.

The Berry-phase theory of polarization (\cref{sec:5.5}) is the other route to $Z^*_\kappa$. The electronic polarization of an insulator is
\begin{equation}
  \vect{P}_{\mathrm{el}} = -\frac{e\,g_{\mathrm{deg}}}{(2\pi)^3}\sum_{v}\int_{\mathrm{BZ}}\dd^3k\,\langle u_{v\vect{k}}|\ii\nabla_{\vect{k}}|u_{v\vect{k}}\rangle ,
  \label{eq:E.4-berry}
\end{equation}
with $g_{\mathrm{deg}} = 2$ for spin, the integrand a Berry connection of the periodic parts $u_{v\vect{k}}$, and the result equal to $-e g_{\mathrm{deg}}/\Omega$ times the sum of the Wannier centers in one cell~\cite[Eqs.~(89)--(90)]{marzari2012maximally}. On a grid of $\vect{k}$ the derivative is replaced by the phase of a product of overlaps $\langle u_{v\vect{k}}|u_{v'\vect{k}+\Delta\vect{k}}\rangle$ along strings of points, which is gauge invariant. Since the choice of which Wannier center belongs to the home cell is arbitrary, $\vect{P}_{\mathrm{el}}$ is defined only modulo the quantum $e g_{\mathrm{deg}}\vect{R}/\Omega$ for a lattice vector $\vect{R}$~\cite[Sect.~V.A.1]{marzari2012maximally}. A Born charge is a difference of two polarizations at small displacements of opposite sign, $\pm\delta u$, and the two must sit on the same branch of the multivalued function; with $\delta u$ a small fraction of a bohr this is automatic. Linear response and finite differences of \cref{eq:E.4-berry} agree within numerical accuracy~\cite[Sect.~II.D.2]{baroni2001phonons}, and the Berry-phase route is the one that also works at finite fields and for ferroelectrics.

The acoustic sum rule for charges and the measured $\varepsilon_\infty$ cost nothing to check. The sum rule, $\sum_\kappa Z^*_\kappa = 0$ (\cref{sec:5.5}), holds exactly only at convergence; its violation on a finite $\vect{k}$ grid is the cleanest measure of how well $\varepsilon_\infty$ and $Z^*$ are converged. In the comparison of $\varepsilon_\infty$ with experiment the functional dominates. Local and semilocal functionals put the gap too low (\cref{sec:3.4}), the energy denominators of \cref{eq:E.4-rcomm} come out too small, and the material is overscreened. The direction follows from the static limit of the Kramers--Kronig relation of \cref{sec:4.3},
\begin{equation}
  \varepsilon_1(\omega = 0) - 1 = \frac{2}{\pi}\int_0^\infty\frac{\varepsilon_2(\omega')}{\omega'}\,\dd\omega' :
  \label{eq:E.4-kk}
\end{equation}
shifting the absorption spectrum $\varepsilon_2$ down in energy by a fraction of an eV, at roughly fixed weight, increases the integral. For silicon, a DFT calculation finds $\varepsilon_\infty = 12.89$~\cite{ponce2018towards}, against the measured room-temperature value of 11.7~\cite{franta2017temperature} (\cref{sec:4.4}) used throughout this book. When $\varepsilon_\infty$ enters a rate through \cref{eq:E.4-na} or the Fr\"ohlich vertex, the measured value is usually the better input; Ref.~\cite{ponce2018towards} rescales the screening of the electron--phonon vertex of silicon to the measured value it adopts, 11.94, for exactly this reason.

\subsection{Force-constant interpolation}

DFPT gives $D^{\mathrm{dyn}}(\vect{Q})$ at any $\vect{Q}$, but a dispersion, a density of states or a Debye--Waller sum needs it at thousands of points. The standard route computes it on a uniform $N_Q\times N_Q\times N_Q$ grid, transforms to real-space force constants $\Phi^{(2)}_{\kappa\kappa'}(\vect{0},\vect{R})$ by the inverse of \cref{eq:A.6-Ddyn},
\begin{equation}
  \Phi^{(2),\alpha\beta}_{\kappa\kappa'}(\vect{0},\vect{R}) = \frac{\sqrt{M_\kappa M_{\kappa'}}}{N_Q^3}\sum_{\vect{Q}\in\mathrm{grid}} D^{\mathrm{dyn},\alpha\beta}_{\kappa\kappa'}(\vect{Q})\,\ee^{-\ii\vect{Q}\cdot\vect{R}} ,
  \label{eq:E.4-ifc}
\end{equation}
and transforms back at any $\vect{Q}$~\cite[Eq.~(99)]{baroni2001phonons}. The discrete transform of \cref{eq:E.4-ifc} returns force constants that are periodic on the $N_Q^3$ supercell, so it is exact only if the true constants vanish beyond that supercell. In practice the vectors $\vect{R}$ are folded into the Wigner--Seitz cell of the supercell, with weights shared among equidistant images, and the grid is chosen large enough that the force constants have decayed at its edge. A rule of thumb is to include neighbors out to two or three bond lengths and to check the interpolated frequencies against direct DFPT runs at a few $\vect{Q}$ off the grid~\cite[Sect.~II.D.3]{baroni2001phonons}. In diamond and zincblende crystals the force constants are long ranged along the $\langle 110\rangle$ chains of bonds~\cite[Sect.~V.A.1]{baroni2001phonons}, so silicon needs larger grids than its short bonds suggest.

Polar crystals need one more step. The nonanalytic term \cref{eq:E.4-na} decays in real space as $R^{-3}$, the dipole--dipole interaction, which no finite supercell captures. It is subtracted from the computed $D^{\mathrm{dyn}}(\vect{Q})$ in its Ewald-summed form, the short-range remainder is interpolated, and the dipole term is added back analytically at the target $\vect{Q}$; this restores the LO--TO splitting and its dependence on the direction of $\vect{Q}\to 0$ (\cref{sec:5.5}).

The acoustic sum rule of \cref{sec:5.2}, $\sum_{\vect{R}\kappa'}\Phi^{(2)}_{\kappa\kappa'}(\vect{0},\vect{R}) = 0$, holds for the exact force constants but not for the computed ones. The exchange--correlation potential is evaluated on a real-space grid that is not invariant under arbitrary translations, so a rigid shift of the crystal changes the energy slightly. Codes impose the rule after the fact, by subtracting the violation from the on-site constants or by a minimal projection. The size of the correction is a diagnostic: the acoustic frequencies at $\Gamma$ before it should be small compared with the lowest frequency the application needs, and a value that grows instead of shrinking with the cutoff signals an unconverged real-space grid. The correction is applied to $D^{\mathrm{dyn}}$ and to $Z^*$, never as a zeroing of $g$ at $\vect{Q}\to 0$: the acoustic vertex vanishes there by its own sum rule (\cref{sec:6.2}), and its leading term, linear in $\vect{Q}$, is physical.

The cost argument for DFPT over frozen phonons is a scaling law. A frozen-phonon calculation of the force constants out to a range $R_{\mathrm{IFC}}$ (in units of the cell size) needs supercells of $\sim R_{\mathrm{IFC}}^3$ cells, $N_{\mathrm{at}}R_{\mathrm{IFC}}^3$ atoms; with ground-state cost cubic in the number of atoms and $3N_{\mathrm{at}}$ independent displacements, the total scales as $3N_{\mathrm{at}}^4R_{\mathrm{IFC}}^9$. DFPT works in the primitive cell on a grid of $\sim R_{\mathrm{IFC}}^3$ wavevectors, each costing $\sim N_{\mathrm{at}}^3$ per perturbation, which totals $\sim 3N_{\mathrm{at}}^4R_{\mathrm{IFC}}^3$~\cite[Sect.~IV.B]{baroni2001phonons}, a factor $R_{\mathrm{IFC}}^6$ less. For a two-atom cell and force constants that extend several shells, the difference is decisive.

\subsection{Displacement correlations from DFPT phonons}

The Debye--Waller exponent $2W_\kappa(\vect{q}) = \vect{q}\cdot B_\kappa\cdot\vect{q}$ and the displacement correlation $\mathcal{C}_{\kappa\kappa'}(t)$ of \cref{sec:5.6} are sums over all phonons of the zone. Inserting the mode expansion \cref{eq:A.6-u} into $B_\kappa = \langle\vect{u}_\kappa\vect{u}_\kappa\rangle$ and using $\langle B_{\vect{Q}\nu}B_{-\vect{Q}\nu}\rangle = 2n_{\vect{Q}\nu} + 1$ in the harmonic thermal state,
\begin{equation}
  B^{\alpha\beta}_\kappa = \frac{1}{N}\sum_{\vect{Q}\nu}\frac{2n_{\vect{Q}\nu} + 1}{2M_\kappa\,\omega_{\vect{Q}\nu}}\,
  e^\alpha_{\kappa\nu}(\vect{Q})\,e^{\beta}_{\kappa\nu}(\vect{Q})^* ,
  \label{eq:E.4-B}
\end{equation}
and $\mathcal{C}_{\kappa\kappa'}(t)$ is the same sum with the phases $\ee^{\mp\ii\omega_{\vect{Q}\nu}t}$ on the emission and absorption parts and the cross products $e_{\kappa\nu}e^*_{\kappa'\nu}$. These are post-processing sums: the frequencies and polarization vectors come from the interpolated force constants on a fine grid, and the cost is negligible. The approximation is $n_A = 0$, with the thermal occupation exact for the harmonic crystal; at high temperature or for light, soft lattices the anharmonic methods of \cref{app:E.8} replace the harmonic phonons in \cref{eq:E.4-B}.

The catch is the zone center, where the acoustic frequencies vanish as $\omega = c_\nu Q$ and the summand of \cref{eq:E.4-B} diverges. The divergence is integrable in three dimensions, but a uniform grid samples it badly. How badly follows from replacing the sum over the grid cell around $\Gamma$ by an integral over a sphere of the same volume, of radius $Q_c$ with $(4\pi/3)Q_c^3 = (2\pi)^3/(N_Q^3\Omega)$. With the acoustic polarization $\vect{e}_{\kappa\nu} = (M_\kappa/M_{\mathrm{cell}})^{1/2}\hat{\vect{e}}_\nu$ of \cref{app:A.6}, an isotropic average over directions, and $\sum_\nu$ over the three acoustic branches, the contribution of that sphere to one diagonal component is
\begin{equation}
  \Delta B^{\alpha\alpha} = \frac{\Omega}{2\pi^2}\int_0^{Q_c}\!Q^2\,\dd Q\;\frac13\sum_\nu\frac{2n_{\vect{Q}\nu}+1}{2M_{\mathrm{cell}}c_\nu Q}
  \simeq
  \begin{cases}
    \dfrac{\Omega\,Q_c^2}{24\pi^2M_{\mathrm{cell}}}\displaystyle\sum_\nu\frac{1}{c_\nu}, & T = 0,\\[2ex]
    \dfrac{\Omega\,Q_c\,T}{6\pi^2M_{\mathrm{cell}}}\displaystyle\sum_\nu\frac{1}{c_\nu^2}, & T \gg c_\nu Q_c,
  \end{cases}
  \label{eq:E.4-dBgamma}
\end{equation}
where the high-temperature form uses $2n + 1 \simeq 2T/\omega$. Since $Q_c \propto 1/N_Q$, the error of a calculation that simply omits the $\Gamma$ cell falls as $N_Q^{-2}$ at zero temperature but only as $N_Q^{-1}$ at room temperature. For silicon, with $\Omega = a^3/4 = \SI{40.05}{\angstrom^3}$, $M_{\mathrm{cell}}$ two silicon masses, a representative sound speed of $\SI{6}{km.s^{-1}}$ for all three branches (the range is 3 to \SI{10}{km.s^{-1}} across materials, \cref{sec:1.3}) and an $8^3$ grid, $Q_c = \SI{0.142}{\angstrom^{-1}}$ and the room-temperature term is $\Delta B^{\alpha\alpha} \approx \SI{3.6e-4}{\angstrom^2}$, while the zero-temperature term is twenty times smaller, \SI{1.9e-5}{\angstrom^2}. (The high-temperature form applies since $c\,Q_c$ corresponds to \SI{65}{K}.) The remedy is to add \cref{eq:E.4-dBgamma}, evaluated with the computed sound speeds, as an analytic correction, or to use a grid fine enough that the residual is below the target accuracy. The same care applies to $\mathcal{C}_{\kappa\kappa'}(t)$ at long times, which are controlled by the smallest frequencies.

\subsection{Wavefunction changes beyond the density}

The vertex of \cref{eq:A.6-g} needs only $\partial_{\kappa,\vect{Q}}V_{\mathrm{KS}}$ and two Bloch states, so every DFPT run already contains it~\cite[Sect.~III.B.3]{giustino2017electron}. Codes write the potentials to disk for each irreducible $\vect{Q}$ of the coarse grid, one complex function on the real-space grid per perturbation, and a later step forms the matrix elements and interpolates them (\cref{app:E.5}). The storage is modest by modern standards: for a two-atom cell, six perturbations per $\vect{Q}$, each a complex function on a real-space grid of $\sim 10^5$ points, are about \SI{10}{MB} per wavevector, so a few hundred irreducible wavevectors fit in a few gigabytes.

A transition amplitude needs more than the density does. The density vertex of \cref{sec:3.2}, $\rho_{cv}(\vect{k},\vect{q}) = \langle\psi_{c,\vect{k}+\vect{q}}|\ee^{\ii\vect{q}\cdot\vect{r}}|\psi_{v\vect{k}}\rangle$, changes under a frozen displacement of crystal momentum $\vect{Q}$ by
\begin{equation}
  \partial\rho_{cv} = \langle\partial\psi_{c}|\ee^{\ii\vect{q}\cdot\vect{r}}|\psi_{v}\rangle + \langle\psi_{c}|\ee^{\ii\vect{q}\cdot\vect{r}}|\partial\psi_{v}\rangle ,
  \label{eq:E.4-drho}
\end{equation}
which is the adiabatic (Herzberg--Teller) form of the phonon-assisted amplitude of \cref{sec:13.2}, valid to relative order $\omega_{\vect{Q}\nu}/\delta_{\mathrm{off}}$, with $\delta_{\mathrm{off}}$ the energy by which the intermediate states are off shell (\cref{sec:13.2}). Both terms need the complete first-order change of an individual state, while \cref{eq:E.4-sth} returns only part of one of them. The full change of an occupied state is
\begin{equation}
  |\partial\psi_{v\vect{k}}\rangle = \sum_{m\notin\mathrm{deg}(v\vect{k})}|\psi_{m,\vect{k}+\vect{Q}}\rangle\,
  \frac{\langle\psi_{m,\vect{k}+\vect{Q}}|\partial_{\kappa,\vect{Q}}V_{\mathrm{KS}}|\psi_{v\vect{k}}\rangle}{\varepsilon_{v\vect{k}} - \varepsilon_{m,\vect{k}+\vect{Q}}} ,
  \label{eq:E.4-dpsi}
\end{equation}
with the states degenerate with $\psi_{v\vect{k}}$ excluded. DFPT computes the part with $m$ empty, by \cref{eq:E.4-sth}, without ever forming the sum. The part with $m$ occupied, $\op{P}^{\mathrm{occ}}_{\vect{k}+\vect{Q}}|\partial\psi_{v\vect{k}}\rangle$, is left out on purpose: these terms drop out of the density change once the $\vect{Q}$ and $-\vect{Q}$ components of the Hermitian perturbation are combined, since exchanging the roles of the two occupied states flips the sign of the denominator~\cite[Sect.~II.C]{baroni2001phonons}, \cite[Eq.~(48)]{giustino2017electron}. In \cref{eq:E.4-drho} they do not cancel, because $\psi_c$ singles out one state. They are, however, cheap. The sum runs over the few occupied bands at $\vect{k}+\vect{Q}$, and the numerators are the occupied--occupied matrix elements of the potential that DFPT has already stored. Two cautions apply. The denominators $\varepsilon_{v\vect{k}} - \varepsilon_{v',\vect{k}+\vect{Q}}$ vanish wherever two valence bands at $\vect{k}$ and $\vect{k}+\vect{Q}$ cross, and near such crossings the adiabatic form fails: there the dynamical denominators $\varepsilon_{v\vect{k}} - \varepsilon_{v',\vect{k}+\vect{Q}} \pm \omega_{\vect{Q}\nu}$ of the full amplitude (\cref{sec:13.2,sec:13.6}) must be used instead. And at $\vect{Q} = 0$ the degenerate subspace has to be treated in the same gauge as the states that enter $\rho_{cv}$ (\cref{app:A.5}).

The change of an empty state is not produced by standard DFPT at all, since empty states do not contribute to the density. Its equation is a Sternheimer equation with the converged potential on the right,
\begin{equation}
  \bigl(H_{\mathrm{KS}} - \varepsilon_{c\vect{k}}\bigr)|\partial\psi_{c\vect{k}}\rangle = -\bigl(1 - \op{P}^{\mathrm{deg}}_{c}\bigr)\,\partial_{\kappa,\vect{Q}}V_{\mathrm{KS}}\,|\psi_{c\vect{k}}\rangle ,
  \label{eq:E.4-dpsic}
\end{equation}
with $\op{P}^{\mathrm{deg}}_{c}$ the projector on the states at $\vect{k}+\vect{Q}$ degenerate with $\psi_{c\vect{k}}$. No self-consistency is needed, because $\partial V_{\mathrm{KS}}$ is already the screened potential of the occupied manifold, and that is the main saving. The operator on the left, however, is indefinite: $\varepsilon_{c\vect{k}}$ lies inside the spectrum of $H_{\mathrm{KS}}$, with all the valence states and part of the conduction bands below it, and the conjugate-gradient solvers used for \cref{eq:E.4-sth} do not converge on an indefinite operator. One way out splits the space: a window $\mathcal{W}$ of bands at $\vect{k}+\vect{Q}$, containing every state below $\varepsilon_{c\vect{k}}$ and the conduction bands up to some energy above it, is summed explicitly with \cref{eq:E.4-dpsi}, where the Wannier-interpolated vertices of \cref{app:E.5} make the terms cheap. The remainder obeys
\begin{equation}
  \op{P}^{\perp}_{\mathcal{W}}\bigl(H_{\mathrm{KS}} - \varepsilon_{c\vect{k}}\bigr)\op{P}^{\perp}_{\mathcal{W}}\,|\partial\psi^{\perp}_{c\vect{k}}\rangle = -\op{P}^{\perp}_{\mathcal{W}}\,\partial_{\kappa,\vect{Q}}V_{\mathrm{KS}}\,|\psi_{c\vect{k}}\rangle ,
  \label{eq:E.4-dpsiperp}
\end{equation}
with $\op{P}^{\perp}_{\mathcal{W}}$ the projector on the complement of the window. Every state left in the complement lies above $\varepsilon_{c\vect{k}}$, so the projected operator is positive definite and the standard solvers apply. The remainder carries the high-energy intermediate states, which matter for large momentum transfers, where $\ee^{\ii\vect{q}\cdot\vect{r}}$ couples to bands far above the gap. The window is a convergence parameter: the sum of the explicit and the Sternheimer parts must not depend on it.

The same equations at finite frequency give dynamical response without empty states. For a perturbation oscillating as $\ee^{-\ii\omega t}$, first-order time-dependent perturbation theory (\cref{sec:2.4}) leads to two components of the change of each occupied state,
\begin{equation}
  \bigl(H_{\mathrm{KS}} - \varepsilon_{v\vect{k}} \mp \omega\bigr)|\partial\psi^{\pm}_{v\vect{k}}\rangle = -\bigl(1 - \op{P}^{\mathrm{occ}}\bigr)\,\partial V_{\mathrm{KS}}(\pm\omega)\,|\psi_{v\vect{k}}\rangle ,
  \label{eq:E.4-sthw}
\end{equation}
with the density change built from both and the potential updated as in \cref{eq:E.4-dV} with an adiabatic kernel. Below the gap the operators remain positive definite on the empty manifold, because $\varepsilon_c - \varepsilon_v - \omega > 0$ for every pair, and the solution is as cheap as a static one. Above the gap a small broadening $\eta_{\mathrm{b}}$, $\omega \to \omega + \ii\eta_{\mathrm{b}}$, is added and the linear systems become complex and non-Hermitian. Frequency-dependent Sternheimer calculations of this kind give the inverse frequency moments and the drag mass of \cref{sec:12.5}, and the screened interaction in Sternheimer formulations of the methods of \cref{app:E.3}.

\subsection{Elastic constants}

The sound speeds of \cref{sec:11.1} follow from the elastic tensor $C_{ijkl}$ through the Christoffel equation of \cref{app:C.18}; for a cubic crystal along $[100]$ it reduces to $\rho_T c_{\mathrm{LA}}^2 = C_{11}$ and $\rho_T c_{\mathrm{TA}}^2 = C_{44}$ in Voigt notation. Three routes give $C_{ijkl}$. The slopes of the interpolated acoustic branches are the cheapest and the least accurate, since the small-$Q$ part of the dispersion is precisely where the Fourier interpolation of \cref{eq:E.4-ifc} is weakest, and the slopes converge only slowly with $N_Q$. Finite strains are the most reliable: strains of a few tenths of a percent of both signs are applied, the internal coordinates are relaxed at each strain, and the stress is fitted. The internal relaxation matters for any structure with internal degrees of freedom, and in the diamond structure it affects $C_{44}$ because the two sublattices shift against each other under shear. Since relaxation can only lower the energy at fixed strain, the relaxed constants are never larger than the clamped ones. DFPT with strain as the perturbation returns the clamped-ion tensor and the internal-strain couplings in one run, from which the relaxed tensor follows by eliminating the internal displacements~\cite[Sect.~II.E.1]{baroni2001phonons}.

Codes differ in the order of the shear rows. Voigt notation assigns $4 \leftrightarrow yz$, $5 \leftrightarrow xz$, $6 \leftrightarrow xy$, but several codes and databases store the shear components in the order $xy, yz, zx$. For a cubic crystal the error is invisible, since $C_{44} = C_{55} = C_{66}$. For a lower-symmetry crystal the permuted tensor still looks reasonable, but the sound speeds along the low-symmetry directions come out wrong. A cheap check is to compute $\rho_Tc^2$ from the tensor along one direction and compare it with the slope of the computed dispersion along the same direction.

\subsection{Convergence parameters and cost}

Five numerical parameters control a DFPT calculation. Each must be converged on the quantity of interest, not on the total energy, and with a tolerance fixed before the scan. The plane-wave cutoff of \cref{app:E.1} sets the basis; the test is the frequency at $\Gamma$ and at one zone-boundary point at two cutoffs, not the total energy. The $\vect{k}$ grid sets the Brillouin-zone integrals in \cref{eq:E.4-sth}; for an insulator the density and the frequencies converge on modest grids, but $\varepsilon_\infty$ and $Z^*$ need noticeably finer ones, because the response to a field involves the velocity matrix elements and energy denominators of \cref{eq:E.4-rcomm} at every $\vect{k}$, and these converge slowly where the direct gap is small~\cite[Sect.~II.C.2]{baroni2001phonons}. Metals need smearing of the occupations (\cref{app:E.1}) and a correspondingly finer grid. The $\vect{Q}$ grid sets the range of the interpolated force constants and is tested by direct runs off the grid. The self-consistency threshold of \cref{eq:E.4-dV} must be much tighter than in a ground-state run. An error in $\partial n$ enters the force constants linearly, and where a frequency is small, as for the acoustic modes near $\Gamma$ or a soft mode, a small error in the force constant is a large relative error in the frequency, or makes it imaginary.

The lattice constant is a parameter in its own right. Phonons are usually computed at the theoretical equilibrium of the functional. In diamond and zincblende crystals the forces vanish by symmetry at any lattice constant, but only at the equilibrium does the stress vanish, so that the second derivatives are taken at a minimum of the same energy surface. Each frequency moves with the lattice constant through its Gr\"uneisen parameter (\cref{app:E.8}); the generalized-gradient approximation, which expands lattice constants, lowers the branches with positive Gr\"uneisen parameter, and for silicon the local-density approximation already reproduces the measured dispersion closely while the generalized-gradient approximation slightly worsens it~\cite[Sect.~V.A]{baroni2001phonons}.

The cost of a complete phonon calculation is the number of irreducible wavevectors on the $\vect{Q}$ grid, times $3N_{\mathrm{at}}$ perturbations (fewer, by symmetry, at high-symmetry wavevectors), times the cost of a ground-state calculation with the reduced symmetry of the perturbed crystal. That reduced symmetry is the hidden factor: a perturbation at a general $\vect{Q}$ leaves few symmetry operations, so the $\vect{k}$ sum runs over many more points than in the ground state. An $N_Q^3$ grid has about $N_Q^3/48$ irreducible points for a cubic crystal when $N_Q$ is large; an $8^3$ grid for silicon has 29, a count that follows from applying the 48 operations of the cubic group to the grid points and counting the distinct orbits. With six perturbations per wavevector, fewer at the high-symmetry points, and the three field perturbations at $\Gamma$, a complete silicon calculation on that grid costs at most about 180 self-consistent linear-response solutions, each of the order of one ground-state calculation at reduced symmetry. The cost grows with the fourth power of the number of atoms per cell~\cite[Sect.~IV.B]{baroni2001phonons}.

\begin{workedexample}{A DFPT calculation of silicon}\label{we:E-dfpt-si}
We assemble the settings of two published silicon calculations and the checks they pass. The first used the local-density approximation with plane-wave pseudopotentials, a \SI{35}{Ry} cutoff and a $14^3$ grid for the ground-state density, and computed the dynamical matrices and the potential changes on a $6^3$ grid of wavevectors~\cite{noffsinger2012phonon}. The second used the generalized-gradient approximation at \SI{35}{Ry} (\SI{45}{Ry} for the local-density approximation), again a $6^3$ wavevector grid, and found that an $18^3$ grid of $\vect{k}$ points was needed for the Born charges to vanish~\cite{ponce2018towards}. A \SI{35}{Ry} cutoff corresponds to plane waves up to $|\vect{G}| = (2m_e\times\SI{476}{eV})^{1/2} = \SI{22}{keV}$, between five and six times the shortest reciprocal-lattice vector $|\vect{G}_{111}| = \sqrt3\,2\pi/a = \SI{3.95}{keV}$ (\cref{sec:1.3}), which is ample for phonons but, as \cref{app:E.6,app:E.7} explain, not automatically for a density vertex at large momentum transfer.

The vanishing of the Born charges is a symmetry statement that the numerics must reproduce. The two atoms of the diamond structure are exchanged by inversion through the bond center. $Z^*_\kappa$ is the derivative of a polar vector with respect to a polar vector, which is even under inversion, so $Z^*_1 = Z^*_2$. Each atom sits at a site of tetrahedral symmetry, so each tensor is a multiple of the identity. Charge neutrality, $Z^*_1 + Z^*_2 = 0$, then forces both to vanish. On a coarse $\vect{k}$ grid they do not, and since $Z^*$ enters \cref{eq:E.4-na}, a residual Born charge produces a spurious LO--TO splitting in a crystal that has none. The $\vect{k}$ grid is therefore fixed by the Born charge, not by the frequencies.

The dielectric constant is the second check, and here the functional limits the result: DFT gives $\varepsilon_\infty = 12.89$~\cite{ponce2018towards}, the measurement 11.7~\cite{franta2017temperature}. The PBE calculation of \cref{we:E-si-bands} reaches 13.04 at the measured lattice constant on an $18^3$ grid, within about one percent of that value. The overestimate, \SI{10}{\percent} for 12.89 and \SI{11}{\percent} for 13.04, is the overscreening of \cref{eq:E.4-kk}. It is not a convergence error, and it does not shrink with any numerical parameter. The frequencies are the third check. The measured optical frequency at $\Gamma$ is \SI{15.5}{THz}~\cite{nilsson1972study}, and the local-density dispersion agrees with neutron scattering across the zone~\cite[Sect.~V.A]{baroni2001phonons}. A calculation at the generalized-gradient lattice constant comes out soft for the branches with positive Gr\"uneisen parameter and hard for those with negative ones.
\end{workedexample}

The vertices for the coarse grids are the input of the Wannier interpolation in \cref{app:E.5}, which turns these few hundred potential changes into the millions of matrix elements that a rate needs.

\section{Wannier functions in practice}\label{app:E.5}

The first-principles phonon-assisted absorption spectrum of silicon was converged to an energy resolution of \SI{14}{meV} with $40\times40\times40$ grids for both the electron wavevector $\vect{k}$ and the phonon wavevector $\vect{Q}$~\cite{noffsinger2012phonon}. That is $64{,}000^2 \approx 4\times10^{9}$ pairs $(\vect{k},\vect{Q})$, each with a vertex $g_{mn\nu}(\vect{k},\vect{Q})$ for every pair of bands and every branch. A direct DFPT calculation (\cref{app:E.4}) costs about one ground-state calculation per perturbation and wavevector, so $64{,}000$ wavevectors are already out of reach, and the matrix elements between all those states even more so. The same study computed the vertices directly on a $6^3$ grid and produced the rest by Wannier interpolation. \Cref{sec:3.6} showed why such an interpolation is exact for quantities of finite range in the Wannier basis, and \cref{sec:6.2} why the electron--phonon vertex is such a quantity. In practice the work is in choosing and testing the Wannier functions and in keeping every interpolated quantity in one gauge; the result is a rate whose fine-grid vertices, \SI{270}{TB} of them, are never stored.

\subsection{Computed quantities}

A Wannier calculation turns quantities known on a coarse uniform grid into real-space matrices, and those back into the same quantities anywhere in the zone. For the bands, the coarse-grid Hamiltonian in the rotated (Wannier) gauge, $U(\vect{k})^\dagger\,\mathrm{diag}(\varepsilon_{n\vect{k}})\,U(\vect{k})$, is transformed to the real-space Hamiltonian $H^W(\vect{R})$ of \cref{sec:3.6}, and at any $\vect{k}$ the matrix $\sum_{\vect{R}}\ee^{\ii\vect{k}\cdot\vect{R}}H^W(\vect{R})$ is diagonalized. Its eigenvalues are the interpolated bands, and its eigenvectors define $U(\vect{k})$ on the fine grid. For the electron--phonon vertex, the coarse-grid vertices are rotated with the coarse-grid $U$ matrices and the phonon displacement patterns are undone, which leaves a matrix in two lattice vectors, one for the electron hop and one for the displaced atom~\cite[Eqs.~(177)--(179)]{giustino2017electron}. Back on a fine grid,
\begin{equation}
  g_{mn\nu}(\vect{k},\vect{Q}) = \sum_{\vect{R}_e,\vect{R}_p}\ee^{\ii(\vect{k}\cdot\vect{R}_e + \vect{Q}\cdot\vect{R}_p)}
  \sum_{m'n'\kappa\alpha}U_{mm'}(\vect{k}+\vect{Q})\,g^W_{m'n'\kappa\alpha}(\vect{R}_e,\vect{R}_p)\,U^\dagger_{n'n}(\vect{k})\,
  \frac{e^\alpha_{\kappa\nu}(\vect{Q})}{\sqrt{2M_\kappa\omega_{\vect{Q}\nu}}} ,
  \label{eq:E.5-ginterp}
\end{equation}
in the normalization of \cref{eq:A.6-g}. The phonon frequencies and polarization vectors on the fine grid come from the interpolated force constants of \cref{app:E.4}. The same construction interpolates the velocity matrix elements of \cref{sec:3.2}, which the optical limit needs, and, with a restriction discussed below, the density vertex $\rho_{mn}(\vect{k},\vect{q})$.

\subsection{Place on the hierarchy}

The interpolation is a numerical device, not a physical approximation, and it has no component on the truncation label $\mathfrak{T}$ of \cref{sec:8.5}. What the label records is the physics of the quantity being interpolated. The bands are independent-particle bands, $\mathcal{X} = \mathrm{IP}$, unless quasiparticle corrections are interpolated in their place (\cref{app:E.3}). The vertex is first order in the displacement, $n_g = 1$. Phonon-assisted absorption is $n_g = 2$ with $n_T = 1$, first order in the thermal phonon occupations $n_{\vect{Q}\nu}$ that weight its absorption branch, and $\mathcal{X} = \mathrm{IP}$, as \cref{sec:8.5} reads it off from the assumptions of \cref{sec:6.5}, and the phonon linewidth of \cref{sec:6.4} is the lowest order in $g$ of the phonon self-energy. The interpolation error belongs to the numerical part of the error budget of \cref{sec:8.8}, and it is tested and reported separately from the truncation.

\subsection{Projections and disentanglement}

Silicon has two textbook choices. The four valence bands form an isolated group, separated from the conduction bands by the gap at every $\vect{k}$, and they can be rotated into four Wannier functions centered on the bonds, each a bonding combination of two $sp^3$ hybrids on neighboring atoms~\cite[Fig.~2]{marzari2012maximally}. No disentanglement is needed, and these functions are the natural choice for the valence part of a polarization calculation (\cref{eq:E.4-berry}). A dark-matter rate needs the conduction bands as well. Eight $sp^3$ hybrids, four per atom, describe the valence bands and the lowest conduction bands, but the conduction bands of silicon cross higher bands of other character, so the eight-dimensional subspace must be extracted from a larger set at every $\vect{k}$. This is disentanglement~\cite{souza2001maximally}.

Two energy windows control it. The outer (disentanglement) window contains every band that may contribute to the subspace. The inner (frozen) window contains the states that must be reproduced exactly; inside it the Bloch states are kept unchanged, and only outside it does the procedure choose, at each $\vect{k}$, the combination of bands that varies most smoothly across the zone. Projection onto atomic $sp^3$ orbitals alone reproduces the bands of silicon well except where higher bands with some $sp^3$ character mix in; forcing the whole valence manifold into the frozen window removes those deviations~\cite[Fig.~5]{marzari2012maximally}. For an electron-recoil calculation the inner window should contain every state that can be a final state at the energies of interest, plus a margin, since a conduction state outside the frozen window is reproduced only approximately.

The number of Wannier functions grows with the energy range that must be described. For phonon-assisted absorption up to the direct gap and beyond, Ref.~\cite{noffsinger2012phonon} extracted 26 Wannier functions from 34 bands, reproducing the band structure from \SI{10}{eV} below to \SI{30}{eV} above the Fermi level. The range is set by the intermediate states. A second-order amplitude sums over all bands $m$ (\cref{sec:6.5,sec:13.2}), and the Wannier manifold truncates that sum; in \cref{app:E.4} we saw how the remainder above the manifold can be added by a projected Sternheimer equation.

Every set of Wannier functions passes four checks before use. The spread functional $\Omega_W$ of \cref{sec:3.6} must be converged and the individual spreads small compared with the cell. The centers must sit where the symmetry says, on the bonds or the atoms. For a time-reversal-symmetric crystal the maximally localized functions can be chosen real, and a large ratio of imaginary to real parts signals a poor minimum. Above all, the interpolated bands must agree with a direct calculation at points off the coarse grid, inside the frozen window, to the tolerance the application needs. In silicon the last check has a twist: the conduction minimum lies on the $\Delta$ line at about 0.85 of the way from $\Gamma$ to X, not at a high-symmetry point, so a coarse grid that does not resolve it misplaces the valley. Ref.~\cite{ponce2018towards} needed a $12^3$ coarse $\vect{k}$ grid for the conduction bands for this reason, while $6^3$ was enough for the phonon wavevectors.

\subsection{Interpolating vertices, and the long-range parts}

The decay of $g^W(\vect{R}_e,\vect{R}_p)$ decides whether \cref{eq:E.5-ginterp} converges. In the electron vector $\vect{R}_e$ it decays at least as fast as the Wannier functions themselves. In the atom vector $\vect{R}_p$ it decays as the screened potential of a displaced atom. In a nonpolar semiconductor or insulator that potential falls off at least as fast as a quadrupole, $R_p^{-3}$, so the vertices in reciprocal space stay finite as $\vect{Q}\to0$ and the interpolation converges~\cite[Sect.~VI.A.2]{giustino2017electron}. For diamond, coarse grids of $4^3$, $6^3$ and $8^3$ points interpolate the vertex of the highest optical branch along high-symmetry lines onto the direct DFPT values, with the $4^3$ grid already close~\cite[Fig.~6]{giustino2017electron}. For a new material, the same comparison, of $g$ from direct DFPT at a few $(\vect{k},\vect{Q})$ off the coarse grid with the interpolated values, is the only evidence that the interpolation converges.

In a polar crystal the displaced atom carries a dipole, its potential decays as $R_p^{-2}$, and the vertex of the longitudinal optical branch diverges as $1/Q$: the Fr\"ohlich coupling of \cref{sec:6.2}. No finite real-space sum reproduces a divergence, so the long-range part is split off analytically. With the book's normalization,
\begin{equation}
  g^{\mathrm{L}}_{mn\nu}(\vect{k},\vect{Q}) = \ii\,\frac{e^2}{\Omega}\sum_\kappa\frac{1}{\sqrt{2M_\kappa\omega_{\vect{Q}\nu}}}
  \sum_{\vect{G}\ne-\vect{Q}}\frac{\vect{e}_{\kappa\nu}(\vect{Q})\cdot Z^*_\kappa\cdot(\vect{Q}+\vect{G})}{(\vect{Q}+\vect{G})\cdot\varepsilon_\infty\cdot(\vect{Q}+\vect{G})}\,
  \langle\psi_{m,\vect{k}+\vect{Q}}|\ee^{\ii(\vect{Q}+\vect{G})\cdot(\vect{r}-\vect{\tau}_\kappa)}|\psi_{n\vect{k}}\rangle ,
  \label{eq:E.5-gL}
\end{equation}
with $\vect{\tau}_\kappa$ the position of atom $\kappa$ in the cell. It is the potential of the bound charge $-\nabla\cdot\vect{P}$ of the displacement wave, with $\vect{P} = e\,\vect{u}_\kappa\cdot Z^*_\kappa/\Omega$ per cell, screened by $\varepsilon_\infty$ and acting on an electron of charge $-e$, for the displacement per unit $B_{\vect{Q}\nu}$ of \cref{eq:A.6-u}; it is Ref.~\cite[Eq.~(180)]{giustino2017electron} in Heaviside--Lorentz units and the normalization of \cref{app:A.6}~\cite{verdi2015frohlich}. The procedure subtracts $g^{\mathrm{L}}$ from the coarse-grid vertices, interpolates the smooth short-range remainder $g^{\mathrm{S}}$, and adds $g^{\mathrm{L}}$ back on the fine grid. The overlap in \cref{eq:E.5-gL} is needed at the fine-grid points; for small $\vect{Q}+\vect{G}$ it is $[U(\vect{k}+\vect{Q})U^\dagger(\vect{k})]_{mn}$ in the Wannier gauge~\cite[Sect.~VI.A.3]{giustino2017electron}, and the sum over $\vect{G}$ is cut off by a Gaussian factor in practice. The inputs $Z^*_\kappa$ and $\varepsilon_\infty$ come from \cref{app:E.4}, with the caveat about the functional's $\varepsilon_\infty$ stated there.

Besides its dipole, a displaced atom carries a dynamical quadrupole, which is nonzero even in silicon, where every Born charge vanishes, and its potential gives vertices that are finite as $\vect{Q}\to 0$ but depend on the direction of $\vect{Q}$. Such a nonanalytic angular dependence is long-ranged in real space and is not removed by a finer coarse grid: densifying the grid confines the spurious oscillations of the interpolated potential to a smaller region around $\Gamma$ without reducing their amplitude~\cite{brunin2020electron}. In silicon, omitting the quadrupole term changes the computed electron mobility by about \SI{10}{\percent}, and in GaAs by about \SI{30}{\percent}~\cite{brunin2020electron}. Codes that support it subtract the quadrupolar term along with the dipole; where the code does not, the error must be estimated from the small-$Q$ behavior of directly computed vertices.

One rule ties all of this together, the rule of \cref{app:A.5}: every quantity that enters one amplitude must be interpolated with the same $U(\vect{k})$. The vertex $g$, the velocity, the density vertex and the energies of a phonon-assisted amplitude are products of matrices in band space, and a gauge mismatch between two of them leaves the eigenvalues intact while scrambling the amplitude. The safest practice is to compute all coarse-grid quantities from one set of wavefunctions, rotate them with one set of $U$ matrices, and interpolate them with one $H^W(\vect{R})$. Quasiparticle corrections (\cref{app:E.3}) are interpolated on the same Wannier basis, and since they change the energy denominators but not the wavefunctions, velocity matrix elements must be rescaled by the ratio of quasiparticle to Kohn--Sham transition energies, \cref{eq:E.3-vrescale}~\cite{noffsinger2012phonon}.

\subsection{Density vertices at large momentum}

The density vertex $\rho_{mn}(\vect{k},\vect{q})$ is interpolated by the same logic, but with a restriction that matters for dark matter. In the Wannier gauge it is the matrix $\langle w_{m'\vect{0}}|\ee^{\ii\vect{q}\cdot\vect{r}}|w_{n'\vect{R}}\rangle$, Fourier transformed in $\vect{R}$. For small $\vect{q}$ one may expand $\ee^{\ii\vect{q}\cdot\vect{r}} \approx 1 + \ii\vect{q}\cdot\vect{r}$, and the vertex follows from the overlaps and the position matrix elements already computed for the Wannier functions; this is the small-$\vect{q}$ form used for $g^{\mathrm{L}}$ above. The expansion is controlled by $|\vect{q}|$ times the spatial extent of a Wannier function, which is of the order of a bond length, \SI{2.35}{\angstrom} in silicon. At the momentum transfers of electron recoil, a keV and above, $|\vect{q}|\times\SI{2.35}{\angstrom} \approx 1.2$ at \SI{1}{keV} and grows linearly beyond, so the expansion fails. The momentum transfer is a parameter of the vertex, not a lattice wavevector, and for each $\vect{q}$ of interest the real-space matrix must be built anew.

The practical route therefore computes $\rho_{mn}(\vect{k},\vect{q})$ directly on the coarse $\vect{k}$ grid for each target $\vect{q}$, from the plane-wave coefficients of the states at $\vect{k}$ and at $\vect{k}+\vect{q}$ (\cref{app:E.7}), rotates it with the same $U$ matrices into the Wannier gauge, transforms it to real space in $\vect{R}$, and interpolates in $\vect{k}$ alone. The cost is one non-self-consistent calculation at the shifted grid $\vect{k}+\vect{q}$ per target $\vect{q}$, and the number of targets is set by the resolution needed in $|\vect{q}|$ and in direction. For a cubic crystal the directions can be restricted to the irreducible wedge.

\subsection{Counting matrix elements and memory}

\begin{workedexample}{Matrix elements of a converged silicon rate}\label{we:E-counts}
The setup of Ref.~\cite{noffsinger2012phonon} has $N_W = 26$ Wannier functions, coarse $6^3$ grids for $\vect{k}$ and $\vect{Q}$, and fine $40^3$ grids for both, with two atoms and six branches. We count what a phonon-assisted rate needs and where it must live.

\emph{Fine-grid vertices.} There are $N_k N_Q = 64{,}000^2 = 4.1\times10^9$ pairs $(\vect{k},\vect{Q})$, and at each pair $N_W^2\times6 = 4{,}056$ complex vertices. At \SI{16}{bytes} per complex number this is $1.7\times10^{13}$ numbers, or \SI{270}{TB}. Nobody stores it. Codes loop over $\vect{Q}$ (or over $\vect{k}$), interpolate the vertices for all partners at once, accumulate the rate, and discard them. Even one $\vect{Q}$ with all $64{,}000$ partners and all band pairs is \SI{4}{GB}, which is why the band indices are restricted at once to the bands that can be initial, final or intermediate states at the energies of interest.

\emph{Real-space vertices.} The matrix $g^W(\vect{R}_e,\vect{R}_p)$ of \cref{eq:E.5-ginterp} has one entry for each pair of lattice vectors in the Wigner--Seitz supercells of the two coarse grids, about $6^3 = 216$ vectors each, times $N_W^2$ band pairs and six displacement directions: $216^2\times676\times6 \approx 1.9\times10^8$ complex numbers, \SI{3}{GB}. This is the object that is kept, and it must fit in the memory of one node, or be distributed over several.

\emph{Interpolation work.} \Cref{eq:E.5-ginterp} is done in two stages. For each $\vect{Q}$ the sum over $\vect{R}_p$ costs $216\times216\times676\times6 \approx 1.9\times10^8$ complex multiply--adds; for each $\vect{k}$ at that $\vect{Q}$ the sum over $\vect{R}_e$ costs $216\times676\times6 \approx 8.8\times10^5$, plus the $U$ rotations, which are smaller. Over all pairs the second stage dominates: $4.1\times10^9\times8.8\times10^5 \approx 3.6\times10^{15}$ multiply--adds, about $3\times10^{16}$ floating-point operations. Restricting $\vect{Q}$ to the irreducible wedge divides this by the 48 operations of the cubic group, to $6\times10^{14}$, and a pair of energy windows does far more.

\emph{Energy windows.} Near the absorption edge only valence states within a few tenths of an eV of the valence maximum and conduction states within the same range of the six conduction minima contribute. For a parabolic band of mass $m^*$, the states within $\delta$ of the edge fill a sphere of radius $k_\delta = (2m^*\delta)^{1/2}$, a fraction $(4\pi/3)k_\delta^3/V_{\mathrm{BZ}}$ of the zone, with $V_{\mathrm{BZ}} = (2\pi)^3/\Omega = \SI{6.19}{\angstrom^{-3}}$. With $\delta = \SI{0.3}{eV}$ and masses of a few tenths of $m_e$ (\cref{sec:3.3}), say $0.5\,m_e$ for the holes and, per conduction valley, the density-of-states mass $(m^*_\parallel m^{*2}_\perp)^{1/3} = (0.92\times0.19^2)^{1/3}\,m_e = 0.32\,m_e$, the fractions are \num{5.3e-3} for one valence band and $6\times\num{2.7e-3} = \num{1.6e-2}$ for the six valleys. On the $40^3$ grid this leaves about 340 valence and 1{,}040 conduction points per band, so $3.5\times10^5$ pairs instead of $4.1\times10^9$, four orders of magnitude fewer. The window, not the interpolation, is what makes fine grids affordable near an edge.
\end{workedexample}

\subsection{Phonon linewidths and the Eliashberg function}

The phonon linewidth $\gamma^{\mathrm{ep}}_{\vect{Q}\nu}$ and the Eliashberg function $\alpha^2F(\omega)$ of \cref{sec:6.4} are Fermi-surface averages of $|g|^2$: double sums over $\vect{k}$ with two delta functions that place both the initial and the final electron on the Fermi surface. They are the hardest integrals in this section, because a delta function in three dimensions samples a two-dimensional surface, and a product of two samples the intersection of two surfaces. In practice each delta function is replaced by a Gaussian of width $\sigma$. The width must be larger than the spacing of the fine-grid energies near the Fermi surface, of order $|\nabla_{\vect{k}}\varepsilon|\,\Delta k$ with $\Delta k = |\vect{b}_i|/N_{\mathrm{fine}}$ the fine-grid step and $N_{\mathrm{fine}}$ the number of fine-grid points per direction, or the sum is noise. It must also be smaller than the energy scale on which $|g|^2$ and the density of states vary, or the result is biased. Convergence means decreasing $\sigma$ while refining the grid until the result is independent of both.

For metals with complicated Fermi surfaces the requirement is severe. For lead, the DFPT phonons needed $14^3$ $\vect{k}$ points and a $10^3$ wavevector grid with a Gaussian broadening of \SI{0.025}{Ry}, the electron--phonon vertices were interpolated from $6^3$ coarse grids, and the coupling strength $\lambda_{\mathrm{ep}}$ converged only with 400{,}000 quasi-random $\vect{k}$ points and 27{,}000 random wavevectors~\cite[Fig.~12]{ponce2016epw}. Quasi-random (Sobol) sampling converged faster than uniform grids in that study~\cite{ponce2016epw}, since it avoids the systematic aliasing of a regular grid with the Fermi surface. For an undoped insulator at $T = 0$ the electron--phonon linewidth below the gap vanishes (\cref{sec:6.4}), and the calculation is needed only for metals, superconductors (\cref{app:E.9}) and doped semiconductors.

\subsection{Convergence of the phonon-assisted dielectric function}

The phonon-assisted $\varepsilon_2^{\mathrm{PA}}(\omega)$ of \cref{sec:6.5} is a double sum over $\vect{k}$ and $\vect{Q}$ of the squared two-path amplitude, with one energy delta function~\cite[Eqs.~(1)--(3)]{noffsinger2012phonon}. Its convergence parameters are the two fine grids, the broadening of the delta function, the imaginary part $\varGamma_{\mathrm{w}}$ given to the intermediate states, and the number of intermediate bands. The energy resolution is set by the grid: in silicon, $40^3$ grids for both sums resolve features of \SI{14}{meV}, which is what the onset needs, while the smoother spectrum between the indirect and direct gaps converges at $24^3$ with a resolution of \SI{30}{meV}~\cite{noffsinger2012phonon}. In that range the result is insensitive to the intermediate-state width, set to a constant \SI{100}{meV}~\cite{noffsinger2012phonon}, because the intermediate states are far from resonance; near the direct gap they are not, and the width matters (\cref{sec:13.6}).

The comparison with experiment is where the physics outside the calculation shows. The computed spectra of Ref.~\cite{noffsinger2012phonon} had to be shifted rigidly by \SIrange{0.15}{0.23}{eV} to match the measured onsets. The shift absorbs the remaining error of the quasiparticle gap, of order \SI{0.1}{eV}, and the electron--phonon renormalization of the gap, \SIrange{0.05}{0.1}{eV} between 0 and \SI{400}{K}, which the calculation did not include~\cite{noffsinger2012phonon}. The second effect is the subject of \cref{sec:13.7,app:E.8}. A computed absorption edge without such a shift is accurate only to the accuracy of the gap.

\subsection{A phonon-assisted electron-recoil calculation}

A phonon-assisted electron recoil (\cref{sec:13.2}) replaces the photon of \cref{sec:6.5} by the dark-matter density vertex at finite $\vect{q}$, and the four paths of its amplitude contain $\rho$, $g$ and the energies at three crystal momenta. The calculation chains the steps above, and each link has its own check. It starts from the ground state and DFPT on a coarse wavevector grid (\cref{app:E.4}), with the frequencies checked against neutron data and the Born charges against zero. Bands on a coarse $\vect{k}$ grid fine enough to place the conduction minimum fix the Wannier manifold, whose frozen window covers the final and intermediate states and whose bands are tested off the grid. The coarse-grid vertices $g$ are rotated into the Wannier gauge with their dipole and quadrupole parts removed, and tested against direct DFPT off the grid. For each target $\vect{q}$, $\rho_{mn}(\vect{k},\vect{q})$ is computed from the plane-wave coefficients and rotated with the same $U(\vect{k})$, and it too is tested against direct values off the grid. For each $\vect{q}$ the window is the range $\omega < E_{\mathrm{min}}(\vect{q})$ with no zero-phonon final state (\cref{sec:13.1}), cut to the fine-grid pockets of states that can contribute; its edges are in \cref{sec:13.9}. In that window the four-path amplitude is summed over the Wannier manifold, with the remainder above it from the projected Sternheimer equation \cref{eq:E.4-dpsiperp}, into the phonon-assisted loss function on a grid in $(|\vect{q}|,\omega)$.

Three controls close the chain. As $\vect{q}\to0$ the density vertex reduces to its optical limit (\cref{sec:3.2}), and the computed loss function must reproduce the phonon-assisted $\varepsilon_2^{\mathrm{PA}}$ of the same code and, after the gap shift above, the measured indirect absorption (\cref{sec:13.5}). With $g$ set to zero the phonon-assisted rate must vanish identically, which catches any zero-phonon contamination of the window. The third is the $f$-sum rule for the density vertices of one occupied state in its nonlocal form, \cref{eq:E.7-fsum}, whose shortfall measures the truncation of the intermediate-state sum once the nonlocal term $\delta_{\mathrm{NL}}$ is included.

The cost of a silicon run follows from the counts above. The DFPT stage is the one of \cref{we:E-dfpt-si}, a few hundred linear-response solutions. The density vertices need one non-self-consistent calculation per target $\vect{q}$; a grid of $|\vect{q}|$ in steps of $\sim\SI{0.1}{keV}$ up to a few keV and a few tens of directions in the irreducible wedge gives several hundred targets, and this stage dominates the plane-wave part of the work. The amplitude stage is a sum over pocket pairs. Near the edge the relevant states lie within $\sim\SI{0.1}{eV}$ of the band extrema, and the window fractions of \cref{we:E-counts}, which scale as $\delta^{3/2}$, drop by a factor $3^{3/2} \approx 5$ from their values at \SI{0.3}{eV}. On a $120^3$ fine grid this leaves about $1.8\times10^3$ valence points per band and $5.4\times10^3$ conduction points. Three valence bands lie within \SI{0.1}{eV} of the maximum, the heavy-hole, light-hole and split-off bands (\SI{44}{meV} below the others, \cref{sec:3.3}), so there are $\sim3\times10^7$ pairs per $\vect{q}$. Momentum conservation fixes the phonon wavevector of each pair, $\vect{Q} = \vect{k}' - \vect{k} - \vect{q}$ up to a reciprocal-lattice vector, so each pair needs its own interpolated vertices. Restricted to one external band and $N_W$ intermediate bands, the four paths need about $4\times216\times26\times6 \approx 1.3\times10^5$ complex multiply--adds of interpolation per pair, plus a smaller amount for $\rho$, which is $\sim\num{1e6}$ floating-point operations. That count assumes the $\vect{R}_p$ stage of \cref{eq:E.5-ginterp} is already done for the pair's $\vect{Q}$; it runs once per distinct $\vect{Q}$ of the pocket-difference set, at $216^2\times26\times6 \approx 7\times10^6$ multiply--adds per external band, and stays below the $\vect{R}_e$ stage for up to $\sim10^5$ distinct $\vect{Q}$. Per $\vect{q}$ the total is $\sim\num{3e13}$ operations, and for several hundred targets $\sim\num{1e16}$. At a sustained \num{e13} operations per second, typical of one modern accelerator, that is $10^3$\,s, a quarter of an hour; at \num{5e10}, typical of a few processor cores, it is two to three days. The estimate scales as $\delta^{3}$ with the energy range above the edge and as $N_{\mathrm{fine}}^6$ with the fine grid, so these two choices, not the plane-wave stages, decide the budget.

\subsection{Pitfalls}

A coarse grid that looks converged for bands and phonons can still be too coarse for vertices, and a fine grid that resolves a \SI{30}{meV} spectrum can be too coarse for a \SI{14}{meV} onset. Each convergence statement belongs to the quantity, the energy range and the resolution for which it was made. \Cref{app:E.6} turns from the vertices to the many-body response that screens them.

\section{Dielectric response from first principles}\label{app:E.6}

The valence electrons of silicon, eight per primitive cell of volume $\Omega = a^3/4 = \SI{40.05}{\angstrom^3}$, have a plasma frequency $\omega_p = (4\pi\alpha\,\nd_e/m_e)^{1/2} = \SI{16.6}{eV}$, and the loss function $L(\vect{q},\omega)$ of \cref{sec:4.4} has its main peak near that energy at small momentum. Where exactly the peak sits, how sharp it is, and how much spectral weight lies near the gap, where the dark-matter rate is largest for light masses, depend on how the response is computed. \Cref{sec:4.6,sec:4.9} built the hierarchy of many-body approximations: the independent-particle bubble, the random-phase approximation with and without local fields, time-dependent density-functional kernels and the Bethe--Salpeter equation. Each becomes a calculation with its own inputs and cost, and the approximation chosen moves both the peak and the weight near the gap that light dark matter sees.

\subsection{Computed quantities}

Every approximation starts from the crystal bubble of \cref{sec:4.6}, the Adler--Wiser form of the independent-particle response~\cite{adler1962quantum,wiser1963dielectric},
\begin{equation}
  R^0_{\vect{G}\vect{G}'}(\vect{q},\omega) = \frac{g_{\mathrm{deg}}}{N\Omega}\sum_{\vect{k}}\sum_{n,m}
  \bigl(f_{n\vect{k}} - f_{m,\vect{k}+\vect{q}}\bigr)\,
  \frac{\rho_{mn}(\vect{k},\vect{q}+\vect{G})^*\,\rho_{mn}(\vect{k},\vect{q}+\vect{G}')}{\omega + \varepsilon_{n\vect{k}} - \varepsilon_{m,\vect{k}+\vect{q}} + \ii\eta_{\mathrm{b}}} ,
  \label{eq:E.6-R0}
\end{equation}
with $\vect{q}$ in the first zone, $N$ the number of $\vect{k}$ points, $g_{\mathrm{deg}} = 2$ for spin, the density vertices $\rho_{mn}$ of \cref{sec:3.2}, and a broadening $\eta_{\mathrm{b}}$ that tends to $0^+$ in the exact expression and is finite on a grid (below). The random-phase dielectric matrix is $\varepsilon_{\vect{G}\vect{G}'} = \delta_{\vect{G}\vect{G}'} - v_{\mathrm{C}}(\vect{q}+\vect{G})\,R^0_{\vect{G}\vect{G}'}$, the proper response being approximated by the bubble (\cref{sec:4.4}). Time-dependent density-functional theory replaces it by the solution of the matrix Dyson equation
\begin{equation}
  R = R^0 + R^0 f_{\mathrm{Hxc}} R , \qquad (f_{\mathrm{Hxc}})_{\vect{G}\vect{G}'} = v_{\mathrm{C}}(\vect{q}+\vect{G})\,\delta_{\vect{G}\vect{G}'} + (f_{\mathrm{xc}})_{\vect{G}\vect{G}'}(\vect{q}),
  \label{eq:E.6-dyson}
\end{equation}
from which $\varepsilon^{-1} = 1 + v_{\mathrm{C}}R$ in matrix form; with $f_{\mathrm{xc}} = 0$ it reduces to the random-phase result. From the inverse matrix come the macroscopic dielectric function $\varepsilon_{\mathrm{mac}} = 1/(\varepsilon^{-1})_{\vect{0}\vect{0}}$ and the loss function. For a momentum transfer outside the first zone, $\vect{q} = \vect{q}_r + \vect{G}_q$ in the notation of \cref{sec:3.1}, the loss function that enters the rate is the diagonal element of the inverse at $\vect{G}_q$,
\begin{equation}
  L(\vect{q},\omega) = \Imag\bigl[-\bigl(\varepsilon^{-1}\bigr)_{\vect{G}_q\vect{G}_q}(\vect{q}_r,\omega)\bigr] ,
  \label{eq:E.6-L}
\end{equation}
so a calculation at dark-matter momentum transfers of several keV needs the matrix at reciprocal vectors far from the origin, not just its head (\cref{sec:10.2}). The Bethe--Salpeter equation leaves the two-point language altogether; we treat it separately below.

\subsection{Place on the hierarchy}

On the truncation label $\mathfrak{T}$ of \cref{sec:8.5}, these calculations differ only in the many-body approximation $\mathcal{X}$: IP for \cref{eq:E.6-R0} without screening, RPA for the diagonal of $\varepsilon_{\vect{G}\vect{G}'}$, RPA with local fields for the full matrix, TDDFT for \cref{eq:E.6-dyson} with a kernel, and BSE. All are at $n_B = 0$, since the nuclei are clamped and $\varepsilon_\infty$ is the static-lattice constant (the lattice part of $\varepsilon_0$ in a polar crystal comes from \cref{app:E.4}), and at zeroth thermal order, $n_T = 0$, with zero-temperature occupations on the static lattice; the temperature dependence of the gap enters only through a shift (\cref{app:E.8}). As \cref{sec:4.6,sec:8.5} stress, $\mathcal{X}$ is not an expansion in a small number. Its error is estimated from the spread between approximations, and for that estimate to mean anything, every approximation must be converged with the same numerical parameters.

\subsection{Convergence parameters}

The wavefunction cutoff, the local-field cutoff, the band count, the $\vect{k}$ grid and the gap control \cref{eq:E.6-R0,eq:E.6-dyson}, and they interact.

\emph{The wavefunction cutoff} limits the largest momentum transfer at which $\rho_{mn}$ can be trusted. A density vertex at $\vect{q}+\vect{G}$ is an overlap of two sets of plane-wave coefficients displaced by $\vect{q}+\vect{G}$, and when the displacement approaches the cutoff momentum, $k_{\mathrm{cut}} = (2m_eE_{\mathrm{cut}})^{1/2}$, few coefficients overlap. For a cutoff of \SI{500}{eV}, $k_{\mathrm{cut}} = \SI{22.6}{keV}$. Within that range the vertex is still only as good as the pseudo-wavefunction is near the nucleus, which is the subject of \cref{app:E.2,app:E.7}.

\emph{The local-field cutoff} sets the size $N_G$ of the matrix $\varepsilon_{\vect{G}\vect{G}'}$, through a separate, smaller energy cutoff on $|\vect{q}+\vect{G}|$. Building the matrix costs $\propto N_G^2$ per transition and inverting it $\propto N_G^3$ per $(\vect{q},\omega)$. The test is the change of the quantity of interest, $\varepsilon_{\mathrm{mac}}$ near the gap or $L$ at the target $\vect{q}$, as $N_G$ grows. At large $|\vect{q}|$ the matrix must contain reciprocal vectors around $\vect{G}_q$, not only around the origin.

\emph{The number of empty bands} truncates the sum over $m$ in \cref{eq:E.6-R0}. Its test is a sum rule: $\int_0^\infty\dd\omega\,\omega\,\varepsilon_2 = (\pi/2)\omega_p^2$ and the same holds for $\omega L$ (\cref{sec:4.4}), and a truncated band sum misses the weight at high energy. A free-electron count gives the scale. Above the bottom of the valence band, $N_{\mathrm{b}}$ bands per spin fill a sphere of radius $k$ with $\Omega k^3/(6\pi^2) = N_{\mathrm{b}}$, so the highest included band lies at $k^2/2m_e$ above the bottom. In silicon 70 bands reach $k = \SI{4.7}{\angstrom^{-1}}$, which is \SI{84}{eV} above the bottom of the valence band. The weight above that is small but not zero, and the loss function at large $\omega$ falls short of the measured one for exactly this reason~\cite{knapen2021dark}.

\emph{The $\vect{k}$ grid} is usually the bottleneck in a semiconductor~\cite{albrecht1998ab}. The bands disperse strongly, and \cref{eq:E.6-R0} is a sum of sharp poles that must merge into a continuum. The broadening $\eta_{\mathrm{b}}$ must exceed the spacing of the transition energies on the grid, of order $|\nabla_{\vect{k}}(\varepsilon_m - \varepsilon_n)|\,\Delta k$; a spectrum computed with too small an $\eta_{\mathrm{b}}$ for its grid shows spurious peaks. For finite $\vect{q}$, the first-zone part $\vect{q}_r$ must be a difference of grid points, so the $\vect{k}$ grid also fixes the momentum resolution. A coarse grid with a moderate $\eta_{\mathrm{b}}$ suffices where the spectrum is smooth on the scale of $\eta_{\mathrm{b}}$; the region near the gap is not, and it needs much finer sampling, often on a shifted or locally refined grid.

\emph{The gap} enters every denominator. Kohn--Sham gaps are too small (\cref{sec:3.4}), and the usual remedies are a rigid scissor shift $\Delta_{\mathrm{scis}}$ of the empty bands, quasiparticle energies from \cref{app:E.3}, or a functional designed for gaps. A scissor keeps the states and raises the transition energies, and its effect on the $f$-sum check follows from that. The density vertices are unchanged (\cref{app:E.3}), so in the optical limit the spectrum moves rigidly, $\varepsilon_2(\omega) = \varepsilon_2^{\mathrm{KS}}(\omega - \Delta_{\mathrm{scis}})$, and its first moment becomes $\int\dd\omega\,\omega\,\varepsilon_2^{\mathrm{KS}}(\omega - \Delta_{\mathrm{scis}}) = (\pi/2)\omega_p^2 + \Delta_{\mathrm{scis}}\int\dd\omega\,\varepsilon_2^{\mathrm{KS}}$, larger by $\Delta_{\mathrm{scis}}\int\dd\omega\,\varepsilon_2^{\mathrm{KS}}$. The excess is the $f$-sum rule of a nonlocal operator, the scissor, and a band-convergence test on a scissored spectrum subtracts it. Codes that build the optical limit from velocity matrix elements must rescale them as in \cref{eq:E.3-vrescale}; otherwise the spectrum acquires a spurious factor $[(\omega - \Delta_{\mathrm{scis}})/\omega]^2$ and is no longer the rigid shift that unchanged states imply. A density vertex computed directly from plane-wave coefficients at finite $\vect{q}$ needs no such care, since it depends on the states and not on their energies.

\emph{Semicore and core states} matter at large $\omega$ and large $\vect{q}$. The silicon $2p$ and germanium $3d$ shells are frozen into the core by most pseudopotentials. Their contribution to $L$ starts at their binding energies, and their contribution to the $f$-sum adds to the valence one. A loss function for $\omega$ above a few tens of eV, or at $|\vect{q}|$ where core orbitals overlap the probe, needs them in the valence or an all-electron treatment (\cref{sec:10.6,app:E.7}).

\subsection{Local fields and exchange--correlation kernels}

Local fields are the difference between the inverse of the averaged matrix and the average of the inverse (\cref{sec:4.6}). Computationally they cost one matrix inversion per $(\vect{q},\omega)$, and they change the spectrum where the response is inhomogeneous on the scale of the cell. In the optical spectrum of silicon they lower and shift the two main peaks to higher energy, which moves the random-phase result further from experiment~\cite{albrecht1998ab}. In the dark-matter rate of silicon and germanium, by contrast, the difference between the loss function with and without local fields is small~\cite{knapen2021dark}, because the rate integrates over $\vect{q}$ and $\omega$ and is dominated by energies and momenta where the response is smooth.

The adiabatic local-density kernel adds little to the random-phase result for the absorption of semiconductors: its spectra are extremely close to those of the random-phase approximation with local fields~\cite{reining2002excitonic}. For energy-loss spectra of solids the same kernel does well, but largely because the random-phase approximation with local fields already agrees with experiment there~\cite{reining2002excitonic}, which is good news for the loss function of a dark-matter rate. The reason is that the kernel is short-ranged, while the electron--hole attraction that produces excitons is long-ranged. A kernel of the form $f_{\mathrm{xc}}(\vect{r},\vect{r}') = -\alpha_{\mathrm{xc}}/(4\pi|\vect{r} - \vect{r}'|)$, which behaves as $-\alpha_{\mathrm{xc}}/|\vect{q}+\vect{G}|^2$ in reciprocal space, fits the measured optical spectrum of silicon almost as well as the Bethe--Salpeter equation, with quasiparticle energies in the bubble and the empirical value $\alpha_{\mathrm{xc}} = 0.2$~\cite{reining2002excitonic}, a fit for silicon rather than a universal constant. (Ref.~\cite{reining2002excitonic} writes $\alpha$; we add the subscript to keep the fine-structure constant unambiguous.) The cost of any two-point kernel is that of the random-phase calculation: the matrices have the same size, and only the Dyson equation changes. This is the attraction of TDDFT for dark-matter loss functions on fine $(\vect{q},\omega)$ grids, and the reason its kernels are worth testing against the Bethe--Salpeter equation on a few points.

\subsection{Optical conductivity and large momentum}

Bosonic absorption (\cref{sec:10.7}) and the crystal Migdal response (\cref{sec:12.3}) take opposite limits of one calculation. For absorption of a bosonic dark-matter particle (\cref{sec:10.7}) the relevant quantity is the optical conductivity $\sigma_1 = \omega\varepsilon_2$ at $\vect{q}\to0$, where the transverse and longitudinal dielectric functions coincide (\cref{sec:4.4}). Computing it needs the head of \cref{eq:E.6-R0} in the optical limit, with the velocity matrix elements and their nonlocal commutator (\cref{app:E.2,app:E.7}), a fine $\vect{k}$ grid near the absorption edge, and the approximation that the energy range demands: near the gap the Bethe--Salpeter equation, well above it the random-phase approximation with local fields. Measured optical data are available for most candidate targets, and where they are, they fix $\sigma_1$ better than any calculation. The calculation is then needed where the data stop: at low temperature, below the measured range, or for anisotropic crystals where only some tensor components were measured.

The opposite limit is the loss function at momenta far beyond the zone, up to tens of keV, and energies up to about \SI{100}{eV}, which the crystal Migdal response of \cref{sec:12.3} convolves with the nuclear motion. Every convergence parameter above is pushed at once. The cutoff must hold momentum components well above the largest $|\vect{q}|$, the band sum must reach \SI{100}{eV} above the gap and beyond, and the semicore shells must be in the valence. In addition, the local-field matrix must be centered on $\vect{G}_q$, which at $|\vect{q}| = \SI{30}{keV}$ lies more than seven shortest reciprocal vectors from the origin. The pseudo-wavefunctions are no longer adequate there, and the all-electron reconstruction of \cref{app:E.2,app:E.7} becomes part of the response calculation. The cost grows with all of these at once, and in practice such a calculation is done on a coarse grid of $|\vect{q}|$ and a few directions, with the convergence tested on the $f$-sum rule at each $\vect{q}$.

\subsection{The Bethe--Salpeter equation}

The Bethe--Salpeter equation of \cref{sec:4.9} is an eigenvalue problem in the space of transitions $(v,c,\vect{k})$. Its Hamiltonian, in the resonant (Tamm--Dancoff) form, is
\begin{equation}
  H^{\mathrm{exc}}_{vc\vect{k},v'c'\vect{k}'} = \bigl(\varepsilon_{c\vect{k}} - \varepsilon_{v\vect{k}}\bigr)\,\delta_{vv'}\delta_{cc'}\delta_{\vect{k}\vect{k}'}
  + g_{\mathrm{deg}}\,K^{\mathrm{x}}_{vc\vect{k},v'c'\vect{k}'} - K^{\mathrm{d}}_{vc\vect{k},v'c'\vect{k}'} ,
  \label{eq:E.6-Hexc}
\end{equation}
with the energies from quasiparticle bands (\cref{app:E.3}), an exchange kernel $K^{\mathrm{x}}$ built from the Coulomb kernel without its $\vect{G} = 0$ component, and a direct kernel $K^{\mathrm{d}}$ built from the statically screened interaction $W^{\mathrm{scr}}$ at the momentum $\vect{k} - \vect{k}'$ of the scattered pair (\cref{sec:4.9}); the factor $g_{\mathrm{deg}}$ holds for singlets. Its eigenvalues are the exciton energies $E_\lambda$ and its eigenvectors the envelopes $A^\lambda_{cv\vect{k}}$, and the macroscopic dielectric function is a sum over $\lambda$ of oscillator strengths $|\sum_{cv\vect{k}}A^\lambda_{cv\vect{k}}\rho_{cv}(\vect{k},\vect{q})^*|^2$ over $\omega - E_\lambda$. At finite $\vect{q}$ the transitions are $(v\vect{k}) \to (c,\vect{k}+\vect{q})$ and the same structure gives the excitonic loss function (\cref{sec:10.5}).

The cost has two parts. Building the kernel requires the screened interaction $\varepsilon^{-1}_{\vect{G}\vect{G}'}(\vect{k}-\vect{k}')$, which is a random-phase calculation of its own, and the matrix elements of $K^{\mathrm{d}}$ between all pairs of transitions, $\propto (N_vN_cN_k)^2$ entries, each a sum over $N_G^2$ reciprocal vectors when the full inverse is used. Solving it costs $\propto (N_vN_cN_k)^3$ for a full diagonalization, or $\propto (N_vN_cN_k)^2$ per step of an iterative (Lanczos--Haydock) method that delivers the spectrum without the eigenvectors. The transition-basis dimension $N_vN_cN_k$ is the number to watch. For the optical spectrum of silicon, Ref.~\cite{albrecht1998ab} used norm-conserving pseudopotentials at \SI{15}{Ry}, quasiparticle corrections interpolated onto the fine grid, and up to 2{,}048 $\vect{k}$ points in the full zone. On an 864-point grid four valence and two conduction bands proved sufficient below \SI{4.5}{eV}, and neglecting the off-diagonal elements of $\varepsilon^{-1}$ in $K^{\mathrm{d}}$ changed nothing visible, so the production run used six bands and a diagonal $\varepsilon^{-1}$~\cite{albrecht1998ab}. The dimension was then $2{,}048\times4\times2 \approx 1.6\times10^4$. A dense diagonalization of that size takes $\sim 10^{13}$ operations, trivial today, and the authors identified the $\vect{k}$ grid, not the cutoff, as the bottleneck~\cite{albrecht1998ab}.

Two lessons of that calculation carry over: symmetry rotations must keep the phases of the states consistent, and a finite grid exaggerates individual structures. Symmetry can be used to reduce the zone, but the spectrum turned out to be extremely sensitive to inconsistent phases when wavefunctions are rotated, especially for degenerate bands, and only a subgroup of the operations was used safely~\cite{albrecht1998ab}. This is the gauge rule of \cref{app:A.5} once more. And a finite $\vect{k}$ grid overestimates individual structures; the shoulder at \SI{3.8}{eV} in the silicon spectrum came out slightly too strong for that reason~\cite{albrecht1998ab}.

Bound excitons near the gap are harder still, because their envelopes are wide in real space and narrow in $\vect{k}$. The Wannier--Mott relation of \cref{sec:4.9}, $E_x^b = \mu_x\alpha^2/(2\varepsilon_\infty^2)$, with the measured binding energy of \SI{14.7}{meV} and $\varepsilon_\infty = 11.7$, gives $\mu_x \approx 0.15\,m_e$ and a radius $a_x = \varepsilon_\infty a_0 m_e/\mu_x \approx \SI{4.2}{nm}$. Resolving an envelope of that size needs a $\vect{k}$ spacing well below $1/a_x \approx \SI{0.024}{\angstrom^{-1}}$, against a zone of width $|\vect{b}_i| = \sqrt3\,2\pi/a = \SI{2.0}{\angstrom^{-1}}$: close to a hundred points per direction near the band edges, and so of order $10^6$ points if the grid were uniform. Calculations of such excitons therefore use grids refined only around the band extrema, or the effective-mass limit, with the Bethe--Salpeter equation reserved for the continuum and the finite-momentum structure (\cref{sec:10.5}).

\subsection{Molecular \texorpdfstring{$GW$}{GW} and Bethe--Salpeter calculations}

Molecular crystals (\cref{sec:16.4}) are often treated molecule by molecule, with the environment added afterwards. The same $GW$ and Bethe--Salpeter equations are then solved in a basis of atom-centered Gaussian orbitals instead of plane waves. The $GW$ quasiparticle energies of molecules improve the Kohn--Sham gaps to within a few tenths of an eV at a cost that scales as the fourth power of the system size, and the Bethe--Salpeter singlet excitation energies that follow have typical errors of \SIrange{0.1}{0.3}{eV}, with a mean error of \SI{0.2}{eV} on a standard benchmark set of about 200 excitations~\cite{blase2020bethe}. An iterative search for the lowest eigenstates costs the same as in TDDFT, $O(N_{\mathrm{orb}}^4)$ operations with $N_{\mathrm{orb}}$ the number of orbitals, when resolution-of-the-identity techniques are used to build the kernel~\cite{blase2020bethe}. The convergence parameter specific to this setting is the Gaussian basis: diffuse functions are needed for the empty states and for the screening, and results are tested along a family of basis sets of increasing size. The crystal environment enters through its dielectric screening, as a reaction field added to the bare Coulomb interaction~\cite{blase2020bethe}, or through a periodic calculation of the full crystal. \Cref{sec:16.4} compares the periodic result with measurement for anthracene: \SI{3.3}{eV} against the measured \SI{3.1}{eV} for the lowest singlet of the crystal.

\subsection{Silicon}

\begin{workedexample}{The loss function of silicon from Kohn--Sham states}\label{we:E-elf-si}
A published dark-matter calculation of the silicon loss function used a plane-wave basis with a \SI{500}{eV} cutoff, a $\Gamma$-centered $8\times8\times8$ grid of $\vect{k}$ points, 70 bands per spin, a functional designed for band gaps with a scissor shift to the measured gap, and the random-phase dielectric matrix with local fields at all $\vect{q}$ of the grid~\cite{knapen2021dark}. For the rate, the matrix was reduced to an angular average of its diagonal elements over all $\vect{q}_r + \vect{G}$ of equal length~\cite[Eq.~(24)]{knapen2021dark}. We check each setting against the estimates of this section.

The cutoff corresponds to $k_{\mathrm{cut}} = (2m_e\times\SI{500}{eV})^{1/2} = \SI{22.6}{keV}$, the range quoted by the authors~\cite{knapen2021dark}. Kinematics allows momentum transfers up to $2\mchi v_{\mathrm{max}}$, \SI{5.3}{keV} per MeV of dark-matter mass (\cref{sec:1.3}), but the valence density vertices fall off beyond a few times $\alpha m_e = \SI{3.7}{keV}$ (\cref{sec:3.5,sec:10.1}), and the cutoff covers that range, provided the pseudo-wavefunctions are accurate there (\cref{app:E.7}). The 70 bands reach \SI{84}{eV} above the bottom of the valence band by the free-electron count above, about \SI{70}{eV} above the valence maximum since the occupied bands are about \SI{12}{eV} wide (\cref{sec:3.3}). The authors traced the shortfall of the computed loss function at large energy to this truncation~\cite{knapen2021dark}. An $8^3$ grid has spacing $|\vect{b}_i|/8 = \SI{0.25}{\angstrom^{-1}}$, or \SI{0.49}{keV}, which is the resolution in $\vect{q}_r$; it is adequate for a loss function dominated by energies of several eV, and coarse for the structure within a few tenths of an eV of the gap, where excitonic and phonon-assisted effects live (\cref{sec:13.1}).

The $f$-sum rule is the cleanest single check. With $\omega_p = \SI{16.6}{eV}$ from the valence density, the integral $\int_0^\infty\dd\omega\,\omega L(\vect{q},\omega) = (\pi/2)\omega_p^2$ at small $\vect{q}$ is \SI{433}{eV^2}. A calculation that recovers, say, 0.9 of it is missing a tenth of the weight, all at high energy, and that missing weight is a bound on what the band truncation can do to a rate at large $\omega$.
\end{workedexample}

The optical spectrum gives the complementary test near the gap. On top of quasiparticle energies, the random-phase spectrum of silicon misplaces the weight of the two main peaks, local fields alone make the agreement worse, and only the electron--hole attraction of the Bethe--Salpeter equation restores the measured line shape and absolute intensity~\cite{albrecht1998ab}. The same hierarchy matters for a dark-matter rate in proportion to the weight it gives to energies within an eV or two of the gap, which is why the many-body spread of \cref{sec:4.6,sec:4.9} is quoted as an uncertainty for light masses and becomes negligible for heavier ones.

\subsection{Pitfalls}

Three settings are easy to get wrong and hard to spot afterwards: the head element, the reciprocal-lattice index of the momentum transfer, and the broadening. The head element $\varepsilon_{\vect{0}\vect{0}}$ used in place of $1/(\varepsilon^{-1})_{\vect{0}\vect{0}}$ drops the local fields, and the element at $\vect{G} = 0$ used in place of $\vect{G}_q$ for a momentum transfer beyond the first zone computes the response at the wrong momentum. A broadening too large for the grid smears the onset, which is the part of the spectrum that light dark matter sees. The response computed here screens the density vertex, the electron--phonon vertex of \cref{app:E.5} and, through $W^{\mathrm{scr}}$, the quasiparticle energies of \cref{app:E.3}, so every one of these errors propagates into the rates of the later chapters.

\section{Crystal form factors}\label{app:E.7}

An electron-recoil rate of silicon at independent-particle level is a sum of $|\rho_{cv}(\vect{k},\vect{q})|^2$ over valence states $v\vect{k}$, conduction states $c\vect{k}'$ and momentum transfers $\vect{q}$, each term weighted by an energy delta function. With four valence bands, 52 conduction bands, a $10\times10\times10$ grid of crystal momenta and every reciprocal-lattice vector up to $q = \SI{30}{keV}$, the sum has about $5\times10^{11}$ terms, as we count below. Each term is an overlap of two sets of plane-wave coefficients. The work is in restoring the core region that the pseudopotential removed, adding the core shells, and knowing when the sum has converged; the bands come from the calculation of \cref{app:E.1}, and the screening that turns the bare vertex into a measured loss function from \cref{app:E.6}.

\subsection{The density vertex}

The object is the density vertex of \cref{eq:A.5-rho},
\begin{equation*}
  \rho_{mn}(\vect{k},\vect{q}) = \langle\psi_{m,\vect{k}+\vect{q}}|\ee^{\ii\vect{q}\cdot\vect{r}}|\psi_{n\vect{k}}\rangle ,
\end{equation*}
with $\vect{q}$ never reduced to the first zone (\cref{app:A.3,app:A.4}). From it the independent-particle structure factor follows as a sum over pairs (\cref{sec:10.1}); the crystal form factor $f_{\mathrm{crys}}$ of the literature is the same object in another normalization, \cref{eq:A.8-fcrys}; and the vertices of the spin and current couplings of \cref{sec:10.4} are built the same way, with $\ee^{\ii\vect{q}\cdot\vect{r}}$ multiplied by a Pauli matrix or by the velocity.

\subsection{Place on the hierarchy}

On the truncation label $\mathfrak{T}$ of \cref{sec:8.5}, a form-factor calculation sits at the base point $\mathfrak{T}_0$. The nuclei are clamped at their equilibrium positions, $n_B = 0$; there is no electron--phonon insertion, $n_g = 0$, and no anharmonic or thermal correction, $n_A = 0$ and zeroth thermal order, $n_T = 0$; the many-body approximation is that of independent particles, $\mathcal{X} = \mathrm{IP}$. The plane wave is kept whole, so the result is exact to all orders in the multipole parameter $\xi_D$, and the dipole form $\vect{q}\cdot\vect{v}_{mn}/(\varepsilon_{m\vect{k}} - \varepsilon_{n\vect{k}})$ of \cref{eq:A.5-optical} is only a check at small $q$. The Kohn--Sham approximation and the pseudopotential sit inside $\mathcal{X} = \mathrm{IP}$ and are not part of the label. The states are Kohn--Sham states rather than quasiparticles (\cref{sec:3.4}), with a scissor shift $\Delta_{\mathrm{scis}}$ on the conduction energies or with $GW$ energies (\cref{app:E.3}); and they are pseudo-wavefunctions unless the region inside the core radius is reconstructed. What the base point leaves out is screening and local fields ($\mathcal{X} = \mathrm{RPA}$ and beyond, \cref{app:E.6}), excitons ($\mathcal{X} = \mathrm{BSE}$), and the lattice ($n_g \ge 1$, \cref{app:E.5,app:E.8}). Screening divides the squared potential by $|\varepsilon(\vect{q},\omega)|^2$, which is $\varepsilon_\infty^2$ as $q \to 0$ and approaches one only at large $q$ (\cref{sec:4.4}), so a form factor at $\mathfrak{T}_0$ is an input to that correction, not a substitute for it.

\subsection{Plane-wave evaluation}

A plane-wave code returns, for each band $n$ and each $\vect{k}$ of its grid, the coefficients $c_{n\vect{k}}(\vect{G})$ of
\begin{equation}
  u_{n\vect{k}}(\vect{r}) = \frac{1}{\sqrt\Omega}\sum_{|\vect{k}+\vect{G}|^2\le 2m_eE_{\mathrm{cut}}} c_{n\vect{k}}(\vect{G})\,\ee^{\ii\vect{G}\cdot\vect{r}}, \qquad \sum_{\vect{G}}|c_{n\vect{k}}(\vect{G})|^2 = 1 ,
  \label{eq:E.7-pw}
\end{equation}
which is the normalization of \cref{eq:A.5-bloch} and the code convention of \cref{app:A.3} (coefficients of $\ee^{+\ii\vect{G}\cdot\vect{r}}$, the opposite of the book's lattice Fourier components). The cutoff $E_{\mathrm{cut}}$ is the one of \cref{sec:3.5}. With $\vect{k}+\vect{q}$ folded back into the zone, $\vect{k}+\vect{q} = \vect{k}' + \vect{G}_{\mathrm{U}}$ with $\vect{k}'$ a point of the grid and $\vect{G}_{\mathrm{U}}$ a reciprocal-lattice vector, and by \cref{eq:A.5-periodic} the vertex is $\int_\Omega u^*_{m\vect{k}'}\ee^{\ii\vect{G}_{\mathrm{U}}\cdot\vect{r}}u_{n\vect{k}}$, and inserting \cref{eq:E.7-pw} the cell integral of $\ee^{\ii(\vect{G}_{\mathrm{U}} + \vect{G} - \vect{G}')\cdot\vect{r}}$ is $\Omega\,\delta_{\vect{G}',\vect{G}+\vect{G}_{\mathrm{U}}}$:
\begin{equation}
  \rho_{mn}(\vect{k},\vect{q}) = \sum_{\vect{G}} c^*_{m\vect{k}'}(\vect{G} + \vect{G}_{\mathrm{U}})\,c_{n\vect{k}}(\vect{G}), \qquad \vect{q} = \vect{k}' - \vect{k} + \vect{G}_{\mathrm{U}} .
  \label{eq:E.7-rhopw}
\end{equation}
The sum runs over the vectors present in both basis sets. Coefficients are matched by their Miller indices, the integer coordinates of $\vect{G}$ in the basis $\vect{b}_i$, since the two states are stored on different lists of $\vect{G}$. A term with $\vect{G}+\vect{G}_{\mathrm{U}}$ outside the basis of $\vect{k}'$ is zero in the representation, so \cref{eq:E.7-rhopw} is exact for the functions the code represents; its error is the basis truncation itself, which we discuss below. This is the expression of Ref.~\cite[Eq.~(3.16)]{essig2016direct}, in the notation of \cref{app:A.5}.

Each value of \cref{eq:E.7-rhopw} costs one dot product of length $N_{\mathrm{PW}}$, the number of plane waves per state. All the $\vect{G}_{\mathrm{U}}$ for one pair $(m\vect{k}', n\vect{k})$ come at once from a fast Fourier transform: the product $u^*_{m\vect{k}'}(\vect{r})\,u_{n\vect{k}}(\vect{r})$ is formed on the real-space grid and transformed, and the component at $\vect{G}_{\mathrm{U}}$ is $\rho_{mn}$ for $\vect{q} = \vect{k}' - \vect{k} + \vect{G}_{\mathrm{U}}$. The cost per pair drops from $N_{\mathrm{PW}}N_{\vect{G}}$ to of order $N_{\mathrm{FFT}}\log N_{\mathrm{FFT}}$, where $N_{\vect{G}}$ is the number of $\vect{G}_{\mathrm{U}}$ needed and $N_{\mathrm{FFT}}$ the number of grid points~\cite[below Eq.~(3.15)]{griffin2021extended}. The real-space grid has to resolve the product, whose components reach twice the wavefunction cutoff. For a cubic cell it covers a cube of side $4\sqrt{2m_eE_{\mathrm{cut}}}$ in reciprocal space, which holds $8\Omega(2m_eE_{\mathrm{cut}})^{3/2}/\pi^3$ points, while the sphere of \cref{eq:E.7-pw} holds $N_{\mathrm{PW}} = \Omega(2m_eE_{\mathrm{cut}})^{3/2}/(6\pi^2)$, so $N_{\mathrm{FFT}} \approx (48/\pi)N_{\mathrm{PW}}$. The fcc grid runs along the primitive vectors, whose lengths multiply to $\sqrt2\,\Omega$, and holds $\sqrt2$ times more, $N_{\mathrm{FFT}} \approx (48\sqrt2/\pi)N_{\mathrm{PW}} \approx 22\,N_{\mathrm{PW}}$: about $36^3$ points at \SI{60}{Ry}, as in \cref{app:E.1}.

\paragraph{Which $\vect{q}$ a grid can reach.} \Cref{eq:E.7-rhopw} needs the state at $\vect{k}' = \vect{k}+\vect{q} - \vect{G}_{\mathrm{U}}$. If both $\vect{k}$ and $\vect{k}'$ are taken from one uniform grid of $n_k = n_1n_2n_3$ points (\cref{app:C.2}), the first-zone part of $\vect{q}$ is a difference of grid vectors, and $\vect{q}$ is resolved in steps of $|\vect{b}_i|/n_i$. A form factor tabulated on a finer or arbitrary set of $\vect{q}$ needs a second, shifted set of states: one non-self-consistent calculation at the points $\vect{k}+\vect{q}$, as density-functional perturbation theory does for phonon wavevectors (\cref{app:E.4}). The states at $\vect{k}'$ must be computed, or unfolded from the irreducible wedge by a symmetry operation, with the phases that the operation carries. In the diamond structure some operations include a fractional translation $\vect{t}_{\mathrm{f}}$, and each coefficient of a rotated Bloch state picks up a phase $\ee^{-\ii\vect{k}_{\mathrm{rot}}\cdot\vect{t}_{\mathrm{f}}}$, with $\vect{k}_{\mathrm{rot}}$ the rotated wavevector $\vect{k}+\vect{G}$. Computing the full grid without symmetry avoids the bookkeeping at a cost of the order of the symmetry factor, 48 for silicon.

\paragraph{Gauge and degeneracy.} Each Bloch state carries an arbitrary phase, and $\rho_{mn}$ picks up the phase difference of its two states, while $|\rho_{mn}|^2$ does not. A rate that sums $|\rho_{mn}|^2$ is therefore gauge independent, provided that the sum over $m$ and $n$ covers complete degenerate sets: a degenerate subspace may be rotated by any unitary matrix, and only the trace over it is fixed. At the valence maximum of silicon the threefold degeneracy at $\mathrm{\Gamma}$ makes this concrete; a sum restricted to one of the three states there is meaningless. Amplitudes that interfere, such as the second-order amplitudes of phonon-assisted processes in which $\rho_{mn}$ multiplies an electron--phonon vertex, need $\rho$ and $g$ in one gauge, computed from the same stored states (\cref{app:A.5,app:E.5}).

\subsection{Velocity matrix elements and nonlocal potentials}

The momentum matrix element follows from \cref{eq:E.7-pw} in the same way,
\begin{equation}
  \vect{p}_{mn}(\vect{k}) = \sum_{\vect{G}} c^*_{m\vect{k}}(\vect{G})\,(\vect{k} + \vect{G})\,c_{n\vect{k}}(\vect{G}) ,
  \label{eq:E.7-ppw}
\end{equation}
but the velocity is not $\vect{p}_{mn}/m_e$ in a pseudopotential calculation. The velocity operator is $\vect{v} = \ii[H,\vect{r}]$. For $H = \vect{p}^2/2m_e + V_{\mathrm{loc}}(\vect{r}) + V_{\mathrm{NL}}$, with $V_{\mathrm{NL}}$ the nonlocal part of the pseudopotential (\cref{app:A.5}), the local potential commutes with $\vect{r}$ and
\begin{equation}
  \vect{v} = \frac{\vect{p}}{m_e} + \ii[V_{\mathrm{NL}},\vect{r}] .
  \label{eq:E.7-v}
\end{equation}
The commutator is evaluated in the $\vect{k}$ representation. Define $V_{\mathrm{NL}}(\vect{k}) \equiv \ee^{-\ii\vect{k}\cdot\vect{r}}V_{\mathrm{NL}}\ee^{\ii\vect{k}\cdot\vect{r}}$, the operator that acts on $u_{n\vect{k}}$. Differentiating, $\nabla_{\vect{k}}V_{\mathrm{NL}}(\vect{k}) = \ee^{-\ii\vect{k}\cdot\vect{r}}\,\ii[V_{\mathrm{NL}},\vect{r}]\,\ee^{\ii\vect{k}\cdot\vect{r}}$, so that
\begin{equation}
  \vect{v}_{mn}(\vect{k}) = \frac{\vect{p}_{mn}(\vect{k})}{m_e} + \langle u_{m\vect{k}}|\nabla_{\vect{k}}V_{\mathrm{NL}}(\vect{k})|u_{n\vect{k}}\rangle .
  \label{eq:E.7-vNL}
\end{equation}
For the separable form \cref{eq:E.2-KB} that plane-wave codes use, with projectors $\tilde p_{\kappa i}$ centered on the atoms and coefficients $D_{\kappa,ij}$, the gradient acts only on the Fourier transforms of the projectors, $\tilde p_{\kappa i}(\vect{k}+\vect{G})$, and costs about as much as applying $V_{\mathrm{NL}}$ once. The correction is not small. In diamond, the absorption spectrum computed with $\vect{p}/m_e$ instead of $\vect{v}$ differs by 15 to 20 percent~\cite[Fig.~1]{ismailbeigi2001coupling}, and Ref.~\cite[Eq.~(20)]{ismailbeigi2001coupling} shows that a long-wavelength field couples to $\vect{v}$, not to $\vect{p}/m_e$.

For the dark-matter vertices the correction has opposite consequences for the density and for the current. The density vertex needs none. The operator $\ee^{\ii\vect{q}\cdot\vect{r}}$ is local, and \cref{eq:E.7-rhopw} evaluated at small $q$ reproduces $\vect{q}\cdot\vect{v}_{mn}/(\varepsilon_{m\vect{k}} - \varepsilon_{n\vect{k}})$ with the full velocity of \cref{eq:E.7-vNL}, since the derivation of \cref{eq:A.5-optical} used only $\ii[H,\vect{r}] = \vect{v}$. The comparison of the two sides at $q \sim 10^{-3}\,|\vect{b}_i|$ is a test that catches a missing nonlocal term in a velocity routine. A coupling to the electron current, on the other hand, is ambiguous with a nonlocal Hamiltonian. The paramagnetic current $j_e(\vect{q})$ of \cref{sec:4.3} is built from $\vect{p}/m_e$; gauge invariance requires an additional nonlocal term, which Ref.~\cite[Eqs.~(18) and~(19)]{ismailbeigi2001coupling} derives as a phase $\exp(\ii\int\vect{A}\cdot\dd\vect{x})$ along the straight line between the two arguments of $V_{\mathrm{NL}}(\vect{r},\vect{r}')$. Its long-wavelength limit is the commutator of \cref{eq:E.7-v}. At finite $q$ the term depends on $\vect{q}$ through the extent of the projectors, of order $q\,r_c$, and codes that compute current responses at finite $q$ either include it or restrict themselves to $q\,r_c \ll 1$.

\paragraph{The scissor shift.} The scissor operator $\Delta_{\mathrm{scis}}\sum_c|\psi_c\rangle\langle\psi_c|$ is nonlocal too. It rescales each interband velocity by $(\varepsilon_{c\vect{k}} - \varepsilon_{v\vect{k}} + \Delta_{\mathrm{scis}})/(\varepsilon_{c\vect{k}} - \varepsilon_{v\vect{k}})$, the rigid-shift case of \cref{eq:E.3-vrescale}, and leaves the density vertex and its optical limit unchanged. An optical spectrum built from $|\vect{v}_{cv}|^2$ with scissored energies but unscaled velocities violates the $f$-sum rule; first-principles absorption calculations therefore carry the nonlocality of the scissor into the oscillator strengths~\cite[Sect.~VII]{zacharias2016one}, where the inverse factor appears because it multiplies their oscillator-strength form rather than the velocity. For silicon in the local-density approximation, with $\Delta_{\mathrm{scis}} = \SI{0.74}{eV}$~\cite[Sect.~VII]{zacharias2016one} and a direct gap at $\mathrm{\Gamma}$ of \SI{2.57}{eV}~\cite[Table~V]{ponce2015temperature}, the factor is 1.29 at $\mathrm{\Gamma}$, and $|\vect{v}_{cv}|^2$ changes by 66 percent.

\subsection{All-electron reconstruction}

A pseudo-wavefunction agrees with the all-electron one outside a sphere of radius $r_c$ around each nucleus and is smooth and nodeless inside it (\cref{sec:3.5}). The two differ only inside the sphere, so their difference has Fourier components at momenta of order $\pi/r_c$ and above, and the pseudo vertex is reliable only for $q\,r_c \ll 1$. The PAW map \cref{eq:E.2-PAW} restores the difference, in one of two ways. The first route expands the all-electron state itself in plane waves up to a second, larger cutoff and then applies \cref{eq:E.7-rhopw} unchanged. Ref.~\cite[Sect.~2.1.1]{griffin2021extended} does this for silicon and germanium with a second cutoff of \SI{2}{keV}, which by $q_{\mathrm{max}} = (2m_eE)^{1/2}$ reaches $q \approx \SI{45}{keV}$. The second route works with the matrix elements, \cref{eq:E.2-rhoPAW}. With $\vect{r} = \vect{\tau}_\kappa + \vect{x}$ in each sphere, the Rayleigh expansion \cref{eq:C.3-planewave} of $\ee^{\ii\vect{q}\cdot\vect{x}}$ reduces each element of the one-center matrix \cref{eq:E.2-onecentre} to Gaunt coefficients times the radial integrals
\begin{equation}
  \int_0^{r_c}\dd x\,x^2\,j_l(qx)\,\bigl[\phi^{\mathrm{AE}}_{\kappa i}(x)\phi^{\mathrm{AE}}_{\kappa j}(x) - \phi^{\mathrm{PS}}_{\kappa i}(x)\phi^{\mathrm{PS}}_{\kappa j}(x)\bigr] ,
  \label{eq:E.7-radial}
\end{equation}
with $\phi(x)$ the radial parts of the partial waves of \cref{eq:E.2-PAW} and $j_l$ the spherical Bessel functions of \cref{app:C.3}. The integrals live on the radial grid of the atomic calculation, which resolves the nodes of the all-electron partial waves, so the correction has no plane-wave cutoff at all; it is limited by the completeness of the partial waves, usually two per angular momentum. They are computed once per species and per $q$, and the per-state work is in the projections $\langle\tilde p_{\kappa j}|\tilde\psi_{n\vect{k}}\rangle$, which the code already stores.

How much the reconstruction matters depends on which $q$ the rate weights. For a heavy mediator, $\FDM = 1$, and the factor $q$ of the measure in \cref{eq:A.8-rateiso} gives large momentum transfers more weight than any lighter mediator does, and the reconstruction raises the silicon rate even above low thresholds (\cref{app:E.2}). For a light mediator, $\FDM^2 \propto q^{-4}$ weights the smallest momentum transfers, and the reconstruction changes the rate little except at high thresholds, where only large $q$ is kinematically allowed. Norm-conserving pseudopotentials can be reconstructed in the same way when the generator writes out its all-electron and pseudo partial waves; without them the pseudo vertex is all there is, and its validity is limited to $q\,r_c \ll 1$.

\subsection{Core and semicore states}

Below the valence bands of silicon lie the atomic shells $2p$, $2s$ and $1s$, with measured binding energies relative to the valence maximum of \SI{99.4}{eV} ($2p_{3/2}$; $2p_{1/2}$ at \SI{99.8}{eV}), \SI{150}{eV} and \SI{1839}{eV}; in germanium the $3d$ shell sits only \SI{29}{eV} below it~\cite[Table~1-1]{williams2009electron}. Kinematics decides when they matter. A transition from a shell at depth $E_{\mathrm{core}}$ needs $\omega \ge \Eg + E_{\mathrm{core}}$, and since $\vmin = \omega/q + q/2m_\chi \le v_{\mathrm{max}}$ it also needs $q \ge \omega/v_{\mathrm{max}}$ and $m_\chi \ge 2\omega/v_{\mathrm{max}}^2$ (\cref{sec:1.3}). With $v_{\mathrm{max}} = 2.650\times10^{-3}$, the silicon $2p$ shell opens at $m_\chi = \SI{29}{MeV}$ and needs $q \ge \SI{38}{keV}$, and the $2s$ shell at \SI{43}{MeV} with $q \ge \SI{57}{keV}$; the germanium $3d$ shell opens at \SI{8.5}{MeV} and needs $q \ge \SI{11}{keV}$. Silicon $2p$ transitions therefore probe the wavefunctions at momenta beyond the reach of a typical plane-wave basis (\cref{sec:3.5}), and a pseudopotential calculation does not contain the shell at all.

Such shells are either kept in the pseudopotential or taken from the atom. A \emph{semicore} shell, such as germanium $3d$, can be kept in the valence of the pseudopotential, at the price of a higher cutoff, and is then computed with the other bands and reconstructed as above. A deeper shell is taken from the isolated atom. Its Bloch states are the lattice sums
\begin{equation}
  \psi_{\kappa nlm,\vect{k}}(\vect{r}) = \frac{1}{\sqrt N}\sum_{\vect{R}}\ee^{\ii\vect{k}\cdot(\vect{R}+\vect{\tau}_\kappa)}\,\psi^{\mathrm{at}}_{\kappa nlm}(\vect{r} - \vect{R} - \vect{\tau}_\kappa) ,
  \label{eq:E.7-core}
\end{equation}
of the atomic orbitals $\psi^{\mathrm{at}}_{\kappa nlm}$~\cite[Eq.~(2.8)]{griffin2021extended}, labeled by the atomic quantum numbers instead of a band index; the orbitals barely overlap, so the band is flat and its energy is the atomic level. The atomic approximation can be tested where both treatments exist. Ref.~\cite[Sect.~2.2]{griffin2021extended} computes the silicon $2p$ and germanium $3d$ shells both ways and finds that the atomic orbitals reproduce the reconstructed crystal wavefunctions up to the largest momentum the crystal calculation resolves, about \SI{45}{keV}. The vertex between a core state and a conduction state is then $\int\dd^3r\,\psi^*_{c,\vect{k}+\vect{q}}\ee^{\ii\vect{q}\cdot\vect{r}}\psi_{\kappa nlm,\vect{k}}$, with the conduction state from the crystal calculation; above the highest computed band the final states are approximated by plane waves, in the spirit of \cref{eq:2.2-fion-pw}, with or without a Coulomb factor~\cite[Sects.~2.3 and~3.4]{griffin2021extended}. Relativistic corrections to the core orbitals are of relative size $(Z\alpha)^2$, about \SI{1}{\percent} for the silicon $1s$ shell and \SI{5}{\percent} for germanium $1s$. That is the size of each term; the net shift of the $1s$ level is about a quarter of it, \SI{0.26}{\percent} for silicon and \SI{1.46}{\percent} for germanium (\cref{we:3-Zalpha}). The approximation leaves out the relaxation of the other electrons around the core hole, which shifts and reshapes the edge; that final-state effect lies beyond $\mathcal{X} = \mathrm{IP}$.

\subsection{Atomic ionization form factors}

For a target made of isolated atoms, such as a noble liquid, the rate of \cref{sec:2.2} needs the ionization form factor $|f_{\mathrm{ion}}(p_{\mathrm{ej}},q)|^2$ of \cref{eq:2.2-fion} for each occupied shell. The initial orbitals are Hartree--Fock orbitals (\cref{sec:3.4}), almost always the Roothaan--Hartree--Fock solutions of Bunge and collaborators~\cite{bunge1993roothaan}, which expand each radial function in Slater-type orbitals,
\begin{equation}
  R^{\mathrm{at}}_{nl}(r) = a_0^{-3/2}\sum_j C_{jln}\,\frac{(2Z_{jl})^{\tilde n_{jl}+1/2}}{\sqrt{(2\tilde n_{jl})!}}\Bigl(\frac{r}{a_0}\Bigr)^{\tilde n_{jl}-1}\ee^{-Z_{jl}r/a_0} ,
  \label{eq:E.7-sto}
\end{equation}
with tabulated coefficients $C_{jln}$, exponents $Z_{jl}$ and integers $\tilde n_{jl}$~\cite[Eq.~(B33)]{catena2020atomic}. The final state is a continuum state of energy $E_{\mathrm{ej}} = p_{\mathrm{ej}}^2/2m_e$, expanded in partial waves $l'$. In the simplest and most common choice it is the exact Coulomb wave of a hydrogenic potential $-Z_{\mathrm{eff}}\alpha/r$, with $Z_{\mathrm{eff}} = n\,(E_B/\SI{13.6}{eV})^{1/2}$ fixed by the binding energy $E_B$ of the shell~\cite[Eqs.~(B34) and~(B35)]{catena2020atomic}; a better choice solves the radial equation in the self-consistent potential of the ion. The radial integrals of $j_L(qr)$ between the two radial functions are then evaluated for every $(l, l', L)$ allowed by the Gaunt coefficients.

The convergence parameters are the partial-wave cutoff $l'_{\mathrm{max}}$, the grids in $q$ and $E_{\mathrm{ej}}$, and the accuracy of the radial integrals. Ref.~\cite[footnote~11]{catena2020atomic} sums to $l'_{\mathrm{max}} = 7$; the number needed grows with $q$, since a plane wave of momentum $q$ reaches angular momenta up to about $q$ times the orbital radius, and has to be checked at the largest $q$ of the table. The $q$ grid must extend to the kinematic maximum $2m_\chi v_{\mathrm{max}}$ of the heaviest mass considered, \SI{5.3}{MeV} at \SI{1}{GeV}, and is best logarithmic. The radial integrands oscillate rapidly at large $q$, and an integration split into intervals can misread an accidental cancellation over one interval as convergence; Ref.~\cite[Appendix~B\,5]{catena2020atomic} therefore requires two consecutive intervals to agree and switches to a series representation at large $q$. The cost is negligible next to any crystal calculation: a few shells, a few partial waves, and one-dimensional integrals.

The approximation behind the table is that of frozen orbitals. The binding energies are single-particle (Koopmans) energies, the other electrons do not relax when one is removed, and the final state does not feel the polarization of the ion. It is also nonrelativistic, which matters for the inner shells of xenon, where $(Z\alpha)^2 \approx 0.16$. A subtler problem is orthogonality. A continuum state of the model potential $-Z_{\mathrm{eff}}\alpha/r$ is not an eigenstate of the Hartree--Fock Hamiltonian and is not orthogonal to the bound orbital, so that $f_{\mathrm{ion}}$ does not vanish as $q \to 0$, as it must for $\langle f|i\rangle = 0$. Orthogonalizing the continuum state to the occupied orbitals removes the spurious small-$q$ weight, which a light mediator would otherwise amplify. For silicon these atomic tables play a second role: the deep shells of the crystal in \cref{eq:E.7-core} are built from the same orbitals.

\subsection{Convergence parameters and checks}

A crystal form factor is controlled by the cutoff, the number of conduction bands, the $\vect{k}$ grid and the treatment of the energy delta function. The cutoff $E_{\mathrm{cut}}$ sets the largest momentum the vertex contains, $(2m_eE_{\mathrm{cut}})^{1/2}$, which is \SI{31}{keV} at \SI{70}{Ry}; Ref.~\cite[Appendix~C]{essig2016direct} varies it from 30 to \SI{110}{Ry} and finds the silicon rate at $m_\chi = \SI{1}{GeV}$ changed by less than one percent above \SI{70}{Ry}. The number of conduction bands $N_c$ sets the largest energy deposition the spectrum contains; with 52 conduction bands the lowest band left out lies about \SI{57}{eV} above the gap, and the computed spectrum ends there~\cite[Sect.~4.2]{essig2016direct}. A spectrum above the highest band is zero, not small, and a rate integrated over $\omega$ is missing everything beyond it. The $\vect{k}$ grid sets both the resolution in $q$ and the sampling of the onset near $E_{\mathrm{min}}(q)$; Ref.~\cite[Appendix~C]{essig2016direct} finds the grid matters only at low masses and low thresholds, where few transitions are kinematically allowed. It tests the grid by shaking every grid point randomly by up to half the grid spacing and comparing ten such grids. Finally, the delta function of energy needs a broadening or a tetrahedron integration, with the width tied to the grid spacing as described in \cref{app:C.2}, or a binning: Ref.~\cite[Sect.~4.2]{essig2016direct} uses bins of \SI{0.1}{eV} in $\omega$ and $0.02\,\alpha m_e = \SI{75}{eV}$ in $q$.

Most errors are caught by the small-$q$ limit and by a sum rule. The small-$q$ check is the normalization $\rho_{nn}(\vect{k},0) = 1$ with $\rho_{mn}(\vect{k},\vect{q})$ for $m \neq n$ vanishing linearly in $q$, with the slope given by the velocity, as discussed above. The sum rule holds for each occupied state and measures the effect of truncating the bands. Since $\rho_{mn}(\vect{k},\vect{q}) = \langle u_{m,\vect{k}+\vect{q}}|u_{n\vect{k}}\rangle$ (\cref{eq:A.5-rho}), and the $u_{m,\vect{k}+\vect{q}}$ form a complete set of eigenstates of $H(\vect{k}+\vect{q}) = \ee^{-\ii(\vect{k}+\vect{q})\cdot\vect{r}}H\ee^{\ii(\vect{k}+\vect{q})\cdot\vect{r}}$,
\begin{equation}
  \sum_m(\varepsilon_{m,\vect{k}+\vect{q}} - \varepsilon_{n\vect{k}})\,|\rho_{mn}(\vect{k},\vect{q})|^2 = \langle u_{n\vect{k}}|H(\vect{k}+\vect{q}) - H(\vect{k})|u_{n\vect{k}}\rangle
  = \vect{q}\cdot\vect{v}^{\mathrm{g}}_n(\vect{k}) + \frac{q^2}{2m_e} + \delta_{\mathrm{NL}} ,
  \label{eq:E.7-fsum}
\end{equation}
where the sum over $m$ runs over all bands and $\delta_{\mathrm{NL}}$ is the part of $\langle u_{n\vect{k}}|V_{\mathrm{NL}}(\vect{k}+\vect{q}) - V_{\mathrm{NL}}(\vect{k})|u_{n\vect{k}}\rangle$ beyond first order in $q$. The kinetic part of $H(\vect{k}) = (\vect{p}+\vect{k})^2/2m_e + V_{\mathrm{loc}} + V_{\mathrm{NL}}(\vect{k})$ gives $\vect{q}\cdot(\vect{p}+\vect{k})/m_e + q^2/2m_e$, and the first-order part of the nonlocal term combines with it into the group velocity, by the Hellmann--Feynman relation $\nabla_{\vect{k}}\varepsilon_{n\vect{k}} = \langle u_{n\vect{k}}|\nabla_{\vect{k}}H(\vect{k})|u_{n\vect{k}}\rangle$. Summed over a filled band the linear term cancels. Transitions between two occupied states cancel in pairs, since time reversal gives $|\rho_{nm}(-\vect{k}-\vect{q},\vect{q})| = |\rho_{mn}(\vect{k},\vect{q})|$ with the opposite energy difference. What remains is the valence $f$-sum rule of a pseudo-Hamiltonian, the band-summed form of \cref{eq:E.2-dc}: $q^2/2m_e$ per valence electron, plus the nonlocal term, which is absent for the all-electron Hamiltonian of \cref{eq:A.7-fsum}. The right side of \cref{eq:E.7-fsum} is a ground-state expectation value and costs nothing; the left side with $N_c$ conduction bands falls short of it by the weight of the bands left out. With a nonlocal pseudopotential, a test that compares the truncated sum with $q^2/2m_e$ alone confuses the two effects, and the nonlocal term must be included before the shortfall is read as band truncation. The shortfall falls slowly with $N_c$, because at momentum transfer $q$ much of the missing weight sits near the free-electron energy $q^2/2m_e$.

\subsection{Cost and pitfalls}

The count of vertices sets the cost. With both $\vect{k}$ and $\vect{k}'$ on a grid of $n_k$ points, $N_v$ valence and $N_c$ conduction bands, and $N_{\vect{G}}$ reciprocal-lattice vectors within the largest $q$, there are $n_k^2N_vN_cN_{\vect{G}}$ vertices, computed with $n_k^2N_vN_c$ fast Fourier transforms of $N_{\mathrm{FFT}}$ points each. The number of reciprocal-lattice vectors within $q_{\mathrm{max}}$ is $\frac{4\pi}{3}q_{\mathrm{max}}^3\,\Omega/(2\pi)^3$. Directional averages can be taken over the irreducible wedge of $\vect{k}$, with the full set of $\vect{k}'$. The vertices themselves are far too many to store, and they are accumulated into bins of $(q,\omega)$, or of $(\vect{q},\omega)$ for a directional response, as they are computed; a tabulated crystal form factor is that binned sum. The all-electron correction adds, per state, the projections onto the atomic projectors, and per $q$ the radial integrals of \cref{eq:E.7-radial}.

Three errors are not caught by the checks above. A $\vect{q}$ reduced to the first zone assigns an Umklapp vertex to the wrong momentum transfer (\cref{app:A.3}). A spectrum cut by the highest band, or a $q$ range cut by $E_{\mathrm{cut}}$, looks converged in its own window and silently loses everything outside it. And normalizations differ between codes and papers: plane-wave codes often fold the spin factor 2 into the $\vect{k}$ weights~\cite[Eq.~(4.2)]{essig2016direct}, and the conversions between $f_{\mathrm{crys}}$ and $\bar S$ are in \cref{eq:A.8-fcrys}. For germanium, which semilocal functionals leave without a gap~\cite[Sect.~2.1]{griffin2021extended}, the band order must be fixed before any rigid shift (\cref{app:E.1}).

\begin{workedexample}{The silicon form factor}\label{we:E-siff}
We set up a valence-to-conduction calculation for silicon at $\mathfrak{T}_0$ and count its cost. The cell volume is $\Omega = a^3/4 = \SI{40.05}{\angstrom^3}$ at $a = \SI{5.431}{\angstrom}$. A cutoff of \SI{70}{Ry}, \SI{952}{eV}, reaches $(2m_eE_{\mathrm{cut}})^{1/2} = \SI{31.2}{keV} = \SI{15.8}{\angstrom^{-1}}$ and holds $N_{\mathrm{PW}} = \Omega(2m_eE_{\mathrm{cut}})^{3/2}/6\pi^2 \approx 2700$ plane waves per state. The non-self-consistent run computes 4 valence and 52 conduction bands on a $10\times10\times10$ grid without symmetry, $n_k = 1000$; the coefficients take $1000\times56\times2700$ complex numbers, \SI{2.4}{GB} in double precision, which fits in memory. The Kohn--Sham gap is about \SI{0.6}{eV} in the representative PBE calculation (\cref{we:E-si-bands}), and a scissor $\Delta_{\mathrm{scis}} = \Eg - E_{\mathrm{g}}^{\mathrm{KS}}$ moves the conduction bands up to the measured low-temperature gap, \SI{1.17}{eV} (\cref{sec:1.3}).

The reciprocal-lattice vectors within $q_{\mathrm{max}} = \SI{30}{keV}$ number $\frac{4\pi}{3}q_{\mathrm{max}}^3\Omega/(2\pi)^3 \approx 2400$, since the zone of silicon has volume $(2\pi)^3/\Omega = \SI{47.6}{keV^3}$. The vertex count is $n_k^2N_vN_cN_{\vect{G}} = 10^6\times4\times52\times2400 \approx 5\times10^{11}$, the number quoted at the start of this section. They come from $n_k^2N_vN_c = 2.1\times10^8$ transforms on grids of $N_{\mathrm{FFT}} \approx 22N_{\mathrm{PW}} \approx 6\times10^4$ points. At about $5N_{\mathrm{FFT}}\log_2N_{\mathrm{FFT}} \approx 5\times10^6$ floating-point operations per transform, the total is $1\times10^{15}$ operations, about three hours on a node that sustains $10^{11}$ operations per second, and less with the irreducible wedge. The calculation is cheap; the binned output, at the bin widths quoted above, has about $500\times400$ entries up to \SI{30}{keV} (the published table of Ref.~\cite{essig2016direct} has $500\times900$ and reaches \SI{67}{keV}).

The calculation is then tested at small $q$, against the sum rule and against the cutoff. At $q = 10^{-3}\,|\vect{b}_i|$ the ratio $|\rho_{cv}|/q$ must equal $|\hat{\vect{q}}\cdot\vect{v}_{cv}|/(\varepsilon_c - \varepsilon_v)$ with the velocity of \cref{eq:E.7-vNL}. For each valence state the truncated left side of \cref{eq:E.7-fsum}, with its nonlocal term on the right, measures the band truncation. And the rate at $m_\chi = \SI{1}{GeV}$ for a heavy mediator, recomputed at a higher cutoff and with all-electron reconstruction, fixes where the pseudo vertex stops being adequate. Kinematics tells us where to look: at \SI{10}{MeV}, with $\omega$ just above the gap, momentum transfers between \SI{0.45}{keV} and \SI{53}{keV} are allowed, but most of the valence weight lies at a few keV, and the cutoff test decides how much of the rate comes from beyond \SI{30}{keV}; the $2p$ shell, which would need $q \ge \SI{38}{keV}$, is still closed below \SI{29}{MeV}.
\end{workedexample}

The independent-particle form factor is where every electronic rate of \cref{chap:10} starts. The next corrections, screening and local fields on the many-body axis $\mathcal{X}$, use the same vertices as input; \cref{app:E.6} treated them.

\section{Anharmonic and finite-temperature methods}\label{app:E.8}

The band gap of silicon near zero temperature is about \SI{60}{meV} smaller than the gap of the same crystal with its nuclei clamped at their equilibrium positions: the zero-point motion of the lattice closes it, and the measured values are 62 and \SI{64}{meV}~\cite[Table~VI]{ponce2015temperature}. Between zero and room temperature the gap closes by a further \SI{50}{meV} (\cref{sec:1.3}), and the phonons themselves shift and broaden. A detector at \SI{10}{mK} sits deep in the zero-point regime, where none of these effects are thermal and all of them are quantum. We will see how a gap shift of tens of meV and the phonon shifts behind it are computed: some methods average the energy surface itself, and others average an electronic property over the zero-point distribution of the nuclei. All of them start from the harmonic phonons, Born charges and electron--phonon potentials of density-functional perturbation theory (DFPT, \cref{app:E.4}).

\subsection{Place on the hierarchy}

Free energies, effective force constants, band energies and absorption spectra at temperature $T$ are all averages over the distribution of nuclear positions, either of the energy surface itself or of an electronic property computed at fixed nuclei. In the language of the truncation label $\mathfrak{T}$ of \cref{sec:8.5}, the lattice-side methods move the anharmonic order $n_A$ and the thermal order $n_T$ away from the base point $\mathfrak{T}_0$, while keeping the Born--Oppenheimer order at $n_B = 0$. The quasi-harmonic approximation keeps every mode harmonic and lets the frequencies depend on volume, the lowest-order effect of the cubic terms on the free energy. The self-consistent harmonic approximation resums the even-order anharmonic terms into an effective harmonic problem and is not a fixed order in $n_A$, and the temperature-dependent effective potential fits effective second- and third-order force constants to a sampled ensemble. The electron-side methods move the electron--phonon order: AHC theory is the $n_g = 2$ correction to the band energies, with two first-order vertices (the Fan--Migdal term) or one second-order vertex (the Allen--Heine term), in the harmonic approximation. The special-displacement method evaluates the adiabatic average of the electronic property to all orders in the displacements, at $n_A = 0$ and in the adiabatic limit. The anharmonic small parameter $\eta_{\mathrm{anh}}$ of \cref{sec:5.7} is small for silicon, and the lattice-side corrections are correspondingly small; the electron-side ones are not.

\subsection{Quasi-harmonic approximation}

In the quasi-harmonic approximation the free energy per cell, $\mathcal{F}$, at cell volume $\Omega$ is the harmonic one, with frequencies computed at that volume~\cite[Eq.~(4)]{baroni2010thermal},
\begin{equation}
  \mathcal{F}(\Omega,T) = E^{\mathrm{ad}}_0(\Omega) + \frac1N\sum_{\vect{Q}\nu}\frac{\omega_{\vect{Q}\nu}(\Omega)}{2} + \frac TN\sum_{\vect{Q}\nu}\ln\bigl[1 - \ee^{-\omega_{\vect{Q}\nu}(\Omega)/T}\bigr] ,
  \label{eq:E.8-qha}
\end{equation}
where $E^{\mathrm{ad}}_0(\Omega)$ is the clamped-nucleus energy per cell, the minimum of the energy surface of \cref{sec:5.1} at that volume, and the sums run over the $N$ points of the phonon grid. The equilibrium volume minimizes $\mathcal{F}$ at each $T$. Differentiating, the volume thermal expansion coefficient is
\begin{equation}
  \alpha_{\mathrm{th}} = \frac{1}{N B\Omega}\sum_{\vect{Q}\nu}\omega_{\vect{Q}\nu}\,\gamma^{\mathrm{G}}_{\vect{Q}\nu}\,\frac{\partial n_{\vect{Q}\nu}}{\partial T}, \qquad
  \gamma^{\mathrm{G}}_{\vect{Q}\nu} = -\frac{\Omega}{\omega_{\vect{Q}\nu}}\frac{\partial\omega_{\vect{Q}\nu}}{\partial\Omega} ,
  \label{eq:E.8-beta}
\end{equation}
with $B$ the isothermal bulk modulus and $\gamma^{\mathrm{G}}_{\vect{Q}\nu}$ the mode Grüneisen parameters~\cite[Eqs.~(6) and~(9)]{baroni2010thermal}. A perfectly harmonic crystal has $\gamma^{\mathrm{G}} = 0$ and no thermal expansion. The calculation is a set of harmonic calculations: DFPT phonons at five to seven volumes around the minimum, a fit of $E^{\mathrm{ad}}_0(\Omega)$ and of each $\omega_{\vect{Q}\nu}(\Omega)$, and a minimization of \cref{eq:E.8-qha}. The cost is that number of DFPT runs. In silicon the transverse acoustic branch has negative Grüneisen parameters, and its large weight in the density of states at low frequency makes the thermal expansion negative at low temperature, a behavior the quasi-harmonic approximation reproduces~\cite[Sect.~5.1]{baroni2010thermal}.

The approximation holds as long as the frequencies depend on the volume but not on the amplitude of vibration, that is, as long as the phonons are well-defined harmonic excitations at each volume. It fails for modes that are unstable in the harmonic approximation, near structural transitions and near melting, which is where the next two methods take over. The zero-point expansion and the sensitivity of the phonons to the lattice constant matter even at \SI{10}{mK}. The zero-point term of \cref{eq:E.8-qha} expands the lattice at $T = 0$, and this expansion changes the gap: for silicon the zero-point lattice expansion alone shifts the gap by $+\SI{9}{meV}$~\cite[Table~S2]{miglio2020predominance}, opposite to the electron--phonon renormalization. And the same Grüneisen parameters set how sensitive a phonon calculation is to the lattice constant it uses: by \cref{eq:E.8-beta}, a relative error $\delta a/a$ in the lattice constant shifts each frequency by $\delta\omega_{\vect{Q}\nu}/\omega_{\vect{Q}\nu} = -3\gamma^{\mathrm{G}}_{\vect{Q}\nu}\,\delta a/a$. In silicon the transverse acoustic modes, with $\gamma^{\mathrm{G}} < 0$, and the optical modes, with $\gamma^{\mathrm{G}} > 0$, move in opposite directions, so that no single choice between the theoretical and the measured lattice constant improves both. A calculation that needs the transverse acoustic branch accurately, as the phonon-assisted thresholds of \cref{chap:13} do, has to test both lattice constants and quote the spread.

\subsection{Self-consistent harmonic approximation}

When the anharmonicity is not small, the harmonic expansion about the minimum is the wrong starting point, and the self-consistent harmonic approximation (SCHA) replaces it by the best harmonic description of the actual distribution of nuclei. Its basis is the Gibbs--Bogoliubov inequality: for any trial density matrix $\tilde\varrho$, the functional $\mathcal{F}[\tilde\varrho] = \langle H\rangle_{\tilde\varrho} - TS_{\mathrm{ent}}[\tilde\varrho]$, with $S_{\mathrm{ent}}$ the entropy of $\tilde\varrho$, is an upper bound on the true free energy~\cite[Eq.~(2)]{monacelli2021stochastic}. The trial density matrices are those of harmonic Hamiltonians,
\begin{equation}
  H_{\mathrm{trial}} = (\text{nuclear kinetic energy}) + \frac12\sum_{ab}(X_a - \bar X_a)\,\Phi^{(2)}_{\mathrm{aux},ab}\,(X_b - \bar X_b) ,
  \label{eq:E.8-scha}
\end{equation}
with $X_a$ the $3N_{\mathrm{at}}$ Cartesian coordinates of the nuclei of a supercell of $N_{\mathrm{at}}$ atoms, and the centroids $\bar X_a$ and the auxiliary force constants $\Phi^{(2)}_{\mathrm{aux}}$ as variational parameters~\cite[Eqs.~(4) and~(7)]{monacelli2021stochastic}. Minimizing $\mathcal{F}$ over $\bar X$ and $\Phi^{(2)}_{\mathrm{aux}}$ gives the SCHA free energy, the equilibrium positions at temperature $T$, and an effective harmonic problem whose quantum zero-point motion is included by construction. At the minimum $\Phi^{(2)}_{\mathrm{aux}}$ satisfies the self-consistency condition of \cref{sec:5.7}: it equals the average of the true second derivatives of the energy surface over the Gaussian distribution it generates.

The averages are computed stochastically. A population of $N_{\mathrm{conf}}$ configurations is drawn in a supercell from the Gaussian distribution of the current trial Hamiltonian, each together with its mirror image about the centroids to reduce the statistical error; DFT supplies the energy, forces and stress of each; and the gradients of $\mathcal{F}$ with respect to $\bar X$ and $\Phi^{(2)}_{\mathrm{aux}}$ follow from these forces. As the parameters move, the same configurations are reweighted by the ratio of the new to the old Gaussian density, until the effective number of samples, $N_{\mathrm{eff}} = (\sum_jw_j)^2/\sum_jw_j^2$ for weights $w_j$, drops below about half of $N_{\mathrm{conf}}$; then a new population is drawn~\cite[Eqs.~(21), (38) and~(39)]{monacelli2021stochastic}. The cost is entirely in the DFT forces: a population holds several hundred to several thousand configurations, and most materials converge within three populations~\cite[Sect.~III\,B]{monacelli2021stochastic}. The convergence parameters are the supercell size, $N_{\mathrm{conf}}$, and the threshold on $N_{\mathrm{eff}}/N_{\mathrm{conf}}$.

The auxiliary frequencies of $\Phi^{(2)}_{\mathrm{aux}}$ are not the physical phonon energies. They describe the Gaussian distribution of the nuclei; the frequencies measured by a probe at small energy follow from the curvature of the free energy with respect to the centroids, which adds a term with third-order force constants~\cite[Sect.~IV]{monacelli2021stochastic}. Mistaking one for the other is the most common error in reading SCHA output. For silicon the SCHA changes the phonons little, since $\eta_{\mathrm{anh}}$ is small; it is needed for light atoms, soft modes and quantum crystals, the targets of \cref{sec:11.5}.

\subsection{Temperature-dependent effective potential}

The temperature-dependent effective potential (TDEP) method fits effective force constants to forces sampled at temperature $T$. Molecular dynamics in a supercell of $N_{\mathrm{at}}$ atoms supplies, at each of $N_t$ time steps, the displacements $\vect{u}_i$ and the DFT forces $\vect{F}_i$ on every atom $i$ of the supercell. The effective second-order force constants $\tilde\Phi^{(2)}_{ij}$ are those that minimize the mismatch
\begin{equation}
  \Delta_{\mathrm{fit}} = \frac{1}{N_t}\sum_{t,i}\Bigl|\vect{F}_{i}(t) + \sum_j\tilde\Phi^{(2)}_{ij}\,\vect{u}_j(t)\Bigr|^2 ,
  \label{eq:E.8-tdep}
\end{equation}
since a harmonic energy surface gives $\vect{F}_i = -\sum_j\tilde\Phi^{(2)}_{ij}\vect{u}_j$. This is a linear least-squares problem, overdetermined once $N_t > 3N_{\mathrm{at}}$~\cite[Eq.~(3), with the opposite sign convention for the force constants]{hellman2011lattice}. Symmetry reduces the unknowns to a few per coordination shell, and the same fit extends to third and fourth order. The fitted constants absorb the anharmonicity sampled at that temperature into effective harmonic ones; the phonons and free energy follow as in the harmonic case~\cite{hellman2013temperature}. The convergence parameters are the number of time steps, which Ref.~\cite[Fig.~1]{hellman2011lattice} tracks through the change in the fitted constants, the supercell size, and the range of the force constants kept in the fit.

The method inherits the statistics of its sampling. Classical molecular dynamics samples a classical distribution, with mean squared displacement $T/(M_\kappa\omega_{\vect{Q}\nu}^2)$ per mode, while the quantum harmonic value is $(2n_{\vect{Q}\nu} + 1)/(2M_\kappa\omega_{\vect{Q}\nu})$. Their ratio is $2T/\omega_{\vect{Q}\nu}$ for $T \ll \omega_{\vect{Q}\nu}$. For silicon phonons of 10 to \SI{60}{meV} at \SI{10}{mK}, the ratio is below $2\times10^{-4}$, so a classical sample sees an almost static lattice and misses the zero-point motion that sets every renormalization at detector temperatures. At low temperature the configurations must be drawn from the quantum distribution instead, as the SCHA does, or from path-integral molecular dynamics (\cref{app:E.11}). The cost of TDEP is that of the molecular dynamics, $N_t$ force calculations of which only those separated by more than a correlation time are independent.

Third-order force constants also give the three-phonon linewidths $\gamma^{\mathrm{anh}}_{\vect{Q}\nu}$ of \cref{sec:5.7}. They are fitted in the same way, or computed by finite differences in supercells, one supercell calculation per pair of displaced atoms allowed by symmetry, with a cutoff radius on the interacting triplets. The linewidth then needs a sum over pairs of final phonons on a dense $\vect{Q}$ grid with an energy delta function, and the same rules for broadening and grid spacing apply as for electronic spectra (\cref{app:C.2}).

\subsection{Allen--Heine--Cardona calculations}

The AHC theory of \cref{sec:6.3} gives the shift of a band energy at second order in the electron--phonon coupling, summed over all phonons with the thermal weight $2n_{\vect{Q}\nu} + 1$. In the adiabatic form, which neglects the phonon energy next to electronic energy differences, the Fan--Migdal part is
\begin{equation}
  \Delta\varepsilon^{\mathrm{FM}}_{n\vect{k}}(T) = \frac{1}{N}\sum_{\vect{Q}\nu}(2n_{\vect{Q}\nu} + 1)\sum_m\frac{|g_{mn\nu}(\vect{k},\vect{Q})|^2}{\varepsilon_{n\vect{k}} - \varepsilon_{m,\vect{k}+\vect{Q}}} ,
  \label{eq:E.8-fan}
\end{equation}
with the vertex of \cref{eq:A.6-g}, and the Allen--Heine part adds the diagonal element of the second-order vertex $g^{\mathrm{AH}}$ with the same weight. In the rigid-ion approximation the second-order vertex is rewritten through first-order matrix elements at $\vect{Q} = 0$, using the translation invariance of the crystal, so that no second derivative of the potential is computed~\cite[Eqs.~(11) and~(15)]{ponce2015temperature} (\cref{sec:6.3}). The AHC literature calls this term the Debye--Waller term; in this book Debye--Waller refers only to $\ee^{-2W_\kappa}$ (\cref{app:A.6}). The nonadiabatic form keeps $\pm\omega_{\vect{Q}\nu}$ in the denominators of \cref{eq:E.8-fan}, separately for phonon emission and absorption, and gives the renormalization a lifetime as well~\cite[Eq.~(16)]{ponce2015temperature}.

The sum over $m$ runs over all bands. In practice it is split: the bands within an active window are summed explicitly, and the contribution of all higher bands is obtained from a Sternheimer equation for the first-order change of the wavefunction projected on the space above the window, as DFPT does for phonons (\cref{app:E.4}). The split is exact in the adiabatic limit; in the nonadiabatic form the Sternheimer part drops the phonon energy, and the active window must be large enough for that to be harmless~\cite[below Eq.~(16)]{ponce2015temperature}.

The sum over $\vect{Q}$ is the hard part. The Fan--Migdal denominators vanish when $\varepsilon_{m,\vect{k}+\vect{Q}} = \varepsilon_{n\vect{k}}$, which happens for $m = n$ at $\vect{Q} \to 0$ and wherever another band is degenerate with the state. Near these points the integrand is singular, and the result converges slowly with the number $N_Q^3$ of grid points, at a rate that depends on the band and on the form of the theory; a small imaginary part $\ii\delta$ in the denominators smooths the integrand, at the price of a bias that must be extrapolated to $\delta \to 0$ together with $1/N_Q \to 0$~\cite[Sect.~III and Table~II]{ponce2015temperature}. A nonzero Born effective charge makes the vertex diverge as $1/Q$ at small $Q$. In a nonpolar crystal the charges vanish by symmetry, but a DFPT calculation always leaves small residual charges that violate the charge-neutrality sum rule, and these must be removed before the $\vect{Q}$ sum, or the integral diverges at band extrema~\cite[Sect.~III]{ponce2015temperature}. In a polar crystal the divergence is physical: the adiabatic integral does not converge at all, and only the nonadiabatic form gives a finite result~\cite{ponce2015temperature}. For polar materials the nonadiabatic, Fröhlich-like part often dominates the zero-point renormalization~\cite{miglio2020predominance}.

The cost is that of DFPT at every $\vect{Q}$ of a dense grid: for silicon, Ref.~\cite{ponce2015temperature} goes up to a $100\times100\times100$ grid, 22\,776 irreducible points, to converge the values of its Table~VI. Wannier interpolation of the vertices (\cref{app:E.5}) reduces this to DFPT on a coarse grid. The results for silicon are in \cref{we:E-zpr} below. The calculations use Kohn--Sham states; since a semilocal functional underestimates the gap, it overscreens the electron--phonon interaction, and DFT-based renormalizations come out smaller than $GW$-based ones~\cite[Sect.~VII\,B]{zacharias2020theory}.

\subsection{Special displacements}

The perturbative route needs vertices and a converged sum over phonon wavevectors. The non-perturbative alternative computes the electronic property directly in a supercell whose nuclei are displaced, and averages over displacements. In the adiabatic and harmonic approximations, the Williams--Lax expression for a temperature-dependent property, here the imaginary part of the dielectric function, is a Gaussian average over the normal coordinates $x_\nu$ of the supercell,
\begin{equation}
  \varepsilon_2(\omega;T) = \prod_{\vect{Q}\nu}\int\dd x_{\vect{Q}\nu}\,\frac{\ee^{-x_{\vect{Q}\nu}^2/2\sigma_{\vect{Q}\nu}^2}}{\sqrt{2\pi}\,\sigma_{\vect{Q}\nu}}\;\varepsilon_2(\omega;\{x\}), \qquad \sigma_{\vect{Q}\nu}^2 = \frac{2n_{\vect{Q}\nu} + 1}{2\omega_{\vect{Q}\nu}} ,
  \label{eq:E.8-wl}
\end{equation}
where the $x_{\vect{Q}\nu}$ are the mass-weighted normal coordinates of the supercell, labeled by the wavevectors $\vect{Q}$ of the primitive-cell phonons it contains, and $\varepsilon_2(\omega;\{x\})$ is computed with the nuclei frozen at the displaced positions~\cite[Eqs.~(3) and~(4), with mass-weighted coordinates]{zacharias2016one}. The same formula holds for a band energy or any other property computed at fixed nuclei. Earlier calculations evaluated it by Monte Carlo over configurations; for the optical spectrum of silicon fewer than ten samples sufficed, and a single sample came close~\cite[Sect.~II]{zacharias2016one}.

Zacharias and Giustino explained why, and replaced the sampling by one deterministic configuration. In a large supercell, the average of a smooth function over the Gaussian distribution of \cref{eq:E.8-wl} is dominated by configurations in which every normal coordinate has the magnitude $\sigma_{\vect{Q}\nu}$, and the cross terms between pairs of modes cancel if the signs alternate in a suitable order~\cite[Sect.~III]{zacharias2016one}. The resulting special, or ZG, displacement of nucleus $\kappa$ in cell $\vect{R}$ of a supercell of $N$ primitive cells is
\begin{equation}
  \vect{u}^{\mathrm{ZG}}_{\vect{R}\kappa} = \sum_{\vect{Q}\in\mathcal{B},\nu}\varsigma_{\vect{Q}\nu}\left[\frac{2n_{\vect{Q}\nu} + 1}{2NM_\kappa\omega_{\vect{Q}\nu}}\right]^{1/2}2\,\Real\bigl[\ee^{\ii\vect{Q}\cdot\vect{R}}\,\vect{e}_{\kappa\nu}(\vect{Q})\bigr] ,
  \label{eq:E.8-zg}
\end{equation}
in the conventions of \cref{eq:A.6-u}~\cite[Eq.~(2)]{zacharias2020theory}. The set $\mathcal{B}$ contains one wavevector from each pair $\pm\vect{Q}$ of the grid commensurate with the supercell, and the signs $\varsigma_{\vect{Q}\nu} = \pm1$ alternate from one mode to the next in a prescribed order. Wavevectors equal to their own time-reversal partners, such as $\vect{Q} = 0$, are left out: they describe standing waves that break the crystal symmetry, and their weight vanishes as the supercell grows~\cite[Sects.~II and~VII\,B]{zacharias2020theory}. The mean squared displacement of each atom, averaged over the supercell, then equals its thermal value from \cref{eq:A.6-u}. The method needs the phonons of the primitive cell on the commensurate grid, from DFPT, and one DFT calculation of the displaced supercell with an electronic $\vect{k}$ sampling appropriate to its size.

The convergence parameter is the supercell size. Including or excluding the wavevectors with $\vect{Q} = -\vect{Q}$ modulo $\vect{G}$ brackets the converged value from above and below, and the width of the bracket is a direct error estimate~\cite[Sect.~VII\,B]{zacharias2020theory}. The phases of the eigenvectors must be fixed consistently from mode to mode before the signs are assigned; Ref.~\cite[Sect.~II]{zacharias2016one} chooses the first nonzero component of each eigenvector positive. Since the displaced supercell is a static calculation, it needs no sum over phonon wavevectors with vanishing denominators, and the adiabatic divergence of polar materials is absent: for GaAs the zero-point renormalization converges smoothly with supercell size, though a $6\times6\times6$ supercell, 432 atoms, is needed to come within ten percent~\cite[Sect.~VII]{zacharias2020theory}. The cost is one supercell calculation, which scales as the cube of the number of atoms for the diagonalization; an $8\times8\times8$ supercell of silicon holds 1024 atoms and costs about $8^3 = 512$ times more than a $4\times4\times4$ one at fixed accuracy per atom.

\paragraph{Phonon-assisted absorption.} With \cref{eq:E.8-wl} applied to $\varepsilon_2$, a single ZG supercell gives the temperature-dependent absorption spectrum including indirect transitions. The displaced supercell breaks translational symmetry, so that a photon can connect states that differ in crystal momentum by a phonon wavevector of the supercell grid. For silicon at room temperature, an $8\times8\times8$ supercell with 40 random $\vect{k}$ points, a scissor of \SI{0.74}{eV} and a Gaussian broadening of \SI{30}{meV} reproduces the measured absorption coefficient over seven orders of magnitude~\cite[Sect.~II and~VII]{zacharias2016one}. Being adiabatic, the method does not distinguish the absorption of a phonon from its emission; it cannot reproduce the two thresholds $\Eg \pm \omega_{\vect{Q}\nu}$ of the Hall--Bardeen--Blatt theory near the onset~\cite[Sect.~II]{zacharias2016one}, which is precisely the region that sets the phonon-assisted thresholds of \cref{chap:13}. At \SI{10}{mK} the phonon-absorption branch is absent, and the onset is at $\Eg + \omega_{\vect{Q}\nu}$ (\cref{sec:6.5}). Special displacements therefore give the spectrum and its renormalization away from the edge, while the edge itself needs the nonadiabatic treatment with explicit phonon energies of \cref{app:E.5} (\cref{sec:13.7}).

\paragraph{Zero-point renormalization from special displacements.} The band energies of the displaced supercell, unfolded onto the primitive zone as a spectral function, give the renormalized bands $\varepsilon^{\mathrm{ren}}_{n\vect{k}}(T)$ directly. Unlike the perturbative sum, this includes all orders in the displacements, within the adiabatic approximation, and needs no electron--phonon vertices. It misses the nonadiabatic correction; in silicon the adiabatic and nonadiabatic AHC renormalizations differ by \SI{8}{meV}~\cite[Table~VI]{ponce2015temperature}.

\begin{workedexample}{The zero-point renormalization of the silicon gap}\label{we:E-zpr}
Perturbative AHC theory, the ZG displacement and experiment each give the indirect-gap renormalization of silicon at $T = 0$. Perturbative AHC theory with the local-density approximation, Sternheimer equations for the high bands and grids up to $100\times100\times100$ phonon wavevectors, extrapolated to $\delta \to 0$, gives $-\SI{64.3}{meV}$ in the adiabatic form and $-\SI{56.2}{meV}$ in the nonadiabatic one~\cite[Table~VI]{ponce2015temperature}. The ZG displacement in an $8\times8\times8$ supercell, 1024 atoms, gives $-\SI{57}{meV}$, made of a rise of the valence maximum by \SI{35}{meV} and a drop of the conduction minimum, near $0.87$ of the way from $\mathrm{\Gamma}$ to $\mathrm{X}$, by \SI{22}{meV}; a $10\times10\times10$ supercell brackets the converged value between 55 and \SI{65}{meV}~\cite[Sect.~VII\,B]{zacharias2020theory}. A $4\times4\times4$ supercell, 128 atoms, is already within ten percent of the converged value~\cite[Sect.~VII]{zacharias2020theory}. The measured renormalization is $-62$ or $-\SI{64}{meV}$, depending on how the data are extrapolated~\cite[Table~VI]{ponce2015temperature}, and the zero-point lattice expansion contributes $+\SI{9}{meV}$ of the opposite sign~\cite[Table~S2]{miglio2020predominance}.

The measured value contains both the electron--phonon renormalization and the lattice expansion. The lattice term is therefore added to the computed fixed-volume value before the two are compared~\cite{miglio2020predominance}: $-56 + 9 = -\SI{47}{meV}$ for the nonadiabatic AHC result and $-57 + 9 = -\SI{48}{meV}$ for the ZG displacement, against $-62$ to $-\SI{64}{meV}$. Equivalently, the electron--phonon part implied by the measurement is $-71$ to $-\SI{73}{meV}$, and the calculations underestimate it by about \SI{20}{\percent}. The shortfall is consistent with the overscreening of the vertex by the small Kohn--Sham gap: a calculation of the same kind on $GW$ bands gives \SI{66}{meV}~\cite[Sect.~VII\,B]{zacharias2020theory}. The cost differs by orders of magnitude: tens of thousands of DFPT calculations in the primitive cell for the perturbative route, against DFPT phonons on an $8\times8\times8$ grid and one 1024-atom supercell for the special displacement. A measured low-temperature gap, such as the \SI{1.17}{eV} of silicon (\cref{sec:1.3}), already contains the zero-point renormalization. A calculation that scissors Kohn--Sham bands to it and then adds a computed renormalization counts it twice; the renormalization is added only to a clamped-nucleus gap, from $GW$ or from a hybrid functional. So the window and threshold calculations of \cref{chap:13} take the measured gap and do not add the renormalization again. Between \SI{10}{mK} and \SI{1}{K} all occupations $n_{\vect{Q}\nu}$ of silicon are of order $\ee^{-\omega_{\vect{Q}\nu}/T}$ with $\omega_{\vect{Q}\nu}/T \gtrsim 10^2$ for every phonon above \SI{10}{meV}, and the thermal part of the renormalization is negligible; only the acoustic modes at small $Q$ are thermally occupied, and their contribution to \cref{eq:E.8-fan} vanishes with their phase space. What remains temperature dependent at those temperatures is the population of the final states, not their energies.
\end{workedexample}

\subsection{Pitfalls}

Three numerical effects are easy to mistake for physics: the band splitting from the $\vect{Q} = -\vect{Q}$ modes, the imaginary broadening of the AHC denominators, and the missing edges of phonon-assisted absorption. Special displacements with the $\vect{Q} = -\vect{Q}$ modes included split degenerate bands, such as the threefold valence maximum of silicon, and the splitting is an artifact~\cite[Sect.~VII\,B]{zacharias2020theory}. An imaginary broadening $\ii\delta$ in the AHC denominators is a numerical device; its effect is extrapolated to $\delta \to 0$, never quoted as a lifetime. The adiabatic special displacements also cannot give the edges of phonon-assisted absorption, which is why the edges of \cref{chap:13} use the Wannier route of \cref{app:E.5}.

\section{Magnetic and superconducting parameters}\label{app:E.9}

The magnon band of NiO is set almost entirely by one number, the antiferromagnetic exchange between next-nearest-neighbor nickel ions, and the density-functional values of it collected in one comparison span a factor of two and a half, from the same crystal structure, depending on the functional and on the Hubbard correction used~\cite[Table~III]{fischer2009exchange}. The superconducting gap of aluminum, by contrast, is computed from first principles to within a few percent of the measured \SI{0.18}{meV}~\cite[Table~III]{savrasov1996electron}, but only after one parameter, the Coulomb pseudopotential, is fixed by the measured critical temperature. These two cases mark the range of what first-principles methods deliver for the magnetic and superconducting targets of \cref{sec:7.1,sec:7.2} and of \cref{chap:14,chap:15}. We will see how DFT~$+\,U$ total energies, spin spirals and Green's functions lead to the exchange constants $J^{\mathrm{ex}}_{ij}$ of the Heisenberg model, and how the Eliashberg equations turn the spectral function $\alpha^2F(\omega)$ of \cref{sec:6.4}, built from the vertices of \cref{app:E.4,app:E.5}, into a gap and a critical temperature. Silicon, the running example of the book, has neither magnetic order nor a Fermi surface; a spin-polarized calculation returns zero moment, and $\alpha^2F$ vanishes identically for an undoped insulator at $T = 0$ (\cref{sec:6.4}). Both statements are themselves computed results, and the worked cases here are NiO and aluminum.

\subsection{Place on the hierarchy}

Both calculations map a first-principles electronic structure onto a model whose parameters are then used elsewhere in the book. For magnets the model is the Heisenberg Hamiltonian, which assumes a rigid local moment on each magnetic ion and keeps only the energy of their relative orientation. Its parameters are static: they come from total energies or from static response, with the electrons treated in a mean field corrected on the $d$ shell. On the many-body axis $\mathcal{X}$ of \cref{sec:8.5}, DFT~$+\,U$ is a static, Hartree--Fock-like treatment of the localized shell, below the dynamical treatments that would describe the paramagnetic phase or the multiplet structure of the ion. The spin-wave theory that uses the constants is the lowest order in $1/|\vect{S}_j|$ (\cref{sec:7.1}). For superconductors, Eliashberg theory is the self-consistent electron self-energy at second order in the electron--phonon vertex, $n_g = 2$, summed to all orders in Nambu space, with vertex corrections dropped by Migdal's theorem (\cref{sec:6.3,sec:6.4}) and with the Coulomb repulsion reduced to one parameter. The weak-coupling BCS theory of \cref{sec:7.2} is its limit for $\lambda_{\mathrm{ep}} \ll 1$ and a constant interaction.

\subsection{DFT~$+\,U$ and the Hubbard parameter}

Semilocal functionals spread the $d$ electrons of transition-metal oxides over too many orbitals. In NiO they put nickel $d$ states at the top of the valence band and give a small gap of the wrong character, while DFT~$+\,U$ puts oxygen $p$ states there and gives the charge-transfer gap that experiments see~\cite[Sect.~V]{cococcioni2005linear}. The correction acts on the occupation matrices $n^{I\sigma}_{mm'} = \sum_{v\vect{k}}f_{v\vect{k}}\langle\psi^\sigma_{v\vect{k}}|P^I_{mm'}|\psi^\sigma_{v\vect{k}}\rangle$ of the localized orbitals of site $I$ with spin $\sigma$, where $P^I_{mm'}$ projects on the atomic orbitals $m$ and $m'$ of the shell~\cite[Eqs.~(5) and~(8)]{cococcioni2005linear}. In its simplest rotationally invariant form it adds
\begin{equation}
  E_U = \frac{U_{\mathrm{H}}}{2}\sum_{I\sigma}\Tr\bigl[n^{I\sigma}(1 - n^{I\sigma})\bigr]
  \label{eq:E.9-U}
\end{equation}
to the energy functional~\cite[Eq.~(9)]{cococcioni2005linear}. In the basis that diagonalizes $n^{I\sigma}$ the trace is $\sum_i\lambda_i(1 - \lambda_i)$ over its eigenvalues $0 \le \lambda_i \le 1$, so the correction penalizes fractional occupations and drives each orbital toward full or empty. The physical content is a restoration of the piecewise linearity of the energy as a function of the occupation of the shell, which semilocal functionals violate.

The parameter $U_{\mathrm{H}}$ can be computed rather than fitted. In the linear-response method of Cococcioni and de Gironcoli, a small potential shift $\delta V^U_I$ is applied to the localized orbitals of one site in a supercell, and the response of the occupations $n_J$ is computed twice: once self-consistently, giving the response matrix $R^U_{IJ} = \partial n_I/\partial(\delta V^U_J)$, and once at fixed Kohn--Sham potential, giving the bare response $R^{U,0}_{IJ}$. The Hubbard parameter is the difference of their inverses,
\begin{equation}
  U_{\mathrm{H}} = \bigl[(R^{U,0})^{-1} - (R^U)^{-1}\bigr]_{II} ,
  \label{eq:E.9-Ulr}
\end{equation}
which removes the curvature of the energy that comes from rehybridization of the noninteracting bands~\cite[Eqs.~(15), (18) and~(19)]{cococcioni2005linear}. The convergence parameter is the supercell size, since the perturbation of one site must not interact with its images. Density-functional perturbation theory replaces the supercells by a sum over wavevectors in the primitive cell and reaches the same $U_{\mathrm{H}}$ with better numerical control: for nickel in NiO, the two routes agree to $5\times10^{-3}$\,eV, and the value converges from 7.90 to \SI{8.17}{eV} as the supercell grows from $2\times2\times2$ to $4\times4\times4$ or the equivalent $\vect{q}$ grid is refined~\cite[Table~I]{timrov2018hubbard}.

The value is not unique. It depends on the projectors that define $n^{I\sigma}$, on the functional and on the pseudopotential~\cite[Sect.~V]{timrov2018hubbard}, and on whether $U_{\mathrm{H}}$ is recomputed self-consistently from the DFT~$+\,U$ ground state. For the same nickel shell in NiO, published linear-response values include \SI{4.6}{eV} with nonorthogonal atomic projectors~\cite[Sect.~V]{cococcioni2005linear}, \SI{8.17}{eV} with Löwdin-orthogonalized projectors from a semilocal ground state~\cite[Table~I]{timrov2018hubbard}, and \SI{6.26}{eV} computed self-consistently~\cite[Sect.~II]{dossantos2026comparative}. A $U_{\mathrm{H}}$ is meaningful only together with its projectors, and a calculation that takes $U_{\mathrm{H}}$ from one code and uses it in another changes the physics.

\subsection{Exchange constants}

We write the Heisenberg model as
\begin{equation}
  H_{\mathrm{H}} = \sum_{\langle ij\rangle}J^{\mathrm{ex}}_{ij}\,\vect{S}_i\cdot\vect{S}_j ,
  \label{eq:E.9-heis}
\end{equation}
with each pair of ions counted once and $J^{\mathrm{ex}}_{ij} > 0$ antiferromagnetic. The literature uses every combination of sign, of single or double counting of pairs, and of spin vectors of length $|\vect{S}|$ or of unit length. The conversions between them are factors of 2 and of $|\vect{S}|^2$~\cite[Table~1]{he2021tb2j}, and a value of $J^{\mathrm{ex}}$ quoted without its Hamiltonian is not a number.

\paragraph{Energy mapping.} The most direct route computes the total energies of several collinear magnetic configurations in DFT~$+\,U$ and equates their differences to those of \cref{eq:E.9-heis} with classical spins. In the fcc lattice of nickel ions in NiO, with $J_1$ between the 12 nearest neighbors and $J_2$ between the 6 next-nearest neighbors, three configurations fix both. In the ferromagnet each ion has all its neighbors parallel. In the antiferromagnet of type I, with ferromagnetic $(001)$ planes alternating in sign, 4 of the 12 nearest neighbors are parallel and 8 antiparallel, and all 6 next-nearest neighbors, at $\pm a$ along the cube axes, lie in planes of the same sign. In the type-II structure of the ground state, with ferromagnetic $(111)$ planes alternating, half of the nearest neighbors are parallel and all next-nearest neighbors antiparallel. Counting each pair once, the energies per ion are
\begin{equation}
  E_{\mathrm{FM}} = (6J_1 + 3J_2)|\vect{S}|^2, \qquad E_{\mathrm{AFI}} = (-2J_1 + 3J_2)|\vect{S}|^2, \qquad E_{\mathrm{AFII}} = -3J_2|\vect{S}|^2 ,
  \label{eq:E.9-map}
\end{equation}
relative to a common reference, and therefore
\begin{equation}
  J_1 = \frac{E_{\mathrm{FM}} - E_{\mathrm{AFI}}}{8|\vect{S}|^2}, \qquad J_2 = \frac{E_{\mathrm{FM}} + 3E_{\mathrm{AFI}} - 4E_{\mathrm{AFII}}}{24|\vect{S}|^2} .
  \label{eq:E.9-J12}
\end{equation}
These are Ref.~\cite[Eqs.~(6) and~(7)]{fischer2009exchange} converted to \cref{eq:E.9-heis}: that reference uses unit vectors, counts each pair twice and puts a minus sign in front, so that its $J_i$ equal $-J_i/2$ of ours for $|\vect{S}| = 1$. The two measured values for NiO that it collects, $J_2 = -\SI{8.66}{meV}$ (Shanker and Singh) and $-\SI{9.51}{meV}$ (Hutchings and Samuelsen~\cite{hutchings1972measurement}) in its convention, are $J_2 = 17.3$ and \SI{19.0}{meV} in ours.

More configurations resolve more neighbors. A cleaner variant flips one chosen pair of spins in a large supercell, in the four combinations $\uparrow\uparrow$, $\uparrow\downarrow$, $\downarrow\uparrow$ and $\downarrow\downarrow$, and isolates the coupling of that pair from all others; Ref.~\cite[Sect.~II]{dossantos2026comparative} uses it for NiO in $4\times4\times4$ supercells. The mapping has three weaknesses. The magnetic moment computed in DFT is not $2|\vect{S}|$ Bohr magnetons, and the classical $|\vect{S}|^2$ in \cref{eq:E.9-map} must be replaced by a choice, $|\vect{S}|^2$ or $|\vect{S}|(|\vect{S}| + 1)$, that changes the constants by up to a factor of two for $|\vect{S}| = 1$. The moments may change size from one configuration to the next, which a Heisenberg model cannot describe. And the result inherits the full dependence on $U_{\mathrm{H}}$: in NiO the superexchange scales roughly as $t_{\mathrm{hop}}^2/U_{\mathrm{H}}$ (\cref{sec:7.1}), so the spread of $U_{\mathrm{H}}$ above becomes a spread of $J_2$, and the density-functional values of Ref.~\cite[Table~III]{fischer2009exchange} range from 11 to \SI{27}{meV} in our convention (the unrestricted Hartree--Fock value in the same table is \SI{4.6}{meV}).

\paragraph{Spin spirals.} A spin spiral of wavevector $\vect{Q}$, with every moment tilted by a cone angle $\theta$ from a fixed axis and rotated in phase by $\vect{Q}\cdot\vect{R}$ from cell to cell, can be computed in the primitive cell when spin--orbit coupling is neglected, by a generalized Bloch theorem. For a lattice with one magnetic ion per cell, \cref{eq:E.9-heis} has $\vect{S}_{\vect{R}}\cdot\vect{S}_0 = |\vect{S}|^2[\sin^2\theta\cos(\vect{Q}\cdot\vect{R}) + \cos^2\theta]$ and therefore
\begin{equation}
  E(\vect{Q},\theta) - E(0,\theta) = \frac{|\vect{S}|^2\sin^2\theta}{2}\bigl[J^{\mathrm{ex}}(\vect{Q}) - J^{\mathrm{ex}}(0)\bigr], \qquad J^{\mathrm{ex}}(\vect{Q}) = \sum_{\vect{R}\neq0}J^{\mathrm{ex}}_{0\vect{R}}\,\ee^{\ii\vect{Q}\cdot\vect{R}} ,
  \label{eq:E.9-spiral}
\end{equation}
per ion, the single-sublattice case of Ref.~\cite[Eq.~(7)]{jacobsson2013exchange}. A set of spirals on a grid of $\vect{Q}$ maps out $J^{\mathrm{ex}}(\vect{Q})$, and a Fourier transform gives the real-space constants out to a range set by the grid. Small cone angles keep the perturbation of the ground state small, which reduces the change of the moments that plagues collinear mapping~\cite[Sect.~II]{jacobsson2013exchange}. For a ferromagnet, linear spin-wave theory in the convention of \cref{eq:E.9-heis} has $\omega(\vect{Q}) = |\vect{S}|[J^{\mathrm{ex}}(\vect{Q}) - J^{\mathrm{ex}}(0)]$, so the spiral energy is the magnon energy directly: $\omega(\vect{Q}) = 2[E(\vect{Q},\theta) - E(0,\theta)]/(|\vect{S}|\sin^2\theta)$, the frozen-magnon relation. For a collinear antiferromagnet with two sublattices, the same theory has
\begin{equation}
  \omega(\vect{Q}) = |\vect{S}|\Bigl\{\bigl[J^{\mathrm{ex}}_{\mathrm{AB}}(0) - J^{\mathrm{ex}}_{\mathrm{AA}}(0) + J^{\mathrm{ex}}_{\mathrm{AA}}(\vect{Q})\bigr]^2 - |J^{\mathrm{ex}}_{\mathrm{AB}}(\vect{Q})|^2\Bigr\}^{1/2} ,
  \label{eq:E.9-afm}
\end{equation}
with $J^{\mathrm{ex}}_{\mathrm{AA}}(\vect{Q})$ and $J^{\mathrm{ex}}_{\mathrm{AB}}(\vect{Q})$ the sums of $J^{\mathrm{ex}}_{0\vect{R}}\,\ee^{\ii\vect{Q}\cdot\vect{R}}$ over the sites $\vect{R} \ne 0$ of the same and of the other sublattice; with $J^{\mathrm{ex}}_{\mathrm{AB}} = 0$ it reduces to the ferromagnetic form. The spirals then need one cone angle per sublattice~\cite[Eq.~(8)]{jacobsson2013exchange}, and the two sums follow from their energies.

\paragraph{Green's functions.} The third route needs only the ground state. Rotating two moments by small angles changes the energy, by the magnetic force theorem, by the change of the single-particle energies at fixed potential, which second-order perturbation theory expresses through the Green's function of the ground state. The result is the formula of Liechtenstein, Katsnelson, Antropov and Gubanov, which in a Wannier basis reads
\begin{equation}
  J^{\mathrm{LKAG}}_{ij} = -\frac{1}{4\pi}\int_{-\infty}^{\mu_{\mathrm{ch}}}\dd\varepsilon\,\Imag\Tr\bigl[\Delta_i\,\GF^{\uparrow}_{ij}(\varepsilon)\,\Delta_j\,\GF^{\downarrow}_{ji}(\varepsilon)\bigr]
  \label{eq:E.9-lkag}
\end{equation}
in the convention of its source~\cite[Eq.~(19) and Table~1]{he2021tb2j}, $H = -\sum_{i\neq j}J^{\mathrm{LKAG}}_{ij}\,\hat{\vect{e}}_i\cdot\hat{\vect{e}}_j$ with unit vectors and each pair counted twice, so that $J^{\mathrm{ex}}_{ij} = -2J^{\mathrm{LKAG}}_{ij}/|\vect{S}|^2$ in the convention of \cref{eq:E.9-heis}. Here $\Delta_i$ is the exchange splitting of the on-site Hamiltonian of ion $i$, $\GF^\sigma_{ij}$ the block of the spin-$\sigma$ Green's function between the Wannier orbitals of ions $i$ and $j$, the trace runs over orbitals, and the integral runs up to the chemical potential $\mu_{\mathrm{ch}}$ (\cref{sec:2.5}). One ground-state calculation yields every pair at once, ferromagnetic or antiferromagnetic, and decomposes each constant by orbital. The formula describes infinitesimal rotations about the ground state, while energy mapping measures large ones, and the two differ where the moments are not rigid. For NiO, linear spin-wave theory with constants from energy mapping overestimates the measured magnon energies by about 14 percent, and with constants from the Green's-function formula by up to 40 percent~\cite[Sect.~III]{dossantos2026comparative}. The convergence parameters of \cref{eq:E.9-lkag} are the Wannier basis, which must contain the magnetic $d$ orbitals and the ligand $p$ orbitals that mediate superexchange, and the $\vect{k}$ grid of the Green's function, which sets the range of $|\vect{R}_i - \vect{R}_j|$ that is free of images.

\paragraph{Spin response.} The most complete route skips the Heisenberg model. Time-dependent density-functional perturbation theory computes the transverse spin susceptibility of the DFT~$+\,U$ ground state as a function of $(\vect{Q},\omega)$, and its poles are the magnons; constants are then fitted to the dispersion only for comparison~\cite[Sect.~II]{dossantos2026comparative}. This is the spin analogue of the dielectric response of \cref{app:E.6}, and it costs correspondingly more.

\subsection{Eliashberg theory}

The gap and critical temperature of a phonon-mediated superconductor follow from the Eliashberg equations, the Nambu--Gor'kov form of the electron self-energy with the phonon propagator of \cref{sec:6.4} and the Coulomb repulsion. On the imaginary axis, at the fermionic Matsubara frequencies $\omega^{\mathrm{F}}_j$ of \cref{sec:4.7}, and for an isotropic gap restricted to the Fermi surface, they are two coupled equations for a mass-renormalization function $Z^{\mathrm{El}}(\ii\omega^{\mathrm{F}}_j)$ and a gap function $\Delta^{\mathrm{El}}(\ii\omega^{\mathrm{F}}_j)$,
\begin{align}
  Z^{\mathrm{El}}(\ii\omega^{\mathrm{F}}_j) &= 1 + \frac{\pi T}{\omega^{\mathrm{F}}_j}\sum_{j'}\frac{\omega^{\mathrm{F}}_{j'}}{\sqrt{(\omega^{\mathrm{F}}_{j'})^2 + \Delta^{\mathrm{El}}(\ii\omega^{\mathrm{F}}_{j'})^2}}\;\lambda_{\mathrm{ep}}(\omega^{\mathrm{F}}_j - \omega^{\mathrm{F}}_{j'}) , \label{eq:E.9-Z}\\
  Z^{\mathrm{El}}(\ii\omega^{\mathrm{F}}_j)\,\Delta^{\mathrm{El}}(\ii\omega^{\mathrm{F}}_j) &= \pi T\sum_{j'}\frac{\Delta^{\mathrm{El}}(\ii\omega^{\mathrm{F}}_{j'})}{\sqrt{(\omega^{\mathrm{F}}_{j'})^2 + \Delta^{\mathrm{El}}(\ii\omega^{\mathrm{F}}_{j'})^2}}\Bigl[\lambda_{\mathrm{ep}}(\omega^{\mathrm{F}}_j - \omega^{\mathrm{F}}_{j'}) - \mu^*_{\mathrm{C}}\Bigr] , \label{eq:E.9-Delta}
\end{align}
with the frequency-dependent coupling
\begin{equation}
  \lambda_{\mathrm{ep}}(\omega^{\mathrm{B}}) = \int_0^\infty\dd\omega\,\frac{2\omega\,\alpha^2F(\omega)}{\omega^2 + (\omega^{\mathrm{B}})^2} ,
  \label{eq:E.9-lambda}
\end{equation}
evaluated at the bosonic differences $\omega^{\mathrm{B}} = \omega^{\mathrm{F}}_j - \omega^{\mathrm{F}}_{j'}$, so that $\lambda_{\mathrm{ep}}(0)$ is the coupling constant $\lambda_{\mathrm{ep}}$ of \cref{sec:6.4}. These are the isotropic forms of the Fermi-surface-restricted equations of Ref.~\cite[Eqs.~(41), (42) and~(48)]{lee2023electron}; the literature writes $Z$, $\Delta$, $\lambda$ and $\mu^*$, symbols that this book uses for other things. The highest temperature at which \cref{eq:E.9-Delta} has a nonzero solution is $T_c$, and the gap on the real axis follows by analytic continuation of $\Delta^{\mathrm{El}}(\ii\omega^{\mathrm{F}}_j)$, usually with Padé approximants (\cref{app:C.5}).

The Coulomb pseudopotential $\mu^*_{\mathrm{C}}$ is the weak point. It replaces the screened Coulomb repulsion between paired electrons, renormalized down from its bare Fermi-surface value by the retardation of the phonon interaction, and in most calculations it is not computed but chosen, with values from 0.1 to 0.2 giving reasonable agreement with experiment~\cite[Sect.~4]{lee2023electron}. Since $T_c$ depends on it exponentially, a computed $T_c$ with a chosen $\mu^*_{\mathrm{C}}$ is a consistency check on $\alpha^2F$ rather than a prediction. Two alternatives remove the parameter: computing $\mu^*_{\mathrm{C}}$ as a Fermi-surface average of the screened interaction, and density-functional theory for superconductors~\cite[Sect.~4]{lee2023electron}.

A quick estimate replaces the equations by the McMillan formula in the form of Allen and Dynes,
\begin{equation}
  T_c = \frac{\omega_{\log}}{1.2}\exp\left[-\frac{1.04(1 + \lambda_{\mathrm{ep}})}{\lambda_{\mathrm{ep}} - \mu^*_{\mathrm{C}}(1 + 0.62\lambda_{\mathrm{ep}})}\right], \qquad
  \omega_{\log} = \exp\left[\frac{2}{\lambda_{\mathrm{ep}}}\int_0^\infty\dd\omega\,\frac{\alpha^2F(\omega)}{\omega}\ln\omega\right] ,
  \label{eq:E.9-AD}
\end{equation}
which is Ref.~\cite[Eq.~(30)]{lee2023electron}, a fit to numerical solutions of the Eliashberg equations useful for preliminary estimates.

\paragraph{Convergence and cost.} The input $\alpha^2F$ is a double Fermi-surface average of $|g|^2$, with two energy delta functions, and it converges only on fine grids of $\vect{k}$ and $\vect{Q}$, which in practice means Wannier interpolation of the vertices (\cref{app:E.5}). The anisotropic equations, with a gap that varies over the Fermi surface, need the vertices on commensurate fine grids, restricted to an energy window of a few times the largest phonon energy around the Fermi level~\cite[Sect.~4]{lee2023electron}. The Matsubara sum needs frequencies up to a cutoff, commonly ten times the largest phonon frequency, with convergence tests recommended~\cite[Sect.~4]{lee2023electron}. Since the spacing is $2\pi T$, the number of frequencies grows as $1/T$, and low-$T_c$ superconductors are expensive: for aluminum, with phonons up to \SI{40.5}{meV} ($\omega_{\mathrm{max}} \approx \SI{470}{K}$~\cite[Sect.~III]{savrasov1996electron}), a cutoff of ten times that holds about 750 positive frequencies at \SI{1}{K} and 7500 at \SI{0.1}{K}. Sparse sampling of the Matsubara axis reduces the count~\cite[Sect.~4]{lee2023electron}. The analytic continuation to the real axis is the last numerical hazard: Padé approximants can produce spurious structure near the gap edge, which is where the pair-breaking absorption of \cref{chap:15} begins, and a real-axis solution or a check against the imaginary-axis data is needed there.

\begin{workedexample}{Aluminum from its Eliashberg function}\label{we:E-al}
Linear-response calculations of aluminum give $\lambda_{\mathrm{ep}} = 0.44$, against 0.42 from tunneling, and a logarithmic phonon frequency $\omega_{\log} = \SI{270}{K}$~\cite[Tables~II and~III]{savrasov1996electron}. Solving the Eliashberg equations with this $\alpha^2F$, the measured $T_c = \SI{1.18}{K}$ is reproduced with $\mu^*_{\mathrm{C}} = 0.12$, and the zero-temperature gap is then \SI{0.18}{meV}, equal to the measured value~\cite[Table~III]{savrasov1996electron}. The weak-coupling ratio of \cref{sec:7.2} agrees, $\Delta_{\mathrm{BCS}} = 1.764\,T_c = \SI{0.179}{meV}$, as it should for $\lambda_{\mathrm{ep}} = 0.44$.

\Cref{eq:E.9-AD} with the same three numbers gives $T_c = (270/1.2)\,\ee^{-5.213}\,\si{K} = \SI{1.22}{K}$, within four percent of the full solution. Its sensitivity to $\mu^*_{\mathrm{C}}$ is the logarithmic derivative of the exponent, $\partial\ln T_c/\partial\mu^*_{\mathrm{C}} = -1.04(1 + \lambda_{\mathrm{ep}})(1 + 0.62\lambda_{\mathrm{ep}})/[\lambda_{\mathrm{ep}} - \mu^*_{\mathrm{C}}(1 + 0.62\lambda_{\mathrm{ep}})]^2 = -23$. A change of $\mu^*_{\mathrm{C}}$ by $\pm0.02$, well inside the range used in the literature, moves $T_c$ by a factor of 0.60 or 1.53, from \SI{0.74}{K} at $\mu^*_{\mathrm{C}} = 0.14$ to \SI{1.87}{K} at 0.10. The gap and $T_c$ of aluminum are measured, and the first-principles $\alpha^2F$ enters where a measured quantity is not available, in the frequency dependence of the phonon-assisted pair-breaking absorption of \cref{chap:15}.
\end{workedexample}

Exchange constants and Eliashberg gaps carry uncertainties of different kinds. Exchange constants depend on the treatment of the localized shell, and their uncertainty is a spread over choices of $U_{\mathrm{H}}$ and of mapping; Eliashberg gaps depend on the electron--phonon vertices of \cref{app:E.5}, which are well controlled, and on one phenomenological parameter. Wherever a magnetic or superconducting target is ranked by a computed parameter, the corresponding spread is the theory uncertainty to carry (\cref{sec:8.8}).

\section{Defects}\label{app:E.10}

A carbon pair with a bound hydrogen, the T center of silicon, emits its zero-phonon line at \SI{935}{meV}; a hybrid-functional calculation in a 512-atom supercell puts it at \SI{985}{meV}~\cite{dhaliah2022first}. Displacing a silicon atom from its site takes between 17.5 and \SI{63.5}{eV}, depending on the direction in which it is struck~\cite{kadribasic2018directional}. Both numbers are defect properties that enter the book: the first in the color centers and dopants that serve as targets and readouts (\cref{sec:7.4,sec:16.6}), the second in the energy that a nuclear recoil stores in the lattice (\cref{chap:19}). Both come from supercells whose defect must be isolated from its own images, electrostatically for charged states and dynamically for a recoil; the effective-mass theory of shallow donors, which no supercell reaches, is in \cref{sec:7.4,app:D.23}.

\subsection{Place on the hierarchy}

A defect calculation is a calculation for a crystal whose translational symmetry is broken, and periodic codes restore it by repeating the defect in a supercell of side $L_{\mathrm{sc}}$. Each defect configuration is then an ordinary clamped-nucleus calculation at the base point $\mathfrak{T}_0$ of \cref{sec:8.5}, with one control that the perfect crystal does not need: the result must converge in $L_{\mathrm{sc}}$, and for charged defects it converges only after an electrostatic correction. The ground state of each charge state is a DFT ground state; excited states of color centers are obtained from constrained occupations, which go beyond ground-state DFT without reaching a correlated treatment of the many-body approximation $\mathcal{X}$. The vibronic side of an optical transition is computed in the harmonic approximation about each minimum, as the displaced-oscillator problem of \cref{sec:5.6}, and summed to all orders in the Huang--Rhys factor, which is the multiphonon end of the lattice regime $\mathcal{K}$. Displacement thresholds, finally, come from molecular dynamics with classical nuclei on the Born--Oppenheimer surface, $n_B = 0$, with the electronic excitations that a fast recoil produces left out.

\subsection{Formation energies and charge states}

The formation energy of a defect $D$ in charge state $z_{\mathrm{D}}$, in units of $e$, compares two supercells, with and without the defect,
\begin{equation}
  E^{\mathrm{f}}_{\mathrm{D}}(z_{\mathrm{D}}) = E_{\mathrm{D}}(z_{\mathrm{D}}) - E_{\mathrm{host}} - \sum_s n_s\,\mu^{\mathrm{at}}_s + z_{\mathrm{D}}\bigl(\varepsilon_{\mathrm{VBM}} + \mu_{\mathrm{ch}}\bigr) + E_{\mathrm{corr}} ,
  \label{eq:E.10-form}
\end{equation}
where $E_{\mathrm{D}}$ and $E_{\mathrm{host}}$ are the total energies of the defective and the host supercell, $n_s$ atoms of species $s$ are added ($n_s > 0$) or removed ($n_s < 0$) and exchanged with a reservoir at chemical potential $\mu^{\mathrm{at}}_s$, the electron chemical potential $\mu_{\mathrm{ch}}$ is measured from the valence maximum $\varepsilon_{\mathrm{VBM}}$, and $E_{\mathrm{corr}}$ is the finite-size correction of the next subsection~\cite[Eq.~(3)]{kim2020quickstart}. The atomic chemical potentials are bounded by the stability of the host against its elements and against competing phases, and $\mu_{\mathrm{ch}}$ ranges over the gap~\cite[Sect.~3]{kim2020quickstart}. Where two charge states have equal formation energy, $E^{\mathrm{f}}_{\mathrm{D}}(z_{\mathrm{D}}) = E^{\mathrm{f}}_{\mathrm{D}}(z'_{\mathrm{D}})$, \cref{eq:E.10-form} gives the thermodynamic transition level
\begin{equation}
  \mu_{\mathrm{ch}}(z_{\mathrm{D}}/z'_{\mathrm{D}}) = \frac{E^{\mathrm{f}}_{\mathrm{D}}(z_{\mathrm{D}})\bigr|_{\mu_{\mathrm{ch}}=0} - E^{\mathrm{f}}_{\mathrm{D}}(z'_{\mathrm{D}})\bigr|_{\mu_{\mathrm{ch}}=0}}{z'_{\mathrm{D}} - z_{\mathrm{D}}} ,
  \label{eq:E.10-level}
\end{equation}
the electron chemical potential at which the defect changes charge with full relaxation of the lattice. It is the quantity that photoluminescence quenching and capture experiments measure, and it differs from a Kohn--Sham eigenvalue, which carries no relaxation and no correction.

The defect is relaxed in a supercell as close to cubic as possible, built from the conventional cell, at the volume of the host~\cite[Sect.~3]{kim2020quickstart}. The work lies in the number of charge and spin states, which must all be tried, since a defect in a nonmagnetic host may have a spin-polarized ground state. The gap problem of semilocal functionals is more serious for defects than for bulk bands, because defect levels sit inside the gap and their position relative to both band edges matters; hybrid functionals, which give silicon a gap of \SI{1.114}{eV}~\cite{dhaliah2022first}, are the standard for levels and optical energies.

\subsection{Charged defects in a periodic cell}

A charged defect in a supercell is a periodic array of charges, neutralized by a uniform background so that the electrostatic energy is finite. Both the interaction with the images and the background are artifacts. The leading term is the energy of a simple-cubic lattice of point charges $z_{\mathrm{D}}e$ in a compensating background and a dielectric medium. In Heaviside--Lorentz units, with $e^2 = 4\pi\alpha$, the Coulomb energy of two charges $e$ at distance $r$ in vacuum is $\alpha/r$, and the Madelung energy per charge of the array is $-z_{\mathrm{D}}^2\alpha\,\alpha_{\mathrm{M}}/(2\varepsilon L_{\mathrm{sc}})$, with $\alpha_{\mathrm{M}} = 2.8373$ the Madelung constant of the simple-cubic lattice~\cite[Sect.~II]{komsa2009comparison}. The correction removes it:
\begin{equation}
  E_{\mathrm{corr}} \simeq \frac{z_{\mathrm{D}}^2\,\alpha\,\alpha_{\mathrm{M}}}{2\varepsilon\,L_{\mathrm{sc}}} + O(L_{\mathrm{sc}}^{-3}) ,
  \label{eq:E.10-mp}
\end{equation}
where the next term, of order $z_{\mathrm{D}}L_{\mathrm{sc}}^{-3}$ times the quadrupole moment of the defect charge, comes from the shape of that charge; this is the Makov--Payne form, quoted in Gaussian units in Ref.~\cite[Eq.~(8)]{kim2020quickstart}. The dielectric constant is $\varepsilon_0$ for a relaxed defect, whose screening includes the ions, and $\varepsilon_\infty$ for a vertical transition at frozen nuclei; in silicon the two are equal, 11.7. The bare point-charge term often overcorrects: for the doubly charged carbon vacancy in diamond it does so clearly~\cite[Sect.~III\,B]{komsa2009comparison}, because the defect charge is not a point. The quadrupole term is hard to determine, and a second artifact remains: the average electrostatic potential of a charged cell is set by convention, so that band edges and defect levels of different cells must be aligned. The schemes in current use, by Freysoldt, Neugebauer and Van de Walle and by Kumagai and Oba, model the defect charge as a Gaussian, compute its periodic energy analytically, and align the potential far from the defect, read from the plane-averaged or the atomic-site potentials, with the model potential~\cite[Sect.~II and Eq.~(1)]{kumagai2014electrostatics}. The Kumagai--Oba version also handles an anisotropic dielectric tensor, and shows that the conventional potential-alignment term is already contained in the image-charge correction~\cite[Sect.~IV]{kumagai2014electrostatics}. Without a correction, formation energies of charged defects converge as $1/L_{\mathrm{sc}}$, and their spread between correction schemes translates into defect densities that differ by orders of magnitude~\cite[Sect.~3]{kim2020quickstart}.

\subsection{Color centers}

The optical transition of a color center connects two states of a few localized electrons. Its zero-phonon energy $E_{\mathrm{ZPL}}$ is the difference of the total energies of the relaxed excited and ground states. When both states are well described by single Slater determinants, the excited state is computed by constraining the occupations, promoting an electron from one Kohn--Sham orbital to another and relaxing the structure in that configuration; this is the $\Delta$SCF method. For the negatively charged nitrogen--vacancy center of diamond, with its triplet transition between single-determinant states, a 512-site supercell gives $E_{\mathrm{ZPL}} = \SI{1.757}{eV}$ with a semilocal functional and \SI{2.035}{eV} with the HSE hybrid functional, against the measured \SI{1.945}{eV}~\cite[Sect.~4.1]{alkauskas2014first}. For the T center of silicon, the same approach in a 512-atom supercell with HSE and $\Gamma$-point sampling gives \SI{985}{meV} against \SI{935}{meV}, and the $(0/{-1})$ transition level at \SI{1.07}{eV} above the valence maximum~\cite{dhaliah2022first}. The radiative lifetime, from the transition dipole of the two Kohn--Sham orbitals and the Wigner--Weisskopf rate, comes out as 1.3 to \SI{5}{\micro s}, against a measured \SI{0.9}{\micro s}~\cite{dhaliah2022first}.

The same two relaxed geometries give the vibronic coupling. The Huang--Rhys factor $S_{\mathrm{HR}}$ of \cref{sec:5.6} is the number of phonons emitted on average when the lattice relaxes after the transition, $S_{\mathrm{HR}} = \sum_k\omega_k\Delta Q_k^2/2$, with $\Delta Q_k$ the projection of the mass-weighted displacement between the two minima on the vibrational mode $k$ of frequency $\omega_k$; the zero-phonon fraction of the emission is $\ee^{-S_{\mathrm{HR}}}$. Its convergence with supercell size is slow, because the relaxation reaches into long-wavelength acoustic modes that a small cell does not contain. For the nitrogen--vacancy center, a 64-site supercell gives $S_{\mathrm{HR}} = 3.02$, 20 percent below the converged 3.67, and the converged value needs an embedding scheme: the forces of the transition from a 512-site hybrid-functional calculation, and the vibrational modes of supercells of up to 10\,648 sites from force constants computed with a semilocal functional~\cite[Sects.~2.3 and~4.2]{alkauskas2014first}. The measured value is $3.85 \pm 0.05$~\cite[Sect.~3]{alkauskas2014first}.

$\Delta$SCF has a clear limit. States that are not single determinants, such as the singlet states of the nitrogen--vacancy center, need a correlated treatment of an active space of defect orbitals, embedded in the screened host, and the density transition matrix elements that a dark-matter excitation needs, $|\langle F|n_e(\vect{q})|I\rangle|^2$ over keV momenta, need the many-body states themselves rather than their energies (\cref{sec:16.6}). Shallow donors are the opposite limit: their effective-mass envelope extends over several nanometers (\cref{sec:7.4}), far beyond any supercell, and a supercell calculation can only supply the central-cell correction to the effective-mass theory of \cref{app:D.23}.

\subsection{Displacement thresholds}

The displacement threshold $E_{\mathrm{d}}$ of \cref{sec:7.4} is the smallest kinetic energy that, given to one atom in a given direction, leaves a stable defect. It is computed by molecular dynamics over many directions: for silicon it ranges from 17.5 to \SI{63.5}{eV} with the Stillinger--Weber potential, which of three common silicon potentials reproduces the density-functional thresholds best~\cite{kadribasic2018directional}. \Cref{app:E.11} gives the protocol, its cost and its approximations, among them the electronic excitation that Born--Oppenheimer dynamics leaves out (\cref{chap:19}).

\subsection{Tunneling defects in metals}

A defect that tunnels between two configurations $L$ and $R$, such as a light interstitial or an atom in a disordered film, couples to the conduction electrons of a metal through the change of their scattering potential. That coupling damps the tunneling and is described by the ohmic constant $\alpha_{\mathrm{ohm}}$ of \cref{sec:16.8}. In the Born limit, $\alpha_{\mathrm{ohm}} = 4K_{\mathrm{b}}$ with $K_{\mathrm{b}} = (\mathcal{N}_{\mathrm{F}}V_{\mathrm{d}}\Omega)^2$, where $\mathcal{N}_{\mathrm{F}}$ is the density of states per spin and volume at the Fermi level and $V_{\mathrm{d}}$ half the difference of the contact potentials of the two configurations, which sit at $\pm V_{\mathrm{d}}$~\cite[Eq.~(4)]{maillet2023nanomechanical}. Beyond the Born limit the change of potential is described by the scattering matrices $S_L$ and $S_R$ of the conduction electrons at the Fermi energy in the two configurations, with $S = \ee^{2\ii\delta}$ for a channel of phase shift $\delta$, and by the eigenvalues $\ee^{2\ii\tilde\delta_\nu}$ of $S_RS_L^\dagger$, so that $\tilde\delta_\nu$ is the difference of the phase shifts in channel $\nu$. A channel whose phase shift changes by $\tilde\delta$ reduces the long-time overlap of the two Fermi seas as $t^{-(\tilde\delta/\pi)^2}$, the orthogonality exponent of Anderson~\cite{anderson1967infrared}, and an ohmic bath of coupling $\alpha_{\mathrm{ohm}}$ produces $t^{-2\alpha_{\mathrm{ohm}}}$. Equating the two over all channels and both spins,
\begin{equation}
  \alpha_{\mathrm{ohm}} = \frac{1}{2\pi^2}\sum_\sigma\sum_\nu\tilde\delta_\nu^{\,2} ,
  \label{eq:E.10-ohm}
\end{equation}
which in the Born limit, with one channel and phase shifts $\mp\pi\mathcal{N}_{\mathrm{F}}V_{\mathrm{d}}\Omega$ in the two configurations, so that $\tilde\delta = 2\pi\mathcal{N}_{\mathrm{F}}V_{\mathrm{d}}\Omega$, returns $\alpha_{\mathrm{ohm}} = 4(\mathcal{N}_{\mathrm{F}}V_{\mathrm{d}}\Omega)^2 = 4K_{\mathrm{b}}$. A first-principles calculation needs the two relaxed configurations, from a nudged-elastic-band search (\cref{app:E.11}), identical relaxation of the host around both, and the scattering phases at the Fermi energy, either from a Green's-function method or from the overlap of the two many-electron ground states of a supercell as a function of its size. The orthogonality exponent appears only with fixed occupations: $\vect{k}$ sampling and smearing of the Fermi surface blur the logarithm on which it rests and must be removed or extrapolated.

\begin{workedexample}{A charged defect in a silicon supercell}\label{we:E-sidefect}
We estimate the image-charge error for a singly charged defect, $z_{\mathrm{D}} = \pm1$, in cubic supercells of silicon built from the conventional cell, $L_{\mathrm{sc}} = na$ with $a = \SI{5.431}{\angstrom}$, and with $\alpha\hbar c = \SI{14.40}{eV.\angstrom}$ and $\varepsilon_0 = 11.7$. \Cref{eq:E.10-mp} gives
\begin{equation*}
  E_{\mathrm{corr}} = \frac{2.8373\times\SI{14.40}{eV.\angstrom}}{2\times11.7\times L_{\mathrm{sc}}}
  = \SI{0.161}{eV},\ \SI{0.107}{eV},\ \SI{0.080}{eV}
\end{equation*}
for 64, 216 and 512 atoms ($n = 2$, 3, 4). Even the 512-atom cell, the size used for the T center, carries a point-charge correction of \SI{80}{meV}, larger than the \SI{50}{meV} by which its computed zero-phonon line misses the measurement. Since the point-charge term can overcorrect (above), production calculations use the Gaussian-charge schemes; what the estimate shows is that the error does not become negligible at any affordable size. It falls only as $1/L_{\mathrm{sc}}$, that is, as $N_{\mathrm{at}}^{-1/3}$, so that halving it means eight times more atoms. The cost of a hybrid-functional calculation grows faster than the cube of the number of atoms, and a 4096-atom cell is out of reach where a 512-atom cell is routine; the correction, not the cell, is what converges the result. The neutral T center needs none of this, and its error budget is the functional; its $(0/{-1})$ level at \SI{1.07}{eV} does, since the charged state enters it~\cite{dhaliah2022first}.
\end{workedexample}

\Cref{app:E.11} extends these calculations to millions of atoms and to the statistics of nuclear-recoil cascades.

\section{Recoil molecular dynamics}\label{app:E.11}

A silicon atom struck with \SI{100}{eV} leaves its site at \SI{26}{km/s}, comes within about \SI{1}{\angstrom} of a neighbor, and within a few picoseconds has either fallen back into the lattice or left behind a vacancy and an interstitial that store several eV between them. Nothing in this sequence is a small perturbation of the harmonic crystal. The atoms leave the region where the force constants of \cref{sec:5.2} apply, bonds break and re-form, and the outcome is decided by the shape of the Born--Oppenheimer energy surface far from equilibrium. \Cref{sec:19.1} hands the recoil over to this dynamics, and \cref{sec:19.4,sec:19.6} need four of its products: the displacement threshold $E_{\mathrm{d}}$ as a function of recoil direction (\cref{sec:7.4}), the distribution of the stored energy $E_{\mathrm{st}}$ at a given recoil energy $E_R$, the configurations of the defects that survive, and the barriers along the paths by which they recombine or move. Each comes from many trajectories on an energy surface fitted to density-functional theory, with the close collisions patched in by a screened Coulomb repulsion, and from the statistics of those trajectories over recoil directions.

\subsection{Computed quantities}

Molecular dynamics integrates Newton's equations for the nuclei on a single energy surface,
\begin{equation}
  M_\kappa\,\ddot{\vect{R}}_I = -\frac{\partial E^{\mathrm{ad}}(\{\vect{R}\})}{\partial\vect{R}_I},
  \label{eq:E.11-newton}
\end{equation}
where $\vect{R}_I$ is the position of nucleus $I$ of species $\kappa$ and $E^{\mathrm{ad}}(\{\vect{R}\})$ is the electronic ground-state energy at fixed nuclear positions, the Born--Oppenheimer surface of \cref{sec:5.1}. A recoil simulation starts from a crystal equilibrated at the base temperature, gives one atom (the primary knock-on atom) the kinetic energy $E_R$ in a chosen direction $\hat{\vect{r}}_R$, integrates \cref{eq:E.11-newton} until the excess energy has spread through the cell, and records what changed. The quantities read off a single trajectory are whether a defect survived, which one, and the potential energy it stores. Repeated over directions and energies, the trajectories give the probability $p_{\mathrm{FP}}(E_R, \hat{\vect{r}}_R)$ that a Frenkel pair survives, the threshold surface $E_{\mathrm{d}}(\hat{\vect{r}}_R)$, and the distribution of $E_{\mathrm{st}}$ that \cref{sec:19.4} bounds with calibration data.

Everything therefore rests on $E^{\mathrm{ad}}$. Density-functional theory (\cref{app:E.1}) gives $E^{\mathrm{ad}}$ and its gradients at a cost that grows as the cube of the number of atoms. A plane-wave calculation with $N_{\mathrm{at}}$ atoms holds a number of occupied bands and a number of plane waves that are both proportional to $N_{\mathrm{at}}$, and orthogonalizing the bands costs their number squared times the basis size. A recoil cascade needs thousands of atoms, to keep the recoil away from its periodic images, and thousands of trajectories, to sample directions. Direct density-functional molecular dynamics can afford a few directions in a cell of a few hundred atoms, and serves as the reference; the production runs use an interatomic potential that reproduces $E^{\mathrm{ad}}$ at a cost linear in $N_{\mathrm{at}}$.

\subsection{Place on the hierarchy}

In the language of the truncation label $\mathfrak{T}$ of \cref{sec:8.5}, \cref{eq:E.11-newton} is the adiabatic limit, $n_B = 0$: the electrons stay in their ground state for every nuclear configuration and never take energy from the nuclei. It keeps the anharmonicity of $E^{\mathrm{ad}}$ to all orders, which the expansion of \cref{sec:5.7} cannot, and it sits beyond the impulse regime of the lattice label $\mathcal{K}$, since the struck atom does not merely recoil from its site but collides with its neighbors. Two approximations come with it. The nuclei are classical, so zero-point motion is absent and the lattice starts at rest apart from thermal motion. For recoils above \SI{10}{eV} this does not matter for the mean threshold, since the kinetic energies exceed the vibrational energies of a few tens of meV by orders of magnitude~\cite[Sect.~II\,A]{sassi2022energy}; it does set the width of the threshold step, \SI{0.7}{eV} at \SI{12.5}{eV} (\cref{sec:19.1-handover}). It does matter for the slow processes at millikelvin, where barriers are crossed by tunneling rather than by thermal activation; we return to this when we discuss defect paths. The second approximation is that energy given to the electrons is lost from the description. At \SI{1}{keV} a silicon recoil already transfers part of its energy to electronic excitations (\cref{sec:19.2}), and that channel needs the nonadiabatic treatment of \cref{sec:19.2}; molecular dynamics then supplies the nuclear trajectory on which the electronic losses are evaluated.

\subsection{Interatomic potentials}

An interatomic potential replaces $E^{\mathrm{ad}}$ by a sum of local terms,
\begin{equation}
  E^{\mathrm{ad}}(\{\vect{R}\}) \approx \sum_I E^{\mathrm{at}}_\kappa(\mathcal{N}_I),
  \label{eq:E.11-local}
\end{equation}
where $\mathcal{N}_I$ is the set of positions of the neighbors of atom $I$ within a cutoff radius $r_{\mathrm{c}}$, measured from $\vect{R}_I$. Locality is what makes the cost linear in $N_{\mathrm{at}}$. It holds for covalent and metallic solids because the density matrix decays with distance, and it fails for long-range electrostatics in ionic materials, which are then added as an explicit Coulomb term.

The classical forms write $E^{\mathrm{at}}_\kappa$ as a handful of terms with physically motivated shapes. The Stillinger--Weber potential for silicon is a pair term plus a three-body term that penalizes deviations from the tetrahedral angle~\cite{stillinger1985computer}; Tersoff-type potentials let the strength of a bond depend on its local coordination. With some ten parameters fitted to a few properties, they reproduce the crystal well and the configurations far from it less well. For the defect energies and planar faults of silicon, the analytical potentials differ from density-functional values by tens of percent for at least some of the defects~\cite[Sect.~III\,E]{bartok2018machine}.

A machine-learned potential keeps \cref{eq:E.11-local} but lets $E^{\mathrm{at}}_\kappa$ be a flexible function of the neighborhood, fitted to density-functional data. The neighborhood is first encoded in a descriptor that is invariant under rotations, reflections and permutations of like atoms. Examples are symmetry functions of the neighbor distances and angles feeding a neural network~\cite{behler2007generalized}, the smooth overlap of atomic positions feeding a kernel regression~\cite{bartok2010gaussian}, or a systematic expansion in body-ordered invariants, the atomic cluster expansion~\cite{drautz2019atomic}. Equivariant message-passing networks build the descriptor from features that rotate with the structure and are passed between neighbors over several layers~\cite{batzner2022e3,batatia2022mace}. The families differ in cost and in how much data they need; they share \cref{eq:E.11-local} and therefore share its failure modes.

The Gaussian approximation potential for silicon~\cite{bartok2018machine} is a well-documented example and the one we use as the silicon case. Its energy is a fixed pair term plus a kernel sum,
\begin{equation}
  E^{\mathrm{ad}} = \sum_{I<J}V^{(2)}(|\vect{R}_I - \vect{R}_J|) + \sum_I\sum_{s=1}^{M_{\mathrm{rep}}}\alpha_s\,K_{\mathrm{GAP}}(\mathcal{N}_I, \mathcal{N}_s),
  \label{eq:E.11-gap}
\end{equation}
with $K_{\mathrm{GAP}}$ a similarity kernel between the neighborhood of atom $I$ and $M_{\mathrm{rep}}$ representative neighborhoods $\mathcal{N}_s$ drawn from the training data, and coefficients $\alpha_s$ fitted by regularized least squares~\cite[Eq.~(1)]{bartok2018machine}. The pair term $V^{(2)}$ is purely repulsive and is fitted to the density-functional energy of an isolated pair of silicon atoms; it carries the steep exchange repulsion at short distance and leaves the bonding to the kernel term~\cite[Sect.~II\,A]{bartok2018machine}. The cutoff is \SI{5}{\angstrom}.

\subsection{Training sets}

The fit is only as good as the configurations it has seen. Each density-functional calculation in the training set yields a total energy, six components of the stress and three force components per atom, and all of them enter the fit with tolerances that set their relative weight. For silicon the database holds 2475 periodic cells with 171\,815 atomic environments~\cite[Table~I]{bartok2018machine}. They cover the diamond crystal at strains and temperatures up to melting, other crystal phases ($\beta$-tin, simple hexagonal, bc8, fcc, hcp and others), liquid and amorphous cells, low-index surfaces with their reconstructions and decohesion paths, the vacancy, the divacancy and the interstitials in cells of 63 to 217 atoms, crack tips and a screw dislocation. The reference is the PW91 functional, with a plane-wave cutoff of \SI{250}{eV}~\cite[Sect.~II\,B]{bartok2018machine}.

What a recoil needs from this list is the defect region and the short-distance region. The defects are there by design. The short distances are not, and thermal sampling cannot put them there. At \SI{1.0}{\angstrom} two silicon atoms repel each other with about \SI{51}{eV} (\cref{eq:E.11-zbl} below), while the hottest training configurations, liquid silicon near \SI{2000}{K}, have $T \approx \SI{0.17}{eV}$. A thermal ensemble visits such a pair separation with a Boltzmann weight of order $\ee^{-51/0.17}$, which is zero for every practical purpose. A training set assembled from equilibrium and thermal configurations therefore leaves the potential unconstrained exactly where a recoil cascade begins, and the cure is the splice described next.

\subsection{Close collisions and the Ziegler--Biersack--Littmark splice}

At separations well below a bond length the energy of two atoms is dominated by the Coulomb repulsion of the nuclei, screened by the electrons that remain between them. The standard description is the universal screened potential of Ziegler, Biersack and Littmark (ZBL)~\cite[Chap.~2]{ziegler1985stopping},
\begin{equation}
  V_{\mathrm{ZBL}}(r) = \frac{Z_1 Z_2\,\alpha}{r}\,\phi\!\left(\frac{r}{a_{\mathrm{U}}}\right), \qquad
  a_{\mathrm{U}} = \frac{0.8854\,a_0}{Z_1^{0.23} + Z_2^{0.23}},
  \label{eq:E.11-zbl}
\end{equation}
with $\phi(x) = 0.1818\,\ee^{-3.2x} + 0.5099\,\ee^{-0.9423x} + 0.2802\,\ee^{-0.4029x} + 0.02817\,\ee^{-0.2016x}$, $a_0$ the Bohr radius and $Z_{1,2}$ the nuclear charges. It is a fit to the repulsion of many atom pairs computed from their free-atom charge densities, which makes it universal in the sense that one screening function serves all pairs once lengths are measured in $a_{\mathrm{U}}$. For silicon on silicon, $a_{\mathrm{U}} = \SI{0.128}{\angstrom}$, and the unscreened Coulomb energy $Z^2\alpha/r = \SI{2822}{eV\,\angstrom}/r$ is reduced by $\phi$ to \SI{51}{eV} at \SI{1.0}{\angstrom} and \SI{9.6}{eV} at \SI{1.5}{\angstrom}.

In a head-on collision of the recoiling atom with a neighbor at rest, the two equal masses share the kinetic energy $E_R/2$ of the center-of-mass frame, and the turning point $r_{\mathrm{min}}$ solves $V_{\mathrm{ZBL}}(r_{\mathrm{min}}) = E_R/2$. Solved numerically with \cref{eq:E.11-zbl} for silicon, it is 1.49, 1.20, 1.00, 0.83, 0.49 and \SI{0.18}{\angstrom} at $E_R = 20$, 50, 100, 200, 1000 and \SI{10000}{eV} (\cref{tab:E-md-si}), against the bond length of \SI{2.35}{\angstrom}. Head-on collisions are the extreme case, and glancing ones stay farther apart. Near the threshold, $E_R \approx 20$ eV, the atoms come to about \SI{1.5}{\angstrom}, inside the bonding region that a good training set covers, while above \SI{100}{eV} they reach the core region that no thermal training set covers.

The two descriptions are joined by a smooth switch,
\begin{equation}
  E^{\mathrm{ad}} = \sum_{I<J}\bigl[1 - f_{\mathrm{sw}}(r_{IJ})\bigr]V_{\mathrm{ZBL}}(r_{IJ}) + f_{\mathrm{sw}}(r_{IJ})\,E^{\mathrm{ad}}_{\mathrm{ML}},
  \qquad f_{\mathrm{sw}}(r) = \frac{1}{1 + \ee^{-b_{\mathrm{sw}}(r - r_{\mathrm{sw}})}},
  \label{eq:E.11-splice}
\end{equation}
written here schematically: in practice the switch multiplies the pair part of the machine-learned energy and the many-body part is left to vanish at short distance by its training. The center $r_{\mathrm{sw}}$ and steepness $b_{\mathrm{sw}}$ are chosen so that the potential is pure $V_{\mathrm{ZBL}}$ below about \SI{1}{\angstrom} and pure machine-learned energy above about \SI{1.5}{\angstrom}. The design rule follows from the turning points. Threshold dynamics must not feel the splice, so the machine-learned energy must be fully restored above the $r_{\mathrm{min}}$ of a \SI{20}{eV} collision, \SI{1.49}{\angstrom}. It also must not leave a minimum, a shoulder or a jump in the force anywhere in between, which would deflect every collision that passes through it. The splice is checked against density-functional energies of a silicon pair, and of a crystal with one atom pushed toward a neighbor, at separations between 0.8 and \SI{1.8}{\angstrom}. A potential whose own pair term was already fitted to the density-functional dimer, as in \cref{eq:E.11-gap}, is usually close to $V_{\mathrm{ZBL}}$ near \SI{1}{\angstrom}, and the check then confirms rather than repairs.

\subsection{Integrating a recoil}

The equations of motion are integrated with the velocity Verlet scheme, which is time-reversible and conserves a shadow energy close to the true one, so that the total energy of an isolated cell drifts only through round-off and through forces that are not smooth. The time step is set by the fastest atom. A silicon atom with \SI{100}{eV} moves at $v = (2E_R/M_\kappa)^{1/2} = \SI{26.2}{km/s}$, which is \SI{0.26}{\angstrom} per femtosecond; at \SI{1}{keV} it moves \SI{0.83}{\angstrom/fs}. Resolving a collision whose turning point lies within a fraction of an \si{\angstrom} of the steep part of $V_{\mathrm{ZBL}}$ requires steps in which no atom moves more than a few hundredths of an \si{\angstrom}: \SI{0.05}{\angstrom} per step means \SI{0.19}{fs} at \SI{100}{eV} and \SI{0.06}{fs} at \SI{1}{keV}. After a few hundred femtoseconds the fastest atom carries a few eV, and the step can grow to the \SI{1}{fs} that thermal silicon allows, about a sixty-fifth of the shortest phonon period (\SI{65}{fs}, from the \SI{64.1}{meV} optical phonon). Recoil codes therefore adapt the step to the current maximum velocity and force, a practice the silicon and germanium threshold studies follow~\cite[Table~I]{sassi2022energy}. The check is energy conservation: with the thermostat switched off, the total energy of the cell must stay constant to well below the stored energies being measured, which for this application means a small fraction of an eV over the whole run.

The cell has to absorb the recoil. Its excess energy, $E_R$ minus whatever is stored, is spread over the cell as heat, and it must leave through a boundary layer held at the base temperature by a thermostat. The silicon and germanium studies use cells of $8\times8\times8$ conventional cubes (4096 atoms), a thermostatted border \SI{6}{\angstrom} thick, and runs of \SI{20}{ps}~\cite[Sect.~II\,A and Table~I]{sassi2022energy}. Three things set these numbers. The recoil must stop well inside the cell, so that it neither crosses the periodic boundary nor reaches the thermostat while still energetic; at \SI{100}{eV} the primary travels a few \si{\angstrom} to about a nanometer, and $8a = \SI{4.3}{nm}$ is ample. The heat must be removed without reflecting back from the thermostat as a coherent pulse, which a too thin or too stiff border does. And the run must last until the potential energy of the cell has settled; the same study fixed its run lengths from test runs longer than \SI{5}{ps}, by the time the potential energy took to come within \SI{1}{eV} of its final value~\cite[Sect.~II\,A]{sassi2022energy}.

\subsection{Cascade statistics}

A single trajectory answers a yes-or-no question for one direction and one energy. The physical quantities are averages over the direction of the recoil, which a dark-matter or neutron scattering event draws from its own angular distribution, and over the position of the struck atom within the unit cell. The standard protocol draws $N_{\mathrm{dir}}$ directions uniformly on the sphere and, for each, increases $E_R$ in steps of about \SI{1}{eV} through the threshold region and more coarsely above it. The published silicon and germanium scans use more than 1000 directions~\cite[Sect.~II\,A]{sassi2022energy}, and 1003 for silicon in the earlier stored-energy study~\cite{kadribasic2020crystal}.

Whether a defect survived is decided by an occupation count: each atom of the final configuration is assigned to the Wigner--Seitz cell of the nearest perfect-lattice site, and any site with occupation other than one marks a vacancy or an interstitial. The method has no free parameter and does not count the bond-switching defects that relax away~\cite[Sect.~II\,A]{sassi2022energy}. From the counts we obtain $p_{\mathrm{FP}}(E_R)$, the fraction of directions with a surviving defect, with the binomial error $[p_{\mathrm{FP}}(1 - p_{\mathrm{FP}})/N_{\mathrm{dir}}]^{1/2}$, at most \num{0.016} for 1000 directions. The threshold surface needs more care than the word ``threshold'' suggests. Along many directions the defect appears at one energy, disappears at a slightly higher energy because the interstitial now lands close enough to its vacancy to recombine, and reappears above that. Silicon and germanium show some of this recombination; sapphire shows so much that a direction-resolved threshold is ill-defined~\cite[Sect.~II\,B]{sassi2022energy}. The quantity the rate formulas of \cref{sec:19.4} need is therefore $p_{\mathrm{FP}}(E_R, \hat{\vect{r}}_R)$ itself, and $E_{\mathrm{d}}(\hat{\vect{r}}_R)$ is quoted as its lowest crossing. With the Stillinger--Weber potential the minimum threshold over all directions is \SI{18}{eV} in silicon and \SI{13}{eV} in germanium (with a modified Stillinger--Weber potential); above it the defect probability rises steeply, though more slowly than in diamond, because of recombination~\cite[Table~II, Fig.~1 and Sect.~II\,B]{sassi2022energy}. Both numbers are properties of the potentials; \cref{sec:19.4} discusses why the germanium value is model-conditional.

\subsection{Stored energy and the quench protocol}

The stored energy of a trajectory is the potential energy the cell keeps once all motion has died out. Taking the difference of potential energies at the end and at the start of a finite run, as the published scans do~\cite[Sect.~II\,A]{sassi2022energy}, includes whatever heat has not yet left the cell. The size of that error follows from equipartition: in the harmonic crystal the potential energy is $\tfrac32N_{\mathrm{at}}T$, so a residual temperature rise of \SI{1}{K} in 4096 atoms adds $\tfrac32\times4096\times\SI{8.6e-5}{eV} = \SI{0.53}{eV}$, comparable to the uncertainty of the stored energy itself. A quench removes it. The final configuration is relaxed to its local minimum by energy minimization: conjugate gradients, or damped dynamics with the velocities set to zero whenever the power $-\nabla E^{\mathrm{ad}}\cdot\dot{\vect{R}}$ turns negative. The perfect crystal is relaxed the same way, and
\begin{equation}
  E_{\mathrm{st}} = E^{\mathrm{ad}}_{\mathrm{min}}(\text{final}) - E^{\mathrm{ad}}_{\mathrm{min}}(\text{perfect}) .
  \label{eq:E.11-Est}
\end{equation}
The minimum reached is the inherent structure of the final configuration: thermal motion is removed and the defects are kept, as long as the barriers between them exceed the small forces of the minimizer. For the experiments of \cref{sec:19.4}, which measure heat at millikelvin long after the event, \cref{eq:E.11-Est} is the relevant quantity.

A second protocol choice concerns the reference energies. A cascade census produces vacancies, interstitials of several types, close Frenkel pairs and small clusters, each with a formation energy that the potential predicts with some error. The stored energy of a trajectory is roughly the sum of its defects' energies plus their interactions. To make results from different potentials, or from potential and density functional, comparable, one table of configurational energies of the recurring defects is computed once with density-functional theory in large relaxed cells, and every census is checked against it. For the silicon potential of \cref{eq:E.11-gap}, the formation energies of the vacancy and of the hexagonal, tetrahedral and dumbbell interstitials lie within \SI{7}{\percent} of the density-functional reference~\cite[Sect.~III\,E]{bartok2018machine}. A defect storing several eV then carries an error of a few tenths of an eV from the potential alone, which is the floor on how precisely a single-pair stored energy can be quoted.

\subsection{Defect paths and barriers}

Once defects survive, \cref{sec:19.6} needs to know how they move and recombine. The tool is the nudged elastic band: a chain of $N_{\mathrm{im}}$ replicas of the cell, between two relaxed end states, is relaxed with spring forces along the chain and the true force perpendicular to it, until the chain lies on the minimum-energy path. In the climbing-image variant the highest replica is driven uphill along the path and converges onto the saddle point, whose energy relative to the initial minimum is the barrier $V_{\mathrm{b}}$~\cite{henkelman2000climbing}. End states come from the cascade census, which at $E_R = \SI{100}{eV}$ already supplies close pairs at many separations and orientations. Paths are first mapped with the machine-learned potential, at a cost of a few thousand force calls each, and the saddle points that matter are recomputed with density-functional theory in cells of 64 to 216 atoms. The silicon potential reproduces the density-functional ordering of the hexagonal and tetrahedral interstitials, which are related by a short displacement along a diffusion path, though it underestimates their energy difference; the analytical potentials make larger errors, and many reverse the order~\cite[Sect.~III\,E]{bartok2018machine}.

A barrier becomes a rate through harmonic transition-state theory, $\varGamma = \varGamma_0\,\ee^{-V_{\mathrm{b}}/T}$, with an attempt rate $\varGamma_0$ from the vibrational frequencies at the minimum and at the saddle. At the operating temperature of a cryogenic detector this is useless for any barrier of chemical size: at \SI{10}{mK}, $T = \SI{0.86}{\micro eV}$, and a barrier of \SI{0.1}{eV} gives a Boltzmann factor of $\ee^{-1.2\times10^5}$. What moves defects at millikelvin is tunneling through low barriers, or activation by the athermal phonons of a later event. The first needs the barrier shape and the tunneling mass along the path, which the nudged elastic band provides and \cref{sec:19.6} turns into a rate; the second needs the energy that reaches the defect, which the transport of \cref{app:E.12} provides.

\subsection{Convergence parameters and cost}

The parameters a recoil study has to converge are the following. For the potential: the size and coverage of the training set, the cutoff $r_{\mathrm{c}}$, the descriptor resolution, and the splice $(r_{\mathrm{sw}}, b_{\mathrm{sw}})$. For the dynamics: the time-step criterion, the cell size, the border thickness and the run length. For the statistics: $N_{\mathrm{dir}}$, the energy grid, and the sampling of the struck atom over the inequivalent sites of the cell. For the stored energy: the quench tolerance. Each is converged by varying it and checking that $p_{\mathrm{FP}}$ and the mean $E_{\mathrm{st}}$ move by less than their statistical errors.

The cost is set by the number of force evaluations. The silicon Gaussian approximation potential costs about \SI{100}{ms} per atom per force call on one core~\cite[Sect.~V]{bartok2018machine}. A 4096-atom cell then needs \SI{410}{s} per step, and a trajectory of $2\times10^4$ steps ($\sim\SI{20}{ps}$ at a typical \SI{1}{fs} step after the collision phase) needs $8\times10^6$ core-seconds, about $2\times10^3$ core-hours. A scan of 1000 directions at 20 energies is $2\times10^4$ trajectories, or $5\times10^7$ core-hours. This is why the published scans use the cheaper analytical potentials, and why a production study with a machine-learned potential has to use one of the faster descriptor families, on graphics processors, and has to spend its trajectories where $p_{\mathrm{FP}}$ changes fastest. A density-functional reference trajectory in a 216-atom cell costs far more per atom, but only a few directions are needed to validate the potential.

\subsection{Pitfalls}

The most common failure is a potential that is accurate at equilibrium and wrong in the region a recoil explores. Two potentials can agree on elastic constants and defect energies and still predict thresholds that differ by several eV, because the threshold is controlled by the energy along the path of the struck atom between its neighbors, not at the end points. The remedy is validation on the quantity itself: thresholds along a few directions from density-functional molecular dynamics, the approach of the density-functional reference study for silicon~\cite{holmstrom2008threshold}. Two physical omissions of \cref{eq:E.11-newton} remain: electronic energy losses, which grow with $E_R$ and which \cref{sec:19.2} treats, and the quantum motion of light defects at millikelvin, which classical dynamics cannot describe.

\subsection{Silicon}

\Cref{tab:E-md-si} collects the silicon numbers of this section. The analytical potentials are used for the large scans, the Gaussian approximation potential for validation and for paths, and density-functional theory as the reference for both.

\begin{table}[t]
  \caption{Recoil molecular dynamics in silicon: the inputs and the quantities a production study converges.}
  \label{tab:E-md-si}
  \centering
  \begin{tabular}{p{0.40\linewidth}p{0.48\linewidth}}
    \toprule
    Quantity & Silicon \\
    \midrule
    Bond length; shortest phonon period & \SI{2.35}{\angstrom}; \SI{65}{fs} \\
    Head-on turning point $r_{\mathrm{min}}$ at $E_R = 20$, 50, 100, 200, 1000, \SI{10000}{eV} & 1.49, 1.20, 1.00, 0.83, 0.49, \SI{0.18}{\angstrom} (\cref{eq:E.11-zbl}) \\
    Speed of the struck atom at 100 and \SI{1000}{eV} & 0.26 and \SI{0.83}{\angstrom/fs} \\
    Training set of the machine-learned potential & 2475 cells, 171\,815 environments, PW91~\cite{bartok2018machine} \\
    Point-defect formation energies vs density-functional theory & within \SI{7}{\percent}~\cite{bartok2018machine} \\
    Cost per atom per force call & \SI{100}{ms} on one core~\cite{bartok2018machine} \\
    Cell, border, run length & 4096 atoms, \SI{6}{\angstrom}, \SI{20}{ps}~\cite{sassi2022energy} \\
    Directions sampled & $> 1000$~\cite{sassi2022energy} \\
    Minimum threshold (Stillinger--Weber) & \SI{18}{eV}~\cite{sassi2022energy} \\
    Residual heat per K in 4096 atoms & \SI{0.53}{eV} (equipartition) \\
    \bottomrule
  \end{tabular}
\end{table}

The machine-learned potential removes the main systematic of the published silicon scans, the analytical potential, but at the $5\times10^7$ core-hours estimated above for a full scan; that trade decides how far the stored-energy predictions of \cref{sec:19.4} can be pushed.

\section{Phonon and charge transport Monte Carlo}\label{app:E.12}

A \SI{100}{eV} deposit in a silicon crystal ends, after the cascades of \cref{chap:18,chap:19}, as some $10^4$ acoustic phonons of a few meV and, in a biased detector, a few tens of electrons and holes drifting toward the electrodes. The sensors on the surface see neither the deposit nor the cascade. They see phonons that arrive over microseconds after scattering off isotopes, decaying, bouncing off surfaces and crossing into metal films, and charges that arrive after drifting, emitting phonons of their own, and sometimes getting trapped. \Cref{sec:20.1,sec:20.2,sec:20.3,sec:20.4,sec:20.5} derive the physics of each step; combined in one simulation of a whole detector, they turn an event into the arrival distributions that \cref{sec:20.5} calls collection kernels. The method is Monte Carlo solution of the Boltzmann equation for phonons and carriers, and the silicon case is worked with the rate constants and code conventions of the public detector-physics toolkit G4CMP~\cite{kelsey2023g4cmp,leman2012invited}.

\subsection{Computed quantities}

The state of a phonon is its wavevector $\vect{Q}$, its branch $\nu$ and its position; the state of a carrier is its crystal momentum, its valley or band, and its position. Between events each particle moves on a straight line with its group velocity (a carrier also accelerates in the applied field). Events are drawn at random with the rates of the scattering processes, and each event replaces the particle by its products. With many particles, the procedure samples the solution of the linear Boltzmann equation
\begin{equation}
  \partial_t f_i + \vect{v}_i\cdot\nabla_{\vect{r}} f_i + \vect{F}_i\cdot\nabla_{\vect{k}} f_i = \Bigl(\frac{\partial f_i}{\partial t}\Bigr)_{\mathrm{coll}},
  \label{eq:E.12-boltzmann}
\end{equation}
for the occupation $f_i(\vect{r}, \vect{k}, t)$ of each species $i$, with $\vect{v}_i$ its group velocity, $\vect{F}_i$ the force on it ($\pm e\vect{\mathcal{E}}$ for carriers in a field $\vect{\mathcal{E}}$, zero for phonons) and a collision term built from golden-rule rates~\cite{jacoboni1983monte}. The equation is linear because the athermal populations are dilute: phonons scatter off the static disorder and decay spontaneously, carriers emit into an empty phonon bath, and no process depends on the occupation of the other particles. The outputs are tallies at the boundaries: for every sensor, the energy and time of each phonon it absorbs, and for every electrode, the charge induced on it.

\subsection{Place on the hierarchy}

Every rate in the collision term is a lowest-order golden-rule rate. Isotope scattering is the second-order Born rate of a phonon off the random mass defects of \cref{sec:11.6}. Anharmonic decay is the rate of the cubic term of the lattice potential (\cref{sec:5.7}), at the lowest order at which it does not vanish, the cubic vertex squared. Neganov--Trofimov--Luke emission and intervalley scattering are one-phonon emission by a carrier through a deformation potential, $n_g = 2$ in the notation of the truncation label $\mathfrak{T}$ of \cref{sec:8.5}, with $n_B = 0$. The transport itself is incoherent: each particle is a wavepacket with a well-defined momentum, and interference between successive events is dropped. This holds when the scattering rate is small compared with the particle's frequency. For a silicon phonon at \SI{1}{THz} the isotope rate is \SI{2.4e6}{s^{-1}} (\cref{eq:E.12-iso} below) against $2\pi\times\SI{e12}{s^{-1}}$; even at the top of the acoustic band, \SI{15}{THz}, the ratio is $\sim10^{-3}$. It also needs mean free paths long compared with the wavelength, which they are by a factor $4\times10^5$ at \SI{1}{THz} (a wavelength of \SI{5.4}{nm} against a mean free path of \SI{2.2}{mm}). The Boltzmann description is therefore safe for the bulk; its approximations are concentrated in the rates and in the surface models, which is where the pitfalls below lie.

\subsection{Phonon rates}

Isotope scattering and anharmonic decay remove a phonon from its ballistic path. Isotope scattering redirects it, and may change its branch, at the rate~\cite{tamura1983isotope}, \cite[Eq.~(6)]{leman2012invited}
\begin{equation}
  \varGamma_{\mathrm{iso}} = B_{\mathrm{iso}}\left(\frac{E_{\mathrm{ath}}}{h}\right)^4 , \qquad
  B_{\mathrm{iso}} = 4\pi^3\,g_{\mathrm{md}}\,\Omega_{\mathrm{at}}\,\langle c^{-3}\rangle,
  \label{eq:E.12-iso}
\end{equation}
where $E_{\mathrm{ath}}/h$ is the phonon frequency, $\Omega_{\mathrm{at}}$ the volume per atom, $\langle c^{-3}\rangle = \tfrac13\sum_\nu c_\nu^{-3}$ the average of the inverse cubed sound speed over the three acoustic branches $\nu$ (the source writes $1/\langle c^3\rangle$ with the polarization-averaged cubed speed; the two agree when the speeds are equal, and $\langle c^{-3}\rangle$ is the average the Debye density of states contains, \cref{sec:20.2}), and $g_{\mathrm{md}} = \langle (M - \bar M)^2\rangle/\bar M^2$ the mass variance of the isotope mixture over the squared mean mass~\cite[Eqs.~(3)--(5)]{hernandez2025modeling}. It is the Rayleigh law of a point scatterer, valid as long as the wavelength is long compared with the interatomic spacing. For natural silicon, $g_{\mathrm{md}} = \num{2.02e-4}$, $\Omega_{\mathrm{at}} = \Omega/2 = \SI{20.0}{\angstrom^3}$ and, with $c_{\mathrm{LA}} = \SI{9.0}{km/s}$ and $c_{\mathrm{T}} = \SI{5.4}{km/s}$, $\langle c^{-3}\rangle = 1/(\SI{2.13e11}{m^3.s^{-3}}) = \SI{4.69e-12}{s^3.m^{-3}}$, and \cref{eq:E.12-iso} gives $B_{\mathrm{iso}} = 124.0\times\num{2.02e-4}\times\SI{2.00e-29}{m^3}\times\SI{4.69e-12}{s^3.m^{-3}} = \SI{2.35e-42}{s^3}$, in line with the tabulated $\SI{2.33e-42}{s^3}$ and with the value $\SI{2.43e-42}{s^3}$ that most transport codes use~\cite[Table~4]{hernandez2025modeling}, \cite[Table~I]{leman2012invited}. The direction and branch of the scattered phonon are then drawn with the weight $|\hat{\vect{e}}\cdot\hat{\vect{e}}'|^2$ of the polarization vectors of the initial and final modes, divided by the cube of the final frequency~\cite[Eq.~(7)]{leman2012invited}. Repeated isotope scattering drives the branch populations of silicon to slow transverse \SI{55}{\percent}, fast transverse \SI{35}{\percent} and longitudinal \SI{10}{\percent}~\cite[Sect.~III\,E]{leman2012invited}.

Anharmonic decay splits a longitudinal phonon into two of lower energy, at the rate~\cite{tamura1985spontaneous}, \cite[Eq.~(8)]{leman2012invited}
\begin{equation}
  \varGamma_{\mathrm{anh}} = A_{\mathrm{anh}}\left(\frac{E_{\mathrm{ath}}}{h}\right)^5 ,
  \label{eq:E.12-anh}
\end{equation}
with $A_{\mathrm{anh}}$ fixed by the second- and third-order elastic constants in the isotropic continuum~\cite[Eq.~(8)]{leman2012invited}. Transverse phonons hardly decay at low frequency. The two channels, $\mathrm{L}\to\mathrm{L}+\mathrm{T}$ and $\mathrm{L}\to\mathrm{T}+\mathrm{T}$, have energy distributions of the daughters that the isotropic approximation fixes in closed form. For $\mathrm{L}\to\mathrm{L}'+\mathrm{T}'$, with $x = E_{\mathrm{L}'}/E_{\mathrm{L}}$ and $\delta = c_{\mathrm{LA}}/c_{\mathrm{T}}$, the rate is distributed as $x^{-2}(1 - x^2)^2[(1 + x)^2 - \delta^2(1 - x)^2][1 + x^2 - \delta^2(1 - x)^2]^2$ over $(\delta - 1)/(\delta + 1) < x < 1$, and the $\mathrm{T}+\mathrm{T}$ channel has a polynomial distribution of the same kind with coefficients set by the elastic constants~\cite[Eqs.~(9) and~(12)]{leman2012invited}. The daughter directions then follow from energy and momentum conservation; for the longitudinal daughter, $\cos\theta_{\mathrm{L}'} = [1 + x^2 - \delta^2(1 - x)^2]/(2x)$~\cite[Eq.~(10)]{leman2012invited}. For silicon, Tamura's isotropic calculation gives $A_{\mathrm{anh}} = \SI{7.41e-56}{s^4}$, as tabulated by Leman~\cite[Table~I]{leman2012invited}; a recent recalculation from cryogenic elastic constants gives $\SI{1.15e-55}{s^4}$~\cite[Table~3]{hernandez2025modeling}. No direct measurement exists for silicon, so $A_{\mathrm{anh}}$ is uncertain by at least the factor 1.5 between the two.

\Cref{tab:E-phrates} evaluates \cref{eq:E.12-iso,eq:E.12-anh} for silicon. The ratio of the two rates, $\varGamma_{\mathrm{anh}}/\varGamma_{\mathrm{iso}} = (A_{\mathrm{anh}}/B_{\mathrm{iso}})(E_{\mathrm{ath}}/h)$, equals $(E_{\mathrm{ath}}/h)/\SI{33}{THz}$ and is below one everywhere in the acoustic band. Decay is nevertheless what lowers the frequency, since isotope scattering conserves energy. A \SI{10}{THz} longitudinal phonon decays within \SI{0.13}{ns}. Its daughters near \SI{3}{THz} live some \SI{55}{ns} before decaying again, scattering off isotopes about ten times in between, and below \SI{1}{THz} the decay time exceeds \SI{10}{\micro s}, longer than the time a phonon needs to reach a surface of a centimeter-sized crystal.

\begin{table}[t]
  \caption{Bulk phonon rates in natural silicon from \cref{eq:E.12-iso,eq:E.12-anh} with $B_{\mathrm{iso}} = \SI{2.43e-42}{s^3}$ and $A_{\mathrm{anh}} = \SI{7.41e-56}{s^4}$~\cite{leman2012invited}. The anharmonic rate applies to longitudinal phonons. The isotope mean free path uses the transverse speed \SI{5.4}{km/s}~\cite{leman2012invited}.}
  \label{tab:E-phrates}
  \centering
  \begin{tabular}{ccccc}
    \toprule
    $E_{\mathrm{ath}}/h$ (THz) & $E_{\mathrm{ath}}$ (meV) & $\varGamma_{\mathrm{iso}}$ (s$^{-1}$) & isotope mean free path & $\varGamma_{\mathrm{anh}}$ (s$^{-1}$) \\
    \midrule
    0.5 & 2.07 & \num{1.5e5} & \SI{3.6}{cm} & \num{2.3e3} \\
    1 & 4.14 & \num{2.4e6} & \SI{2.2}{mm} & \num{7.4e4} \\
    3 & 12.4 & \num{2.0e8} & \SI{27}{\micro m} & \num{1.8e7} \\
    10 & 41.4 & \num{2.4e10} & \SI{0.22}{\micro m} & \num{7.4e9} \\
    \bottomrule
  \end{tabular}
\end{table}

\subsection{Free flight, focusing and surfaces}

Between events a phonon moves with the group velocity of its branch, obtained from the Christoffel equation of continuum elasticity, $\rho_T\omega^2\hat e_i = C_{ijkl}Q_jQ_k\hat e_l$, for the unit polarization $\hat{\vect{e}}$ of the elastic wave and the elastic tensor $C_{ijkl}$ (\cref{app:C.18}). In a cubic crystal the group velocity is not parallel to $\vect{Q}$, and an isotropic distribution of wavevectors becomes a strongly peaked distribution of energy flux, the phonon focusing that concentrates ballistic phonons along certain crystal directions~\cite[Sect.~2.4]{hernandez2025modeling}. Using the continuum dispersion is an approximation for the phonons near the top of the acoustic band, whose dispersion bends over; it is good for the few-THz phonons that dominate after the first decays.

A phonon that reaches a surface is reflected specularly, reflected diffusely, down-converted into lower-energy phonons, or transmitted into whatever is attached. The bare crystal surface is described by the probabilities of these outcomes, which are inputs of the simulation, fitted to measured pulse shapes for a given surface preparation. An interface to a sensor film is described by a transmission probability. The acoustic mismatch model treats the interface as a welded plane between two elastic media and computes the transmission of each incident mode from continuity of displacement and stress. The diffuse mismatch model assumes that a phonon forgets its direction at the interface and is transmitted with a probability fixed by the ratio of the phonon densities of states on the two sides~\cite{swartz1989thermal}. Both obey the reciprocity derived in \cref{sec:20.3}. Measured silicon-to-aluminum values, $\bar t_{12}P_{\mathrm{abs}} = 0.3$--0.55, are absorption probabilities and not transmissions~\cite{martinez2019measurements}, as \cref{sec:20.3} shows, and a simulation that uses them must use them in that role.

\subsection{Films}

Inside a superconducting film the phonon breaks Cooper pairs, the resulting quasiparticles relax by emitting phonons, and phonons with energy above $2\Delta_{\mathrm{BCS}}$ break further pairs; this is the Kaplan cascade of \cref{sec:20.4}~\cite{kaplan1976quasiparticle}. A full simulation tracks quasiparticles and phonons in the film. The lumped model used in detector codes treats the film as a single object with a thickness $d_{\mathrm{f}}$, a gap $\Delta_{\mathrm{BCS}}$ and an energy-dependent phonon mean free path $\mathfrak{l}_{\mathrm{f}}(E_{\mathrm{ath}})$. An incident phonon escapes back into the substrate unabsorbed with probability $\exp[-4d_{\mathrm{f}}/\mathfrak{l}_{\mathrm{f}}(E_{\mathrm{ath}})]$, the factor 4 being twice the thickness, for the round trip, times two from the average over incidence angles~\cite[Eq.~(17)]{kelsey2023g4cmp}. Otherwise its energy goes to two quasiparticles, which share it with the BCS density of states weight; quasiparticles below $3\Delta_{\mathrm{BCS}}$ are counted as collected, and those above emit a phonon and continue~\cite[Sect.~2.4]{kelsey2023g4cmp}. Phonons below $2\Delta_{\mathrm{BCS}}$ cannot break pairs and are either absorbed directly with a set probability or returned to the substrate. The output is the quasiparticle yield $Y_{\mathrm{qp}}(E_{\mathrm{ath}})$ of \cref{sec:20.4}, sampled event by event, and the phonons the film sends back into the crystal.

\subsection{Carrier transport}

In a biased detector the electrons and holes of the cascade drift in the field $\vect{\mathcal{E}}$, computed beforehand for the electrode geometry by a finite-element solution of the Laplace equation. A carrier accelerates until it moves faster than the longitudinal sound speed, and then emits acoustic phonons, much as a charge faster than light in a medium emits Cherenkov radiation. These are the Neganov--Trofimov--Luke phonons, and they carry the work $eV_{\mathrm{bias}}$ that the bias does on each pair, the Luke gain of \cref{sec:19.5}~\cite{luke1988voltage}. For a carrier with an isotropic mass $m$ and momentum $k$, and the acoustic deformation potential $D^n$ of its band (\cref{sec:6.2}) taken isotropic, the golden-rule emission rate is (\cref{eq:20.1-luke})
\begin{equation}
  \varGamma_{\mathrm{L}}(k) = \frac{c_{\mathrm{LA}}}{3\,l_0}\left(\frac{k}{k_{\mathrm{s}}}\right)^2\left(1 - \frac{k_{\mathrm{s}}}{k}\right)^3 \quad (k > k_{\mathrm{s}}), \qquad
  k_{\mathrm{s}} = m\,c_{\mathrm{LA}}, \qquad l_0 = \frac{\pi\rho_T}{2m^3(D^n)^2},
  \label{eq:E.12-luke}
\end{equation}
in natural units, with $k_{\mathrm{s}}$ the momentum of a carrier moving at the sound speed and $l_0$ a characteristic scattering length; the polar angle $\theta$ of the emitted phonon, measured from $\vect{k}$, is drawn from $P(\theta) \propto (\cos\theta - k_{\mathrm{s}}/k)^2\sin\theta$ for $0 \le \theta \le \arccos(k_{\mathrm{s}}/k)$~\cite[Eq.~(9)]{kelsey2023g4cmp}, a law that integrates to \cref{eq:E.12-luke}; the total rate printed beside it~\cite[Eq.~(8)]{kelsey2023g4cmp} is lower by the factor $k_{\mathrm{s}}/k$. The rate vanishes below $k_{\mathrm{s}}$, which is the Cherenkov condition, and grows as $k^2$ at large $k$. \Cref{eq:E.12-luke} counts longitudinal phonons only; a carrier faster than the transverse sound speed, which is lower, can also emit transverse ones, and a model that omits them shifts the Luke phonons toward the longitudinal branch.

Holes in silicon and germanium are close to isotropic, and \cref{eq:E.12-luke} applies directly. Electrons sit in anisotropic valleys, six along $\langle100\rangle$ in silicon and four along $\langle111\rangle$ in germanium, each with the mass tensor of components $m^*_\parallel$ along the valley axis and $m^*_\perp$ across it (\cref{sec:3.3}). A Herring--Vogt transformation rescales the momentum along each principal axis by $(m_{\mathrm{c}}/m^*_i)^{1/2}$, with the conductivity mass $3/m_{\mathrm{c}} = 1/m^*_\parallel + 2/m^*_\perp$, and turns the valley into an isotropic one. The field and the emission are handled there as for holes, and the results are transformed back~\cite{herring1956transport}, \cite[Eqs.~(11)--(12)]{kelsey2023g4cmp}. The transformation is not unitary, and the emitted phonon's momentum is transformed back by keeping its magnitude, which conserves energy, and transforming only its direction~\cite[Sect.~V\,C]{leman2012invited}. Electrons also scatter between valleys by emitting or absorbing intervalley phonons, which turns their trajectories toward the field direction. The intervalley rates are parametrized, from measurements in silicon and from a model in germanium, rather than computed from first principles~\cite[Sect.~2.3]{kelsey2023g4cmp}; the drift speed against field is then a check of the whole carrier model, which reproduces measurements in germanium~\cite[Fig.~7]{kelsey2023g4cmp}.

Impurities enter as two processes, each with a mean free path that is fitted to the device. One is trapping, in which a carrier is captured by a neutral impurity and stops, and the other is impact ionization of an impurity, in which a carrier frees a second charge from a charged one~\cite[Eqs.~(13)--(16)]{kelsey2023g4cmp}. Trapping reduces the collected charge and the Luke phonons, impact ionization increases them, and both fill the regions between the single-pair peaks of a charge-resolving silicon detector, from whose shape the four or six mean free paths are fitted~\cite[Fig.~6]{kelsey2023g4cmp}. In a uniform field with trapping alone, the simulation must reproduce the Hecht relation of \cref{sec:20.1}; that is the first check of a trapping model.

The step size for carriers is limited by the requirement that the rate stays nearly constant during a step. Taking a step of half the current mean free time conserves energy, in the sense that the energy of the Luke phonons emitted equals the work done by the field, and longer steps produce too few Luke phonons~\cite[Sect.~V\,D]{leman2012invited}. The number of Luke phonons is large: each pair crossing a bias $V_{\mathrm{bias}}$ releases $eV_{\mathrm{bias}}$, and a code that assumes \SI{2}{meV} per Luke phonon in estimating its event size~\cite[Sect.~3.4]{kelsey2023g4cmp} expects $\SI{100}{eV}/\SI{2}{meV} = 5\times10^4$ phonons per pair at \SI{100}{V}.

\subsection{Collection kernels}

The quantity a detector analysis needs is not the history of each phonon but the probability that energy deposited at $\vect{r}_{\mathrm{ev}}$, in phonons of energy $E_{\mathrm{ath}}$, is absorbed in sensor $s$ within a time window. The collection kernel of \cref{sec:20.5} is this probability, tallied by launching phonons of each energy from a grid of event positions and counting where they end. For charge, the analogous kernel is the induced charge on each electrode, computed from the trajectories with Ramo's theorem. A charge $q$ moving from $\vect{r}_1$ to $\vect{r}_2$ induces $q[\varphi_s(\vect{r}_2) - \varphi_s(\vect{r}_1)]$ on electrode $s$, where $\varphi_s$ is the weighting potential, the potential with electrode $s$ at unit voltage and the others grounded (\cref{sec:20.1}). Kernels separate the transport from the event: once tabulated, they are convolved with the deposit and the phonon spectrum of any event, from the created-phonon spectra of \cref{chap:19} to the down-converted spectra of \cref{sec:20.2}, without rerunning the transport.

The statistical error of a kernel bin with $n$ absorbed out of $N_{\mathrm{l}}$ launched phonons is binomial, $[n(1 - n/N_{\mathrm{l}})]^{1/2}/N_{\mathrm{l}}$. For the collection contrasts of \cref{sec:20.5}, which are differences between kernels of similar size, the useful technique is to run the two cases with common random numbers, so that their statistical fluctuations are correlated and cancel in the difference.

\subsection{Validation and the closure checks}

A transport run is only as good as its checks, and four of them cost nothing: energy, charge, rate and step closure. For energy closure, The deposit, plus the work of the bias on the carriers, must equal the energy absorbed in the sensors, plus the energy lost at bare surfaces, plus the energy in trapped carriers and in particles still in flight at the end, to the round-off of the tallies. A missing term in this balance signals a process that loses energy without booking it, which is the most common silent bug in such codes. For charge closure, the charge induced on all electrodes, summed with the trapped charges weighted by their weighting potentials, must equal the charge created. For rate closure, in a large block with surfaces removed and one process switched on at a time, the measured mean free paths must follow \cref{eq:E.12-iso,eq:E.12-anh} with the input constants, with the slopes $E_{\mathrm{ath}}^{-4}$ and $E_{\mathrm{ath}}^{-5}$; this has been done for the materials of the recent rate tabulation~\cite[Sect.~3]{hernandez2025modeling}. For step closure, halving every time step must leave the kernels unchanged within their statistical errors.

Beyond these internal checks, a run must report its agreement with external benchmarks that test the physics rather than the bookkeeping. They are carrier drift velocities against field, the Hecht curve for a known trap density, the ballistic focusing pattern of a thin crystal, the quasiparticle yield $Y_{\mathrm{qp}}(E_{\mathrm{ath}})$ against the Kaplan cascade of \cref{sec:20.4}, and, where it exists, a measured pulse shape. A run that reports kernels without these numbers cannot be distinguished from one with an energy leak.

\subsection{Cost}

The cost is set by the number of particle steps. A \SI{10}{keV} deposit in silicon that has down-converted to \SI{1}{THz} phonons ($\SI{4.1}{meV}$) holds $2.4\times10^6$ phonons. In a crystal a centimeter across, each phonon scatters off isotopes about every \SI{2}{mm} (\cref{tab:E-phrates}) and bounces off the surfaces many times before a sensor covering a few percent of the surface absorbs it. At a hundred steps per phonon the event needs $\sim10^8$--$10^9$ steps, and a biased detector adds the $5\times10^4$ Luke phonons per pair per \SI{100}{V}. With full tracking, a few-keV event takes a day or more of central processing unit (CPU) time, and with the down-sampling described next about an hour per event~\cite[Sect.~3.4 and Fig.~8]{kelsey2023g4cmp}. Collection kernels are much cheaper than events, since each launched phonon is an independent sample and no event needs to be complete.

Down-sampling generates only a fraction $f$ of the phonons and carriers and gives each the weight $1/f$~\cite[Sect.~3.4]{kelsey2023g4cmp}. It conserves the mean of every tally. It does not conserve the fluctuations: a weighted event has a variance $1/f$ times too large for the collected energy, and a variance computed from weighted samples is not the physical variance. Kernels and means may be computed with down-sampling; the measured-charge Fano factor $F_{\mathrm{meas}}$ of \cref{sec:20.1}, and any statement about event-by-event resolution, may not.

\subsection{Pitfalls}

No closure check sees the errors in three inputs: the sound speed used to convert a decay calculation to the fifth-power law, the starting phonon spectrum, and the surface, intervalley and trapping parameters. The fifth-power law of \cref{eq:E.12-anh} amplifies any inconsistency in that sound speed, as the recalculation found for sapphire~\cite[Sect.~3.1]{hernandez2025modeling}. The transport is usually started from phonons at the Debye frequency, on the grounds that the first decays erase the details~\cite[Sect.~III\,B]{leman2012invited}, which is safe for kernels at long times but not for the prompt signal near the event, the subject of the memory-loss arguments of \cref{sec:20.2}. The surface parameters are fits to the device in which they were measured, so that transferring them to a crystal from a different boule is an assumption to be stated.

\subsection{Silicon}

\Cref{tab:E-transport-si} lists the silicon inputs of a phonon and charge transport run, with the checks that a run must report alongside its kernels.

\begin{table}[t]
  \caption{Phonon and charge transport in silicon: inputs, the quantities they control, and the checks a run must report. Rates are from \cref{eq:E.12-iso,eq:E.12-anh} with the cited constants.}
  \label{tab:E-transport-si}
  \centering
  \begin{tabular}{p{0.36\linewidth}p{0.52\linewidth}}
    \toprule
    Input or check & Silicon \\
    \midrule
    $B_{\mathrm{iso}}$; $A_{\mathrm{anh}}$ & \SI{2.43e-42}{s^3}; \SIrange{7.41e-56}{1.15e-55}{s^4}~\cite{leman2012invited,hernandez2025modeling} \\
    Sound speeds L, T & 9.0, \SI{5.4}{km/s}~\cite{leman2012invited} \\
    Isotope mean free path at \SI{1}{THz} & \SI{2.2}{mm} \\
    Equilibrium branch populations slow transverse (ST) : fast transverse (FT) : longitudinal & 55 : 35 : 10~\cite{leman2012invited} \\
    Si $\to$ Al absorption probability $\bar t_{12}P_{\mathrm{abs}}$ & 0.3--0.55~\cite{martinez2019measurements} \\
    Luke phonons per pair at \SI{100}{V} & $5\times10^4$ (at \SI{2}{meV} each) \\
    Carrier model & six $\langle100\rangle$ electron valleys, Herring--Vogt; isotropic holes; fitted intervalley, trapping and impact-ionization mean free paths \\
    Closure checks & energy, induced charge, rate slopes $E_{\mathrm{ath}}^{-4}$ and $E_{\mathrm{ath}}^{-5}$, time step \\
    Benchmarks & drift velocity, Hecht curve, focusing pattern, $Y_{\mathrm{qp}}(E_{\mathrm{ath}})$, pulse shape \\
    \bottomrule
  \end{tabular}
\end{table}

What the transport adds to the physics of \cref{chap:20} is the geometry: the same rates in a crystal of a different shape, with sensors elsewhere, give different kernels, and every statement about what a readout measures is a statement about a detector, not about silicon.

\section{Hot-carrier cascades from first principles}\label{app:E.13}

An electron given \SI{50}{eV} in silicon ends, a picosecond later, as about 13 electron--hole pairs at the band edges and some \SI{35}{eV} of phonons. Which share goes where, and how much the pair number fluctuates from one event to the next, is decided by a competition between two processes at every energy along the way. One is impact ionization, in which the hot carrier gives part of its energy to a valence electron and lifts it across the gap; the other is phonon emission, in which it gives a few tens of meV to the lattice. \Cref{sec:18.1} writes the cascade as a branching process whose kernels are master-formula rates with the carrier as the projectile, and \cref{sec:18.2} defines the mean energy per pair $\varepsilon_{\mathrm{eh}}$ and the Fano factor $F_{\mathrm{eh}}$ as its first two moments. Here we explain how those kernels are computed from the band structure, the screened interaction and the electron--phonon vertex, how the cascade is then sampled, and what each step costs. Silicon serves as the case where both measurements and a calibrated phenomenological model exist to compare against.

\subsection{Computed quantities}

For a seed class $\Lambda$, the set of initial electron and hole states that the primary excitation produces (\cref{sec:18.3}), the calculation produces the distribution $P_{\mathrm{casc}}(N_{\mathrm{eh}}|\Lambda)$ of the final number of pairs. Its moments are the measured quantities,
\begin{equation}
  \varepsilon_{\mathrm{eh}} = \frac{E_0}{\langle N_{\mathrm{eh}}\rangle}, \qquad
  F_{\mathrm{eh}} = \frac{\langle N_{\mathrm{eh}}^2\rangle - \langle N_{\mathrm{eh}}\rangle^2}{\langle N_{\mathrm{eh}}\rangle},
  \label{eq:E.13-moments}
\end{equation}
for a seed of total energy $E_0$, in the limit of large $E_0$ where both become independent of it. The cascade needs the quasiparticle bands $\varepsilon_{n\vect{k}}$ from the band edges up to the highest energy a carrier reaches, tens of eV for the cascades of \cref{chap:18}; the impact-ionization rate $\varGamma_{\mathrm{ii}}(n\vect{k})$ of a carrier in state $n\vect{k}$, with the distribution of its three products over final states; and the phonon emission and absorption rate $\varGamma_{\mathrm{ph}}(n\vect{k})$, with its final states and phonon branches. The cascade Monte Carlo then follows each carrier from state to state until every carrier is below the impact-ionization threshold $\theta_{\mathrm{ii}}\Eg$ (\cref{sec:18.1}), after which it can only lose energy to phonons and cannot add a pair.

\subsection{Place on the hierarchy}

Each kernel is a golden-rule rate at a definite place on the truncation label $\mathfrak{T}$ of \cref{sec:8.5}. The impact-ionization rate is the decay of a quasiparticle into a lower-energy quasiparticle and an electron--hole pair through the dynamically screened interaction $W^{\mathrm{scr}}$ of \cref{sec:4.8}; with $W^{\mathrm{scr}}$ in the random-phase approximation, its many-body approximation is $\mathcal{X} = \mathrm{RPA}$, applied to quasiparticle rather than Kohn--Sham energies. The phonon rate is one-phonon emission or absorption, $n_g = 2$, with $n_B = 0$. The cascade itself is a sequence of independent golden-rule events: the carrier is on shell between events, and interference between successive events is dropped. This is the same incoherent limit as the transport of \cref{app:E.12}, and it holds when the total scattering rate of a carrier is small compared with its energy above the band edge, the condition that the off-shell parameter of \cref{sec:8.5} be small. A rate of \SI{e14}{s^{-1}} corresponds to a width $\varGamma = \SI{0.066}{eV}$ in natural units, so a carrier \SI{5}{eV} above the conduction-band minimum meets the condition even at rates ten times larger. A carrier a few tens of meV above the impact-ionization threshold, whose energy relative to that threshold is comparable to the width, does not, and the threshold region is where the incoherent picture is least reliable.

\subsection{Quasiparticle bands}

The impact-ionization threshold is set by the gap, and the rate just above it rises steeply with the energy above threshold, so the gap must be right before anything else is. Semilocal density-functional theory places the conduction bands of silicon too low, with an indirect gap of about half the measured \SI{1.17}{eV} (\cref{app:E.1}). Quasiparticle corrections in the $G_0W_0$ approximation (\cref{app:E.3}) repair the gap and the ordering~\cite{hybertsen1986electron}. A self-consistent variant gives \SI{1.23}{eV} for silicon without spin--orbit coupling, and a single empirical scaling of the self-energy correction brings the gap to experiment and the other tabulated conduction-band energies of silicon to within \SI{0.2}{eV} of measurement~\cite[Table~I]{kotani2010impact}. A rigid scissor shift fixes the gap by assuming that the correction is the same for every conduction state, which has to be checked over the tens of eV that the cascade spans.

The band window is large. The number of bands below an energy $E$ above the band bottom is about the number of plane waves in a sphere of radius $k_E = (2m_eE)^{1/2}$ in the unit cell, $\Omega k_E^3/(6\pi^2)$. For silicon, with $\Omega = a^3/4 = \SI{40.05}{\angstrom^3}$, this is about 8 bands up to \SI{20}{eV}, 32 up to \SI{50}{eV} and 90 up to \SI{100}{eV}. All of these bands, on a dense $\vect{k}$ grid, enter both kernels.

\subsection{The impact-ionization kernel}

The rate at which a quasiparticle in an empty state $n\vect{k}$ decays is set by the imaginary part of its self-energy $\Sigma$ (\cref{sec:4.7}). In the $GW$ approximation, $\Sigma = \ii\GF W^{\mathrm{scr}}$, the imaginary part for an unoccupied state is~\cite[Eqs.~(1)--(2)]{kotani2010impact}
\begin{gather}
  \varGamma_{\mathrm{ii}}(n\vect{k}) = 2z_{\mathrm{qp}}\,|\Imag\Sigma_{n\vect{k}}|, \nonumber\\
  \Imag\Sigma_{n\vect{k}} = -\pi\sum_{n'\vect{k}'}\iint\dd^3r\,\dd^3r'\,
  \psi^*_{n\vect{k}}(\vect{r})\psi_{n'\vect{k}'}(\vect{r})\nonumber\\
  \times\Imag W^{\mathrm{scr}}(\vect{r}, \vect{r}', \varepsilon_{n\vect{k}} - \varepsilon_{n'\vect{k}'})\,
  \psi^*_{n'\vect{k}'}(\vect{r}')\psi_{n\vect{k}}(\vect{r}'),
  \label{eq:E.13-imsigma}
\end{gather}
with the sum over empty final states $n'\vect{k}'$ between the Fermi level and $\varepsilon_{n\vect{k}}$, and $z_{\mathrm{qp}}$ the quasiparticle residue of \cref{eq:E.3-qp}, typically about 0.8~\cite[Sect.~II]{kotani2010impact}. The physical content is in $\Imag W^{\mathrm{scr}}$: at energy transfer $\omega = \varepsilon_{n\vect{k}} - \varepsilon_{n'\vect{k}'}$ it is the spectral weight of the excitations the carrier can create, which below the plasmon energy are electron--hole pairs and above it are mostly plasmons. In silicon the plasmon is at \SI{16.6}{eV}~\cite[Table~I]{ramanathan2020ionization}, and for a carrier more than that above the band edge \cref{eq:E.13-imsigma} includes plasmon emission automatically; the plasmon then decays into a pair. A cascade code that takes its rates from \cref{eq:E.13-imsigma} must therefore not add a separate plasmon channel.

The older route writes the rate directly as a golden-rule sum over the two electrons and one hole of the final state. It uses the matrix element of the screened interaction between the initial and final pairs, direct and exchange terms combined as $|M|^2 = \tfrac12(|M_{\mathrm{D}}|^2 + |M_{\mathrm{E}}|^2 + |M_{\mathrm{D}} - M_{\mathrm{E}}|^2)$, and momentum conservation $\vect{k} = \vect{k}' + \vect{k}_1 - \vect{k}_2$~\cite{kane1967electron}, \cite[Eq.~(5)]{kotani2010impact}. For weak $\Imag W^{\mathrm{scr}}$ the two expressions differ by the factor $z_{\mathrm{qp}}$ and by the exchange term, which \cref{eq:E.13-imsigma} omits; together these make the two rates differ by a factor between 0.5 and 1~\cite[Sect.~II]{kotani2010impact}. Omitting exchange alone multiplies the rate by no less than 0.75~\cite[Sect.~II]{kotani2010impact}. On the logarithmic scale on which impact-ionization rates are usually compared this is a small difference. For $F_{\mathrm{eh}}$, which depends on how the rate rises above threshold relative to the phonon rate, it is not negligible, and the choice belongs among the stated approximations of a calculation.

The cascade needs more than the total rate. When an ionization event is selected, the energies and momenta of the three products are drawn from the integrand of the rate, resolved by final state. The kernel is then a table of probabilities. For each initial $(n, \vect{k})$ it holds the probabilities of the final $(n', \vect{k}')$ of the scattered carrier and of the transfer $(\vect{q}, \omega)$ to the pair; the pair's electron and hole are then drawn with the weight of the joint density of states at that $(\vect{q}, \omega)$. The screened interaction enters at all transfers up to the carrier energy, so $W^{\mathrm{scr}}$ is needed with full frequency dependence from 0 to tens of eV and with local fields, which is the dielectric calculation of \cref{app:E.6} carried to high frequency.

\subsection{The electron--phonon kernel}

The phonon rate of a carrier in state $n\vect{k}$ is the golden-rule rate of \cref{sec:6.2},
\begin{equation}
  \begin{split}
  \varGamma_{\mathrm{ph}}(n\vect{k}) = 2\pi\sum_{m\nu}\int\frac{\dd^3Q}{V_{\mathrm{BZ}}}\,|g_{mn\nu}(\vect{k}, \vect{Q})|^2
  \Bigl[&(n_{\vect{Q}\nu} + 1)\,\delta(\varepsilon_{n\vect{k}} - \varepsilon_{m,\vect{k}+\vect{Q}} - \omega_{\vect{Q}\nu})\\ &+ n_{\vect{Q}\nu}\,\delta(\varepsilon_{n\vect{k}} - \varepsilon_{m,\vect{k}+\vect{Q}} + \omega_{\vect{Q}\nu})\Bigr],
  \end{split}
  \label{eq:E.13-gph}
\end{equation}
with $V_{\mathrm{BZ}}$ the zone volume, the emission term carrying $n_{\vect{Q}\nu} + 1$ and the absorption term $n_{\vect{Q}\nu}$. At the millikelvin temperatures of a detector, $n_{\vect{Q}\nu} = 0$ for the optical and zone-boundary phonons that dominate, and only emission survives. The vertex $g_{mn\nu}$ comes from density-functional perturbation theory (\cref{app:E.4}) and is interpolated to fine grids with Wannier functions (\cref{app:E.5})~\cite{ponce2016epw,giustino2017electron}. First-principles rates of this kind, combined with $GW$ electron--electron rates, have been used to follow hot carriers in silicon through their first picosecond~\cite{bernardi2014ab}.

The catch for a cascade is the energy range. The Wannier set that interpolates the band edges, eight $sp^3$-like orbitals per silicon cell spanning the valence bands and the lowest four conduction bands (\cref{app:E.5}), does not represent states tens of eV up, and its interpolated $g$ is meaningless there. Either the Wannier manifold is enlarged, with disentanglement over a window that contains tens of entangled bands, or the vertex is computed directly on a coarser grid at high energy, where the rates are smoother functions of $\vect{k}$ because the density of final states averages over many bands. Above a few eV the phonon rate is dominated by the emission of optical and intervalley phonons of about \SI{60}{meV}~\cite[Sect.~II\,B]{ramanathan2020ionization}, so the energy lost per phonon event is nearly fixed, and what the calculation must get right is the rate, not the detailed spectrum.

\subsection{Full-band cascade Monte Carlo}

With the two kernels tabulated, the cascade is sampled as in full-band device simulation~\cite{fischetti1988monte}, without a field. Each carrier is a state $(n, \vect{k})$; its total rate is $\varGamma_{\mathrm{tot}} = \varGamma_{\mathrm{ii}} + \varGamma_{\mathrm{ph}}$. The time to its next event is drawn by self-scattering: a constant ceiling $\varGamma_{\mathrm{max}} \ge \varGamma_{\mathrm{tot}}$ is chosen, a time is drawn from $\varGamma_{\mathrm{max}}\ee^{-\varGamma_{\mathrm{max}}t}$, and at that time a real event is accepted with probability $\varGamma_{\mathrm{tot}}/\varGamma_{\mathrm{max}}$, otherwise nothing happens~\cite{jacoboni1983monte}. Without a field the times only order the events and do not change the outcome, but they are kept because the same code with a field gives the drift of \cref{app:E.12}. An accepted event is impact ionization with probability $\varGamma_{\mathrm{ii}}/\varGamma_{\mathrm{tot}}$, which replaces the carrier by three, or phonon emission, which lowers its energy by $\omega_{\vect{Q}\nu}$ and moves it to $m, \vect{k}+\vect{Q}$. Holes are followed in the valence bands with their own kernels, and the hole band width, about \SI{12}{eV} in silicon, limits the energy a hole can take~\cite[Sect.~II\,A]{ramanathan2020ionization}. A carrier that falls below $\theta_{\mathrm{ii}}\Eg$ is removed from the list with its pair counted. The cascade ends when the list is empty.

The seed matters at low energy. \Cref{sec:18.1} explains why the cascade must start from the final states of the primary process, with their momentum and band structure, rather than from a carrier of given energy. The energy sharing between the first electron and hole sets how much energy is above threshold, and below about \SI{20}{eV} in silicon different sharing assumptions give different pair yields~\cite[Sect.~IV\,B]{ramanathan2020ionization}. The seed distribution of \cref{sec:18.3} is drawn into the initial carrier list, and $P_{\mathrm{casc}}(N_{\mathrm{eh}}|\Lambda)$ inherits the seed class $\Lambda$.

\subsection{The single-parameter model}

Before the first-principles cascade, silicon already has a cascade model that reproduces the measured pair energy, and it is the natural benchmark. The model of Alig and Bloom, as implemented by Ramanathan and Kurinsky~\cite{ramanathan2020ionization}, replaces both kernels by free-particle forms with a single optical phonon of energy $\omega_{\mathrm{AB}} = \SI{63}{meV}$. For a carrier at energy $E$ above its band edge, the ratio of the phonon and impact-ionization rates is then~\cite[Eq.~(8)]{ramanathan2020ionization}
\begin{equation}
  \frac{\varGamma_{\mathrm{ph}}(E)}{\varGamma_{\mathrm{ii}}(E)} = A_{\mathrm{AB}}\,\frac{105}{2\pi}\,\frac{(E - \omega_{\mathrm{AB}})^{1/2}}{(E - \Eg)^{7/2}} ,
  \label{eq:E.13-ab}
\end{equation}
with energies in eV and one constant $A_{\mathrm{AB}}$, the ratio of the squared phonon and Coulomb matrix elements times phase-space factors, and the probability that the next event is an ionization is $[1 + \varGamma_{\mathrm{ph}}/\varGamma_{\mathrm{ii}}]^{-1}$~\cite[Eq.~(10)]{ramanathan2020ionization}. The value $A_{\mathrm{AB}} = 5.2$, in the units of the source (quoted there as eV$^2$, with which the ratio in \cref{eq:E.13-ab} carries a leftover eV$^{-1}$), was chosen in the original model to reproduce the measured pair energy at high energy, and hot-electron injection data confirm the model up to \SI{5}{eV}~\cite[Sect.~II and Table~1]{ramanathan2020ionization}. It makes the ionization probability per event 0.05 at \SI{3}{eV}, 0.36 at \SI{5}{eV}, 0.77 at \SI{8}{eV} and 0.88 at \SI{10}{eV}, from \cref{eq:E.13-ab} with $\Eg = \SI{1.17}{eV}$. Phonon emission wins below about \SI{6}{eV}, where the two probabilities cross, and the many \SI{63}{meV} losses a carrier suffers on its way down from there are what puts $\varepsilon_{\mathrm{eh}}$ at three times the gap.

Simulated with this model, the high-energy moments follow the linear fits $\varepsilon_{\mathrm{eh}} = 1.7\Eg + 0.084A_{\mathrm{AB}} + 1.3$ and $F_{\mathrm{eh}} = -0.028\Eg + 0.0015A_{\mathrm{AB}} + 0.14$, with energies in eV and $A_{\mathrm{AB}}$ in the same units~\cite[Eqs.~(12)--(13)]{ramanathan2020ionization}. At \SI{123}{K}, where the gap is \SI{1.160}{eV}~\cite[Eq.~(11)]{ramanathan2020ionization}, they give $\varepsilon_{\mathrm{eh}} = \SI{3.71}{eV}$ and $F_{\mathrm{eh}} = 0.115$. The measurement at the same temperature with a charge-resolving silicon sensor is $\varepsilon_{\mathrm{eh}} = \SI{3.749 \pm 0.001}{eV}$ and $F_{\mathrm{eh}} = 0.119 \pm 0.002$~\cite{rodrigues2021absolute}. The pair energy agrees to \SI{1}{\percent} because $A_{\mathrm{AB}}$ was tuned to it; the Fano factor is a prediction and falls two standard deviations low. This is the comparison the first-principles cascade must improve on. Since its kernels contain no fitted constant, matching $\varepsilon_{\mathrm{eh}}$ is a test of the kernels, and $F_{\mathrm{eh}}$ is a second, independent one.

\subsection{Convergence parameters}

The parameters fall into three groups. For the bands: the quasiparticle correction, the number of bands (\cref{eq:E.13-moments} needs every band a carrier can reach) and the $\vect{k}$ grid. For $W^{\mathrm{scr}}$: the plane-wave cutoff of the dielectric matrix, which must resolve momentum transfers up to about twice the largest carrier momentum, the number of bands in the polarizability, the frequency grid from 0 to the largest carrier energy, and the $\vect{q}$ grid. For the kernels: the $\vect{k}$ grid on which $\varGamma_{\mathrm{ii}}$ and $\varGamma_{\mathrm{ph}}$ are tabulated, and the broadening of the energy-conserving delta functions. Near threshold the impact-ionization rate is a sum over few final states that conserve both energy and momentum. On a finite grid it scatters strongly from one $\vect{k}$ to the next; with 1728 $\vect{k}$ points in the zone the numerical error is about a factor two at rates below \SI{e10}{s^{-1}}~\cite[Fig.~2]{kotani2010impact}. The scatter is a physical feature of the band structure smeared by the grid, and it averages out at higher energies, where many final states contribute~\cite[Sect.~III]{kotani2010impact}.

The statistical parameter is the number of simulated cascades $N_{\mathrm{casc}}$. The mean $\langle N_{\mathrm{eh}}\rangle$ converges quickly, since its relative error is $[F_{\mathrm{eh}}/(\langle N_{\mathrm{eh}}\rangle N_{\mathrm{casc}})]^{1/2}$. The variance converges slowly: for a nearly Gaussian pair number, the relative error of a sample variance is $(2/N_{\mathrm{casc}})^{1/2}$. Resolving $F_{\mathrm{eh}}$ to $\pm0.002$ at 0.119, the precision of the silicon measurement, needs a relative error of \SI{1.7}{\percent} and therefore $N_{\mathrm{casc}} \approx 2/0.017^2 \approx 7\times10^3$ cascades at each seed energy, and more if the comparison is to be made at the two-standard-deviation level.

\subsection{Cost}

The dominant cost is the full-frequency screened interaction. For carriers up to \SI{50}{eV} in silicon, the momentum transfers reach $2k_E = \SI{7.2}{\angstrom^{-1}}$, and a dielectric matrix that resolves them holds $N_G \approx \Omega(2k_E)^3/(6\pi^2) \approx 260$ plane waves. Stored at 200 frequencies on the 29 irreducible points of an $8^3$ $\vect{q}$ grid, in double-precision complex numbers, the matrices $W^{\mathrm{scr}}_{\vect{G}\vect{G}'}(\vect{q}, \omega)$ occupy $29\times260^2\times200\times\SI{16}{byte} \approx \SI{6}{GB}$, and computing them requires the polarizability summed over about 30 empty bands on a $\vect{k}$ grid of similar density. For carriers up to \SI{100}{eV} the band count triples, $N_G$ grows to about 730 and $N_G^2$ eightfold. The kernel tables add a sum over final states for every tabulated initial state: with $N_k$ points and $N_b$ bands it scales as $N_k^2N_b^2$ before symmetry, which for the 1728-point grid above and 30 bands is about $3\times10^9$ terms, each a contraction with $W^{\mathrm{scr}}$. The phonon kernel is a standard Wannier-interpolated calculation (\cref{app:E.5}), cheaper than the electronic one once the high-energy vertex problem is solved. The cascade Monte Carlo is cheap: a \SI{5}{keV} seed makes some 1300 pairs and a few times $10^4$ phonon events, so that $10^4$ cascades amount to some $6\times10^8$ events, negligible next to the tabulation of the kernels.

\subsection{Pitfalls}

The gap must match the detector temperature, \SI{1.17}{eV} at millikelvin against \SI{1.12}{eV} at room temperature~\cite[Table~I]{ramanathan2020ionization}, since every rate near threshold depends on the energy above $\theta_{\mathrm{ii}}\Eg$. Two double countings are easy to make: the $GW$ rate already contains plasmon emission, so no plasmon channel is added; and the first-principles phonon rate of \cref{eq:E.13-gph} replaces, and is not added to, a phenomenological one. Two models that reproduce the same $\varepsilon_{\mathrm{eh}}$ do not validate each other's $F_{\mathrm{eh}}$, and the agreement of the cascade of \cref{sec:18.4} with the single-parameter model is agreement of two models, not independent support.

\subsection{Silicon}

\Cref{tab:E-cascade-si} summarizes the silicon inputs, the reference values and the sizes derived in this section.

\begin{table}[t]
  \caption{The hot-carrier cascade in silicon. Band and matrix sizes are free-electron estimates from the text; the model values follow from \cref{eq:E.13-ab} and its fits at \SI{123}{K}~\cite{ramanathan2020ionization}; the measurement is at \SI{123}{K} and \SI{5.9}{keV}~\cite{rodrigues2021absolute}.}
  \label{tab:E-cascade-si}
  \centering
  \begin{tabular}{p{0.43\linewidth}p{0.45\linewidth}}
    \toprule
    Quantity & Silicon \\
    \midrule
    Gap at low temperature; quasiparticle gap without spin--orbit & \SI{1.17}{eV}; \SI{1.23}{eV}~\cite{kotani2010impact} \\
    Plasmon energy; hole band width & \SI{16.6}{eV}; \SI{12}{eV}~\cite{ramanathan2020ionization} \\
    Bands up to 20, 50, \SI{100}{eV} & 8, 32, 90 \\
    $N_G$ and storage of $W^{\mathrm{scr}}$ for \SI{50}{eV} carriers & 260; \SI{6}{GB} (29 $\vect{q}$, 200 frequencies) \\
    $\vect{k}$ grid; threshold-region error & 1728 points; factor 2 below \SI{e10}{s^{-1}}~\cite{kotani2010impact} \\
    Ionization probability per event at 3, 5, \SI{8}{eV} & 0.05, 0.36, 0.77 (\cref{eq:E.13-ab}) \\
    $\varepsilon_{\mathrm{eh}}$, $F_{\mathrm{eh}}$: model & \SI{3.71}{eV}, 0.115 \\
    $\varepsilon_{\mathrm{eh}}$, $F_{\mathrm{eh}}$: measured & \SI{3.749 \pm 0.001}{eV}, $0.119 \pm 0.002$ \\
    Cascades for $\delta F_{\mathrm{eh}} = 0.002$ & $7\times10^3$ per seed energy \\
    \bottomrule
  \end{tabular}
\end{table}

What a first-principles cascade adds is a prediction for the materials of \cref{chap:17}, where no measurement exists, with kernels that were not tuned to any of them; silicon, whose $\varepsilon_{\mathrm{eh}}$ the measurement fixes, is where that prediction is first put to the test.


\section*{Summary}

Every method in this appendix starts from one base calculation, at clamped nuclei with independent particles, and moves one entry of the truncation label at a time. Every numerical parameter is converged on the quantity that is used, not on the total energy: the cutoff that converges the energy of silicon to \SI{0.1}{meV} per atom leaves the density vertex at large momentum transfer unconverged, and the $8^3$ grid that converges the energy to \SI{2}{meV} per atom leaves $\varepsilon_\infty$ nine percent too high. Three rules of bookkeeping recur: a measured gap already contains the zero-point renormalization; every quantity that enters one amplitude is computed in one gauge; and a sum rule tested with a nonlocal Hamiltonian carries its nonlocal term. For silicon the dominant theory uncertainties are few: the gap (about \SI{0.1}{eV}), the screening ($\varepsilon_\infty$ of 13.0 computed against 11.7 measured), the vertex beyond $q \sim 1/r_c$, and, for recoils, the interatomic potential. The rest are controlled numerically.

\printbibliography[heading=subbibliography]
